
% Chapter 9
% Required bibliography keys used in this chapter:
%   BunkeNikolausVoelkl
%   HopkinsSinger2005
%   RobaloNoncommutativeMotivesI
% The chapter also cites the internal results of Chapters 1--8.
%
% This source is written to be included in the existing manuscript.  It does
% not introduce a second document preamble.

\providecommand{\dR}{\mathrm{dR}}
\providecommand{\op}{\mathrm{op}}
\providecommand{\Sp}{\mathrm{Sp}}
\providecommand{\SH}{\mathrm{SH}}
\providecommand{\CAlg}{\operatorname{CAlg}}
\providecommand{\Mod}{\operatorname{Mod}}
\providecommand{\Fun}{\operatorname{Fun}}
\providecommand{\Map}{\operatorname{Map}}
\providecommand{\map}{\operatorname{map}}
\providecommand{\colim}{\operatorname*{colim}}
\providecommand{\lim}{\operatorname*{lim}}
\providecommand{\fib}{\operatorname{fib}}
\providecommand{\cofib}{\operatorname{cofib}}
\providecommand{\Sym}{\operatorname{Sym}}
\providecommand{\SymSeq}{\operatorname{SymSeq}}
\providecommand{\DiffSp}{\operatorname{DiffSp}}
\providecommand{\Disc}{\operatorname{Disc}}
\providecommand{\Nis}{\mathrm{Nis}}
\providecommand{\cn}{\mathrm{cn}}
\providecommand{\fp}{\mathrm{fp}}
\providecommand{\id}{\operatorname{id}}
\providecommand{\RedInf}{\operatorname{Red}_{\mathrm{inf}}}
\providecommand{\CrystInf}{\operatorname{Crys}_{\mathrm{inf}}}
\providecommand{\CrysCoref}{\operatorname{Crys}^{!}_{\mathrm{inf}}}
\providecommand{\DiffPush}{\widehat{(-)}_{!}}

\chapter{Brave-New Differential Cohomology and Differential Umkehr Theory}
\label{chap:brave-new-differential-cohomology}
% Retain the former label so that references from the later draft chapters
% continue to resolve while those chapters are being rewritten.
\label{chap:brave-new-infinitesimal-geometry}

\section{Purpose, scope, and the two levels of the theory}
\label{sec:bn-diff-purpose}

Chapters~1--8 constructed spectral schemes, cotangent and obstruction theory,
smooth and \'etale spectral geometry, the spectral Nisnevich site, unstable
and stable brave-new motivic homotopy theory, Thom twists, compactification,
and the coherent six operations.  The present chapter adds the layer in which
a stable brave-new coefficient can retain differential and secondary
information while its homotopy-invariant shadow remains the already constructed
category \(\SH(S)\).

The construction follows the structural lesson of
Bunke--Nikolaus--V\"olkl: differential cohomology begins with a sheaf valued in
a stable category and separates its homotopy-invariant part from the secondary,
cycle, and deformation parts by adjunctions and exact fibres and cofibres
\cite[Sections~2--3]{BunkeNikolausVoelkl}.  In the present spectral setting raw
Thom stabilization is imposed before affine-line homotopification so that
virtual differential Thom twists and the smooth Umkehr lift are already defined
inside the Thom-stable differential category.  The simultaneous localization
which imposes both Thom stabilization and affine-line invariance is later shown
to be independent of the order of localization; the corrected order is used to
separate Thom periodicity from the BNV homotopification itself.

Two geometric levels are used and are never identified.

\begin{enumerate}
\item The \emph{finite-infinitesimal level} is space-valued.  It records
the equivalence-closed finite-composite class generated by the underlying
arrows of the structured spectral square-zero construction and produces the
intrinsic spectral de Rham localization \(L_{\dR}\).
\item The \emph{stable differential level} is spectrum-valued.  It is obtained
from stable spectral Nisnevich sheaves in two separate stages:
\[
  \mathcal P_S^{\mathrm{pre}}
  \xrightarrow{\ L_{\mathrm{Th},S}\ }
  \mathcal R_S^0
  \xrightarrow{\ \mathbb H_S\ }
  \SH(S).
\]
The first arrow imposes raw Thom stabilization and produces a Thom-stable
coefficient category.  The second arrow imposes only stable
\(\mathbb A^1\)-invariance.  Its inclusion of local objects is bireflective and
therefore gives the exact BNV triple
\[
  \mathbb H_S\dashv\Disc_S\dashv\mathbb S_S.
\]
\end{enumerate}

The first level is infinitesimal geometry; the second is differential stable
cohomology.  Their separation is essential.  The intrinsic de Rham reflector
is generated by spectral square-zero immersions.  The BNV homotopification
\(\mathbb H_S\) is generated only by affine-line projection equivalences in the
already Thom-stable category \(\mathcal R_S^0\).  No equality between these
localizations is asserted.  A canonical stable realization interface is
constructed below from stabilized relative de Rham-local objects to stable
spectral Nisnevich coefficient sheaves and then to \(\mathcal R_S^0\); this
bridge relates the two levels without identifying their localizations.

The final parts of the chapter construct the bridge to the six operations
without introducing a new class of ``admissible'' differential morphisms or
postulating cycle transports.  For a smooth separated morphism of finite
presentation
\[
  p:X\longrightarrow Y
\]
the smooth direct image is first constructed on the Thom-stable differential
categories, differential Thom twists are constructed from vector bundles, and
one obtains a canonical Thom-stable lift
\[
  \widetilde p_!
  :=
  p_\sharp^0\Sigma^{-T_p,0}
  :\mathcal R_X^0\longrightarrow\mathcal R_Y^0
\]
whose motivic homotopification is the exceptional direct image of Chapter~8:
\[
  \mathbb H_Y\widetilde p_!
  \simeq
  p_!\mathbb H_X.
\]
The chapter then constructs, rather than assumes, the canonical BNV correction
\(\widehat T\) of any such lift \(T\), together with its cycle defect and the
associated BNV obstruction object.  The correction is proved to preserve the
motivic-local sector automatically and to be BNV-idempotent:
\[
  \widehat{\widehat T}\simeq\widehat T.
\]
The obstruction therefore measures whether the original lift already agrees
with its canonical BNV completion, rather than whether such a completion
exists.  For the geometric lift \(\widetilde p_!\), the affine-line
Beck--Chevalley equivalence is constructed from smooth base change.  When the
comparison
\[
  \widetilde p_!\longrightarrow\widehat{\widetilde p_!}
\]
is an equivalence, equivalently when its cycle defect is invertible, the
universal BNV interval integral is automatically natural under the original
geometric lift.  This happens for split finite \'etale spectral morphisms and,
by Nisnevich descent, for finite \'etale morphisms which are
Nisnevich-locally split.

\begin{convention}[Universes and geometric scope]
Fix Grothendieck universes
\[
  \mathcal U\subset\mathcal V\subset\mathcal W.
\]
Affine spectral schemes and spectral schemes used as geometric tests are
\(\mathcal U\)-small; small site skeleta are formed in \(\mathcal V\); sheaf
categories, localizations, slices, and functor categories are formed in
\(\mathcal W\).  The finite-infinitesimal constructions impose no
Noetherianity, characteristic, Postnikov-boundedness, or blanket nilpotence
condition.  The stable part uses precisely the qcqs hypotheses already imposed
in Chapters~4--8, and smooth Umkehr statements use the separated
finite-presentation scope of Chapter~8.
\end{convention}

\subsection{Chapter-level target}

The chapter proves the following chain of constructions:
\[
\boxed{
\begin{gathered}
\text{structured square-zero geometry}
\longrightarrow
\mathcal X_{\Sp}^{\mathrm{inf}}
\longrightarrow
L_{\dR},
\\[2mm]
\mathcal D_{\dR}^{\mathrm{st}}(S)
\xrightarrow{\Gamma_{\dR,S}}
\mathcal N_S^{\Sp}
=
\operatorname{Shv}_{\Nis}(\operatorname{Sm}^{\fp}_{/S};\Sp)
\longrightarrow
\mathcal P_S^{\mathrm{pre}}
\xrightarrow{L_{\mathrm{Th},S}}
\mathcal R_S^0,
\\[2mm]
\mathcal N_S^{\Sp}[(\mathbb T_S^{\mathrm{raw}})^{-1}]
\xrightarrow{\ \Psi_S\ }
\mathcal R_S^0
\xrightarrow{\ \Phi_S\ }
\mathcal N_S^{\Sp}[(\mathbb T_S^{\mathrm{raw}})^{-1}],
\qquad
\Phi_S\Psi_S\simeq\id,
\\[2mm]
\mathcal R_S^0
\xrightarrow{\mathbb A^1\text{-homotopification }\mathbb H_S}
\SH(S),
\qquad
\mathbb H_S\dashv\Disc_S\dashv\mathbb S_S,
\\[2mm]
(\mathsf U,\mathsf S,\mathcal A,\mathcal Z,\mathsf Z)
\longrightarrow
\text{BNV reconstruction}
\longrightarrow
\int_{\mathbb A^1}^{\mathrm{BNV}},
\\[2mm]
\text{base functoriality}
+\text{ differential Thom twists}
+\text{ smooth }p_\sharp^0
\longrightarrow
\widetilde p_!=p_\sharp^0\Sigma^{-T_p,0},
\\[2mm]
\mathbb H_Y\widetilde p_!\simeq p_!\mathbb H_X
\longrightarrow
\left(
  \kappa_p:\widetilde p_!\to\widehat{\widetilde p_!},
  \ \omega_p,
  \ \operatorname{Obs}_{\mathrm{BNV}}(p)
\right),
\\[2mm]
\widehat{\widetilde p_!}\text{ preserves motivic-local objects}
\longrightarrow
\widehat{\widehat{\widetilde p_!}}\simeq\widehat{\widetilde p_!},
\\[2mm]
\operatorname{Obs}_{\mathrm{BNV}}(p)\simeq0
\Longleftrightarrow
\kappa_p\text{ is an equivalence},
\\[2mm]
\kappa_p\text{ an equivalence}
+\text{ constructed affine-line Beck--Chevalley}
\longrightarrow
\text{BNV--Umkehr integration naturality for }\widetilde p_!.
\end{gathered}}
\]
The reconstruction and homotopy-formula parts are the brave-new analogues of
\cite[Proposition~3.5 and Theorem~3.2]{BunkeNikolausVoelkl}.  The smooth
Thom-stable lift and its BNV defect are additional spectral-motivic structures
constructed in this chapter.
\section{The native big affine spectral Nisnevich topos}
\label{sec:infinitesimal-native-topos}

Let
\[
    \mathcal C:=\operatorname{AffSp}
\]
be a $\mathcal V$-small skeleton of the $\infty$-category of
$\mathcal U$-small affine spectral schemes of Chapter~1; thus
\[
    \operatorname{AffSp}=\CAlg(\Sp)_{\cn}^{\op}.
\]
For $T\in\mathcal C$, the affine spectral scheme $T$ is quasi-compact and
quasi-separated, so Chapter~\ref{chap:smooth-spectral-geometry} supplies the
spectral Nisnevich site
\[
    \operatorname{Sm}^{\fp}_{/T}.
\]

\begin{definition}[Finite affine Nisnevich basis cover]
\label{def:infinitesimal-affine-nis-basis}
For $T\in\mathcal C$, let
\[
    \operatorname{Cov}^{\acute et}_{\Nis}(T)
\]
denote the $\infty$-category whose objects are finite families of affine
spectral \'etale morphisms
\[
    \{T_\alpha\longrightarrow T\}_{\alpha\in A}
\]
whose generated sieve on the terminal object $T$ is covering for the spectral
Nisnevich topology on $\operatorname{Sm}^{\fp}_{/T}$.  Morphisms are
refinements together with their coherent comparison data.
\end{definition}

\begin{proposition}[Finite affine basis theorem]
\label{prop:infinitesimal-affine-nis-basis}
The families of Definition~\ref{def:infinitesimal-affine-nis-basis} form a
basis for a Grothendieck topology on $\mathcal C$.  They are stable under
pullback, finite composition, and passage to the finite fibre products
occurring in their \v{C}ech nerves.  Every Nisnevich covering sieve of an
affine spectral scheme contains such a finite affine basis cover.
\end{proposition}

\begin{proof}
Let $\mathcal P_T$ be the smallest collection of finite covering families in
$\operatorname{Sm}^{\fp}_{/T}$ which contains all singleton equivalence
families, the empty cover of the initial object, and all two-member families
attached to spectral Nisnevich distinguished squares and which is closed
under pullback and finite composition.  The identity, pullback, and
composition axioms show that $\mathcal P_T$ is a Grothendieck pretopology,
and its associated Grothendieck topology is precisely the spectral Nisnevich
topology of Chapter~\ref{chap:smooth-spectral-geometry}, because both are
generated by the same distinguished families.  Hence every Nisnevich covering
sieve contains the sieve generated by a finite $\mathcal P_T$-cover.

Every object of $\operatorname{Sm}^{\fp}_{/T}$ is quasi-compact and
quasi-separated, because its structural morphism is of finite presentation
and $T$ is affine.  Hence every such object admits a finite affine Zariski
cover.

A finite Zariski cover by quasi-compact opens is a Nisnevich cover.  Indeed,
for a binary such cover $X=U\cup V$, let
\[
    Z:=(X\setminus U)_{\mathrm{red}}.
\]
Then $Z\subseteq V$, the map $V\to X$ is \'etale, and
\[
    V\times_XZ\xrightarrow{\simeq}Z.
\]
Thus $\{U\to X,V\to X\}$ is the covering family of a Nisnevich distinguished
square.  A finite open cover is obtained by iterating the binary case.

Now take a covering sieve of the terminal object $T$.  It contains a finite
$\mathcal P_T$-cover by the preceding pretopology description.  Refine every
member of that finite family by a finite affine Zariski cover.  Transitivity
gives a finite affine spectral \'etale covering family of $T$ contained in
the original sieve.  This proves the basis assertion.

Pullback stability follows from base-change stability of spectral Nisnevich
distinguished squares and of affine \'etale morphisms.  Finite composition is
the transitivity axiom for the topology.  Finally, finite fibre products of
affine spectral schemes are affine, so all finite intersections appearing in
\v{C}ech nerves remain in the basis category.
\end{proof}

\begin{definition}[Big affine spectral Nisnevich topos]
\label{def:infinitesimal-native-topos}
Equip $\mathcal C$ with the Grothendieck topology generated by the finite
affine basis covers of Proposition~\ref{prop:infinitesimal-affine-nis-basis}
and put
\[
    \mathcal X_{\Sp}
    :=
    \operatorname{Shv}_{\Nis}(\mathcal C;\mathcal S).
\]
For $T\in\mathcal C$, write $h_T$ for its representable Nisnevich sheaf.
\end{definition}

\begin{proposition}[Subcanonicity and compatibility]
\label{prop:infinitesimal-big-nis}
The affine spectral Nisnevich topology is subcanonical.  Let $S$ be a spectral
scheme in the standing motivic scope and let $X\to S$ be an affine object of
$\operatorname{Sm}^{\fp}_{/S}$.  Let
$\operatorname{Et}_{\mathrm{aff}}(X)$ denote the $\infty$-category of affine
spectral schemes \'etale over $X$.  The notion of Nisnevich covering family
on $\operatorname{Et}_{\mathrm{aff}}(X)$ obtained from the absolute affine
spectral Nisnevich topology agrees with the notion obtained by viewing the
same diagrams over $S$ in $\operatorname{Sm}^{\fp}_{/S}$.  Consequently, a
finite family of affine spectral \'etale morphisms
$\{X_\alpha\to X\}$ is a covering family for the topology induced from
$\mathcal C$ if and only if it is a covering family when the same family is
regarded in $\operatorname{Sm}^{\fp}_{/S}$.
\end{proposition}

\begin{proof}
Every affine basis cover is a spectral \'etale covering family.  Spectral
schemes and their morphisms satisfy effective \'etale descent, so every
affine representable is a sheaf for the affine Nisnevich topology.  This
proves subcanonicity.

For compatibility, let $Y\to X$ be affine and spectral \'etale.  Since
$X\to S$ is smooth and of finite presentation, the composite $Y\to S$ is
smooth and of finite presentation.  Thus forgetting the map to $X$, or
composing it with $X\to S$, gives the two ambient realizations of
$\operatorname{Et}_{\mathrm{aff}}(X)$ which are being compared.

A Cartesian square of affine \'etale objects over $X$ is spectral Nisnevich
distinguished exactly when its defining diagram has an open immersion, an
\'etale morphism, the reduced closed complement, and the complementary
equivalence.  These conditions depend only on the four spectral schemes and
the maps between them, and are unchanged whether the diagram is viewed
absolutely, over $X$, or over $S$.  Hence the generating distinguished
covers, and therefore the two induced notions of Nisnevich covering family,
agree on $\operatorname{Et}_{\mathrm{aff}}(X)$.  Applying
Proposition~\ref{prop:infinitesimal-affine-nis-basis} gives the asserted
equivalence for finite affine families.
\end{proof}

\begin{proposition}[Native Yoneda realization of spectral schemes]
\label{prop:infinitesimal-native-yoneda}
For every $\mathcal U$-small spectral scheme $X$, the functor of affine points
\[
    h_X:\mathcal C^{\op}\longrightarrow\mathcal S,
    \qquad
    h_X(T):=\Map(T,X),
\]
is a sheaf for the affine spectral Nisnevich topology.  The assignment
$X\mapsto h_X$ defines a fully faithful functor from the $\infty$-category of
$\mathcal U$-small spectral schemes to $\mathcal X_{\Sp}$.  We henceforth regard
such spectral schemes as native objects of $\mathcal X_{\Sp}$ through this embedding.
\end{proposition}

\begin{proof}
Chapter~1 realizes every spectral scheme as a spectral Zariski stack, hence
as a presheaf on affine spectral schemes.  Let
\[
    \{T_\alpha\longrightarrow T\}_{\alpha\in A}
\]
be a finite affine spectral Nisnevich basis cover.  This is, in particular,
an \'etale covering family.  Effective \'etale descent for spectral schemes
and their morphisms therefore gives a canonical equivalence
\[
    \Map(T,X)
    \simeq
    \lim_{[n]\in\Delta}
    \prod_{(\alpha_0,\ldots,\alpha_n)\in A^{n+1}}
    \Map
    \left(
      T_{\alpha_0}\times_T\cdots\times_TT_{\alpha_n},
      X
    \right).
\]
Thus $h_X$ satisfies descent for every finite affine Nisnevich basis cover and
hence is an object of $\mathcal X_{\Sp}$.

The functor of affine points is fully faithful on spectral schemes by the
spectral-stack Yoneda realization of Chapter~1.  Since
$\mathcal X_{\Sp}$ is a full subcategory of the same affine presheaf category,
the restricted functor remains fully faithful.
\end{proof}

\section{The infinitesimal level and the obstruction to cohesion}
\label{sec:infinitesimal-level-terminology}

The terminology used below is categorical rather than cohesive.  An adjoint
triple
\[
    L\dashv M\dashv R
\]
with $L$ and $R$ fully faithful is an adjoint cylinder.  In the topos-theoretic
situation below, $i^*\dashv i_*$ is a geometric embedding and the additional
left adjoint $i_!$ makes that geometric embedding essential.  The essential
image of $i_*$ is reflective through $i^*\dashv i_*$ and coreflective through
$i_*\dashv i^!$.  The functor $i_!$ gives a second fully faithful copy of the
native topos, and the canonical comparison $i_!F\to i_*F$ measures the failure
of these two copies to agree.  We use the geometric names ``infinitesimal
level'', ``infinitesimal cylinder'', ``reduction'', ``crystallization'', and
``crystalline coreflection'' for this particular spectral instance.

\begin{proposition}[Obstruction to absolute cohesion]
\label{prop:infinitesimal-absolute-obstruction}
Let
\[
    \Gamma:\mathcal X_{\Sp}\longrightarrow\mathcal S
\]
denote the direct-image, equivalently global-sections, functor of the terminal
geometric morphism.  Then $\Gamma$ does not admit a right adjoint.
Consequently the native spectral Nisnevich topos is not equipped by the
present construction with the full absolute cohesive adjoint quadruple over
spaces.
\end{proposition}

\begin{proof}
Put
\[
    B:=\mathbb S[1/2]\times\mathbb S[1/3].
\]
The two principal opens
\[
    D(2)=\Spec\mathbb S[1/2],
    \qquad
    D(3)=\Spec\mathbb S[1/3]
\]
cover $\Spec\mathbb S$.  More precisely, let
\[
    Z:=V(2)_{\mathrm{red}}\subseteq\Spec\mathbb Z
\]
be regarded as the corresponding discrete reduced closed spectral subscheme
of $\Spec\mathbb S$.  Since $3$ is invertible modulo $2$,
\[
    D(3)\times_{\Spec\mathbb S}Z\xrightarrow{\simeq}Z.
\]
Thus
\[
    \{D(2)\to\Spec\mathbb S,\ D(3)\to\Spec\mathbb S\}
\]
is a spectral Nisnevich covering family.

Since spectra of finite product algebras form finite coproducts of affine
spectral schemes,
\[
    \Spec B\simeq D(2)\amalg D(3).
\]
Hence the sieve generated by
\[
    \Spec B\longrightarrow\Spec\mathbb S
\]
contains the preceding covering family.  Therefore
\[
    h_{\Spec B}\longrightarrow 1
\]
is an effective epimorphism in $\mathcal X_{\Sp}$.

A global section of $h_{\Spec B}$ would induce a unital map of
$E_\infty$-rings
\[
    B\longrightarrow\mathbb S.
\]
On $\pi_0$ this would give a unital ordinary ring map
\[
    \mathbb Z[1/2]\times\mathbb Z[1/3]\longrightarrow\mathbb Z.
\]
The image of the idempotent $(1,0)$ is an idempotent of $\mathbb Z$, hence is
$0$ or $1$.  Therefore such a map factors through one of the two projections,
but neither $\mathbb Z[1/2]\to\mathbb Z$ nor
$\mathbb Z[1/3]\to\mathbb Z$ exists as a unital ring map.  Hence
\[
    \Gamma(h_{\Spec B})\simeq\varnothing.
\]

If $\Gamma$ admitted a right adjoint, then $\Gamma$ would be both a left and
a right adjoint and would preserve all small colimits as well as all limits.
In particular it would preserve effective epimorphisms: the \v{C}ech nerve
of an effective epimorphism is formed by finite limits and its geometric
realization is a colimit.  The displayed effective epimorphism would then be
sent to
\[
    \varnothing\longrightarrow *,
\]
which is not an effective epimorphism in spaces.  This contradiction proves
the assertion.
\end{proof}

\begin{remark}[No fractured-topos assertion]
\label{rem:infinitesimal-not-fractured}
The infinitesimal level constructed below is not asserted to be a fractured
$\infty$-topos structure.  In particular, the right adjoint $i^*$ of the
future embedding $i_!$ is not conservative in general.  If
\[
    A'=H(k[\varepsilon]/(\varepsilon^2)),
    \qquad
    A=Hk,
\]
and $F=h_{\Spec A'}$, then the canonical comparison
$i_!F\to i_*F$ becomes an equivalence after applying $i^*$, while its value on
the square-zero thickening $\Spec A\to\Spec A'$ is not an equivalence: the
endomorphisms $\varepsilon\mapsto a\varepsilon$ of $A'$ have the same
restriction to $A$.  Thus the infinitesimal adjunction is a different
categorical structure from a fracture in the sense used for petit--gros
topos comparisons.
\end{remark}

\section{Structured spectral square-zero immersions}
\label{sec:infinitesimal-square-zero}

\begin{construction}[Underlying arrows of the structured square-zero construction]
\label{con:infinitesimal-square-zero-underlying-arrow}
Let \(A\) be a connective commutative ring spectrum, let
\(M\in\Mod_A\) be connective, and let
\[
    \eta:L_A\longrightarrow M[1]
\]
be a degree-one derivation.  Chapter~\ref{chap:cotangent-spectral-deformation}
constructs the connective commutative ring spectrum
\[
    A^\eta
    :=
    A\times_{A\oplus M[1]}A,
\]
where the two maps to \(A\oplus M[1]\) are the derivation classified by
\(\eta\) and the zero derivation, together with its canonical projection
\[
    q_\eta:A^\eta\longrightarrow A.
\]
Let
\[
    \Sigma_{\mathrm{sqz}}^{\mathrm{alg}}
    \subset
    \Fun\!\left(
      \Delta^1,
      \CAlg(\Sp)_{\cn}
    \right)
\]
be the equivalence-closed class spanned by the arrows \(q_\eta\) as
\((A,M,\eta)\) varies.  We call
\[
    (A,M,\eta)\longmapsto q_\eta
\]
the \emph{underlying-arrow construction} for structured spectral square-zero
extensions.
\end{construction}

\begin{definition}[Spectral square-zero extension]
\label{def:infinitesimal-square-zero-extension}
A morphism
\[
    q:A'\longrightarrow A
\]
of connective commutative ring spectra is a \emph{spectral square-zero
extension} if, as an object of the arrow category, it belongs to
\(\Sigma_{\mathrm{sqz}}^{\mathrm{alg}}\).  Equivalently, it is equivalent in
the arrow category to an output \(q_\eta:A^\eta\to A\) of
Construction~\ref{con:infinitesimal-square-zero-underlying-arrow}.  No
structured presentation is retained as auxiliary data of the underlying
arrow: whenever structured deformation theory is used, the argument is made
on the constructed arrows \(q_\eta\) and then transported along equivalences
in the arrow category.
\end{definition}

\begin{definition}[Spectral square-zero immersion]
\label{def:infinitesimal-square-zero-immersion}
The affine morphism
\[
    \Spec A\longrightarrow\Spec A'
\]
opposite to a spectral square-zero extension $A'\to A$ is a
\emph{spectral square-zero immersion}.
\end{definition}

\begin{proposition}[Geometric properties]
\label{prop:infinitesimal-square-zero-properties}
Every spectral square-zero immersion is a spectral closed immersion.  It
induces a homeomorphism on underlying topological spaces, isomorphisms on
corresponding residue fields, and an equivalence of small spectral Zariski
sites.  The class of spectral square-zero immersions is stable under arbitrary
affine base change.
\end{proposition}

\begin{proof}
All assertions are invariant under equivalence in the arrow category, so by
Construction~\ref{con:infinitesimal-square-zero-underlying-arrow} it is enough
to prove them for a constructed arrow
\[
    q_\eta:A^\eta\longrightarrow A
\]
with connective coefficient \(M\).  The square-zero extension theorem of
Chapter~\ref{chap:cotangent-spectral-deformation} gives a canonical fibre
sequence
\[
    M\longrightarrow A^\eta\longrightarrow A.
\]
Hence $\pi_0(A^\eta)\to\pi_0(A)$ is surjective and its kernel is square-zero.
The affine spectral closed-immersion criterion applies.  Every prime of
$\pi_0(A^\eta)$ contains the square-zero kernel, so quotient induces a
homeomorphism
\[
    \Spec\pi_0(A)\xrightarrow{\simeq}\Spec\pi_0(A^\eta).
\]
The kernel vanishes in every residue field, giving the residue-field
identification.  Principal opens correspond by lifting their defining
functions across the nilpotent quotient; these principal opens form a basis,
so the small Zariski sites are equivalent.

For affine base change, let $A^\eta\to B'$ be arbitrary and set
\[
    B:=A\otimes_{A^\eta}B',
    \qquad
    J:=M\otimes_AB.
\]
The algebra $B'$ is canonically a deformation of the $A$-algebra $B$ over the
constructed extension $q_\eta$, through the tautological equivalence
\[
    B'\otimes_{A^\eta}A\xrightarrow{\simeq}B.
\]
The affine deformation correspondence of
Chapter~\ref{chap:cotangent-spectral-deformation} canonically constructs the
induced compatible derivation
\[
    \eta_{B'}:L_B\longrightarrow J[1]
\]
and identifies the canonical morphism $B'\to B$, as an arrow, with
\[
    q_{\eta_{B'}}:B^{\eta_{B'}}\longrightarrow B.
\]
The coefficient can also be read directly from the fibre sequence: exact base
change of
\[
    M\longrightarrow A^\eta\longrightarrow A
\]
along $A^\eta\to B'$ gives
\[
    M\otimes_{A^\eta}B'\longrightarrow B'\longrightarrow B,
\]
and, because the $A^\eta$-action on $M$ factors through $A$, there is a
canonical equivalence
\[
    M\otimes_{A^\eta}B'\simeq M\otimes_AB=J.
\]
Thus the base-changed arrow belongs to
\(\Sigma_{\mathrm{sqz}}^{\mathrm{alg}}\).  Passing to affine spectra proves
the assertion.
\end{proof}

\begin{remark}[Square-zero immersions need not be monomorphisms]
\label{rem:infinitesimal-square-zero-not-monomorphism}
The word ``immersion'' records the spectral closed-immersion property; it does
not imply that a spectral square-zero immersion is a monomorphism of spectral
schemes.  For example, let
\[
    A'=H\bigl(k[\varepsilon]/(\varepsilon^2)\bigr),
    \qquad
    A=Hk.
\]
Then $\Spec A\to\Spec A'$ is a spectral square-zero immersion, but the
derived self-intersection has affine coordinate algebra
\[
    A\otimes_{A'}A,
\]
which has nontrivial higher homotopy.  Hence the multiplication
$A\otimes_{A'}A\to A$ is not an equivalence, so the diagonal is not an
equivalence and the morphism is not a monomorphism.
\end{remark}

\begin{lemma}[Relative structured square-zero promotion]
\label{lem:infinitesimal-relative-square-zero}
Let \(C\) and \(R\) be connective commutative ring spectra, let \(J\) be a
connective \(C\)-module, and let
\[
    \delta:L_C\longrightarrow J[1].
\]
Write
\[
    q_\delta:C^\delta\longrightarrow C
\]
for the constructed absolute square-zero arrow of
Construction~\ref{con:infinitesimal-square-zero-underlying-arrow}.  For every
structural morphism
\[
    R\longrightarrow C^\delta,
\]
the resulting diagram
\[
    R\longrightarrow C^\delta\xrightarrow{q_\delta}C
\]
canonically determines a structured square-zero extension in \(\CAlg_R\) by
the same coefficient \(J\).
\end{lemma}

\begin{proof}
The structural map $R\to C^\delta$ is a lift of $R\to C$ through the
constructed absolute square-zero extension \(q_\delta\).  By the pullback
description of \(C^\delta\), equivalently by the fixed square-zero lifting
theorem over the absolute base \(\mathbb S\), this lift canonically determines
a nullhomotopy \(h\) of the composite
\[
    L_R\otimes_RC
    \longrightarrow
    L_C
    \xrightarrow{\delta}
    J[1].
\]
Applying $\map_C(-,J[1])$ to the transitivity cofibre sequence
\[
    L_R\otimes_RC\longrightarrow L_C\longrightarrow L_{C/R}
\]
gives a fibre sequence of mapping spectra
\[
    \map_C(L_{C/R},J[1])
    \longrightarrow
    \map_C(L_C,J[1])
    \longrightarrow
    \map_C(L_R\otimes_RC,J[1]).
\]
Thus the pair consisting of $\delta$ and the canonically determined
nullhomotopy $h$ determines a canonical point of the homotopy fibre,
equivalently a relative classifying derivation
\[
    \overline\delta:L_{C/R}\longrightarrow J[1]
\]
together with the induced homotopy identifying its composite
$L_C\to L_{C/R}\to J[1]$ with $\delta$.

Form the relative structured square-zero extension in $\CAlg_R$ classified by
$\overline\delta$.  The forgetful functor
\[
    \CAlg(\Sp)_{R/}\longrightarrow\CAlg(\Sp)
\]
creates limits, so its underlying absolute structured extension is the
pullback
\[
    C\times_{C\oplus J[1]}C
\]
classified by the composite $L_C\to L_{C/R}\to J[1]$.  By construction this
composite is identified with $\delta$, so the underlying absolute arrow is
canonically the constructed arrow \(q_\delta:C^\delta\to C\).

It remains to compare the $R$-algebra structures.  Under the absolute
square-zero lifting equivalence, lifts $R\to C^\delta$ of $R\to C$ are
identified with nullhomotopies of
\[
    L_R\otimes_RC\longrightarrow L_C\xrightarrow{\delta}J[1].
\]
The given structural morphism represents the nullhomotopy $h$, and the
relative pullback classified by $\overline\delta$ represents the same
nullhomotopy by construction.  The fibre of the lifting equivalence over
$h$ is contractible.  Hence the absolute identification of the two
square-zero pullbacks refines canonically to an equivalence in
$\CAlg(\Sp)_{R/}$ over $C$.  This supplies the required structured extension
in $\CAlg_R$ with all coherence inherited from the lifting equivalence and
transitivity.
\end{proof}

\begin{remark}[No square-zero presentation datum]
\label{rem:infinitesimal-structured-data-retention}
The preceding lemma starts from the constructed arrow \(q_\delta\), not from a
bare underlying arrow supplemented by a chosen presentation.  Its relative
derivation is constructed functorially from the defining derivation
\(\delta\) and the structural morphism \(R\to C^\delta\).  Passing from
constructed arrows to spectral square-zero extensions means passing to the
equivalence-closed essential image of
Construction~\ref{con:infinitesimal-square-zero-underlying-arrow}; no
presentation is retained as additional input.
\end{remark}

\section{Finite spectral infinitesimal thickenings}
\label{sec:infinitesimal-finite-thickenings}

\begin{definition}[Finite spectral infinitesimal thickening]
\label{def:infinitesimal-finite-thickening}
Let
\[
    \operatorname{Inf}_{\mathrm{fin}}
\]
be the smallest class of morphisms in \(\mathcal C\) which contains all
equivalences and all spectral square-zero immersions and is closed under
finite composition.  A morphism
\[
    j:U\longrightarrow T
\]
in $\mathcal C$ is a \emph{finite spectral infinitesimal thickening} if
\(j\in\operatorname{Inf}_{\mathrm{fin}}\).  No finite factorization,
coefficient module, or classifying derivation is retained as additional data
of \(j\).
\end{definition}

\begin{proposition}[Permanence]
\label{prop:infinitesimal-thickening-permanence}
Finite spectral infinitesimal thickenings are stable under equivalence,
composition, and arbitrary affine base change.
\end{proposition}

\begin{proof}
Stability under equivalence and finite composition is built into
Definition~\ref{def:infinitesimal-finite-thickening}.  For base change, let
\(\mathcal B\) be the class of morphisms in \(\mathcal C\) whose arbitrary
affine base changes belong to \(\operatorname{Inf}_{\mathrm{fin}}\).
Equivalences belong to \(\mathcal B\).  Spectral square-zero immersions belong
to \(\mathcal B\) by
Proposition~\ref{prop:infinitesimal-square-zero-properties}.  Since base
change commutes with composition, \(\mathcal B\) is closed under finite
composition.  The defining minimality of
\(\operatorname{Inf}_{\mathrm{fin}}\) therefore gives
\[
    \operatorname{Inf}_{\mathrm{fin}}\subseteq\mathcal B.
\]
Thus every finite spectral infinitesimal thickening is stable under arbitrary
affine base change.
\end{proof}

\begin{lemma}[Reduced support invariance]
\label{lem:infinitesimal-support-invariance}
Every finite spectral infinitesimal thickening $U\to T$ induces a canonical
homeomorphism
\[
    |U|\xrightarrow{\simeq}|T|
\]
and isomorphisms on corresponding residue fields.
\end{lemma}

\begin{proof}
Let \(\mathcal H\) be the class of morphisms in \(\mathcal C\) inducing a
homeomorphism on underlying spaces and isomorphisms on corresponding residue
fields.  The class \(\mathcal H\) contains equivalences and spectral
square-zero immersions by
Proposition~\ref{prop:infinitesimal-square-zero-properties}, and it is closed
under finite composition.  Hence
\(\operatorname{Inf}_{\mathrm{fin}}\subseteq\mathcal H\) by
Definition~\ref{def:infinitesimal-finite-thickening}.
\end{proof}

\section{The finite-infinitesimal arrow category}
\label{sec:infinitesimal-arrow-category}

\begin{definition}[The finite-infinitesimal category]
\label{def:infinitesimal-arrow-category}
Let
\[
    \mathcal C_{\mathrm{inf}}
    \subset
    \Fun(\Delta^1,\mathcal C)
\]
be the full subcategory spanned by finite spectral infinitesimal thickenings.
Define
\[
    p,q:\mathcal C_{\mathrm{inf}}\longrightarrow\mathcal C
\]
by
\[
    p(U\to T)=U,
    \qquad
    q(U\to T)=T,
\]
and let
\[
    i:\mathcal C\longrightarrow\mathcal C_{\mathrm{inf}},
    \qquad
    i(T)=\id_T.
\]
\end{definition}

\begin{proposition}[Endpoint adjunctions]
\label{prop:infinitesimal-endpoint-adjunctions}
There is an adjoint triple
\[
    \boxed{
    q\dashv i\dashv p.
    }
\]
\end{proposition}

\begin{proof}
Let $j:U\to T$ be a finite thickening and $X\in\mathcal C$.  A morphism
$j\to\id_X$ in the arrow category is a commutative square
\[
\begin{tikzcd}
    U \ar[r] \ar[d,"j"'] & X \ar[d,equal]\\
    T \ar[r] & X,
\end{tikzcd}
\]
and is therefore determined by the lower map $T\to X$.  Hence
\[
    \Map_{\mathcal C_{\mathrm{inf}}}(j,iX)
    \simeq
    \Map_{\mathcal C}(qj,X).
\]
Similarly, a morphism $\id_X\to j$ is determined by its upper map $X\to U$,
giving
\[
    \Map_{\mathcal C_{\mathrm{inf}}}(iX,j)
    \simeq
    \Map_{\mathcal C}(X,pj).
\]
These equivalences are natural, so they exhibit $q\dashv i\dashv p$.
\end{proof}

\begin{proposition}[Finite products]
\label{prop:infinitesimal-products}
The category $\mathcal C_{\mathrm{inf}}$ is closed under finite products in
the arrow category.  Moreover,
\[
    i(T\times S)\simeq i(T)\times i(S).
\]
\end{proposition}

\begin{proof}
The terminal arrow is the identity of the terminal affine spectral scheme.
For two thickenings $U\to T$ and $V\to S$, factor their product canonically as
\[
    U\times V
    \longrightarrow
    T\times V
    \longrightarrow
    T\times S.
\]
The first morphism is the affine base change of \(U\to T\), and the second is
the affine base change of \(V\to S\).  Both are finite spectral infinitesimal
thickenings by
Proposition~\ref{prop:infinitesimal-thickening-permanence}, and their
composite is one by the same proposition.  The identity statement is
immediate.
\end{proof}

\section{The target-pullback Nisnevich topology}
\label{sec:infinitesimal-target-pullback-topology}

\begin{definition}[Target-pullback basis cover]
\label{def:infinitesimal-target-pullback-cover}
Let $j:U\to T$ be an object of $\mathcal C_{\mathrm{inf}}$.  A
\emph{target-pullback Nisnevich basis cover} of $j$ is a finite family
\[
    \{j_\alpha:U_\alpha\longrightarrow T_\alpha\}_{\alpha\in A}
    \longrightarrow
    j
\]
for which
\[
    \{T_\alpha\longrightarrow T\}_{\alpha\in A}
\]
belongs to $\operatorname{Cov}^{\acute et}_{\Nis}(T)$ and every square
\[
\begin{tikzcd}[column sep=large,row sep=large]
    U_\alpha \ar[r] \ar[d,"j_\alpha"']
        & U \ar[d,"j"]\\
    T_\alpha \ar[r]
        & T
\end{tikzcd}
\]
is a pullback square.  Equivalently,
\[
    U_\alpha\simeq U\times_TT_\alpha.
\]
\end{definition}

\begin{proposition}[The target-pullback pretopology]
\label{prop:infinitesimal-target-pullback-pretopology}
The target-pullback basis covers contain identities, are stable under
pullback, and are stable under finite composition.  They therefore generate
a Grothendieck topology on $\mathcal C_{\mathrm{inf}}$.
\end{proposition}

\begin{proof}
Identity families are immediate.  Consider a basis cover of $j:U\to T$ and a
morphism $j':U'\to T'\to j$ in the arrow category.  Pulling back the target
cover gives a Nisnevich basis cover
\[
    \{T_\alpha\times_TT'\to T'\}_\alpha
\]
by Proposition~\ref{prop:infinitesimal-affine-nis-basis}.  Its source is
\[
    U'\times_{T'}(T_\alpha\times_TT')
    \simeq U'\times_TT_\alpha,
\]
which is exactly the pullback in the arrow category.  Hence basis covers are
pullback-stable.

For composition, let $\{T_\alpha\to T\}$ be a basis cover and for each
$\alpha$ let $\{T_{\alpha\beta}\to T_\alpha\}_\beta$ be a basis cover.
Their composites form a Nisnevich basis cover of $T$ by transitivity.  Source
terms are obtained by iterated pullback, so the composite family is again
target-pullback.
\end{proof}

\begin{proposition}[Endpoint continuity]
\label{prop:infinitesimal-endpoint-continuity}
The functors
\[
    p,q:\mathcal C_{\mathrm{inf}}\longrightarrow\mathcal C
\]
carry target-pullback basis covers to affine spectral Nisnevich covers, and
$i$ carries affine spectral Nisnevich basis covers to target-pullback basis
covers.
\end{proposition}

\begin{proof}
For $q$ this is the definition.  For $p$, the source family associated with
$\{T_\alpha\to T\}$ is
\[
    \{U\times_TT_\alpha\to U\}_\alpha,
\]
which is a Nisnevich cover by pullback stability.  For $i$, a basis cover
$\{T_\alpha\to T\}$ gives the pullback family
\[
    \{\id_{T_\alpha}\to\id_T\}_\alpha.
\]
\end{proof}

\section{Infinitesimal rigidity of the affine \texorpdfstring{\'etale}{etale} site}
\label{sec:infinitesimal-etale-rigidity}

\begin{lemma}[One-stage \'etale rigidity]
\label{lem:infinitesimal-one-stage-etale-rigidity}
Let $j:U\to T$ be a spectral square-zero immersion of connective affine
spectral schemes.  Pullback induces an equivalence
\[
    j^*:\operatorname{Et}(T)
    \xrightarrow{\simeq}
    \operatorname{Et}(U)
\]
between the $\infty$-categories of spectral \'etale schemes.  It restricts to
an equivalence
\[
    j^*:\operatorname{Et}_{\mathrm{aff}}(T)
    \xrightarrow{\simeq}
    \operatorname{Et}_{\mathrm{aff}}(U)
\]
on affine \'etale objects.  Under the affine equivalence, Nisnevich basis
covers correspond.
\end{lemma}

\begin{proof}
Write $T=\Spec A'$ and $U=\Spec A$.  By
Proposition~\ref{prop:infinitesimal-square-zero-properties},
\[
    \pi_0(A')\longrightarrow\pi_0(A)
\]
is a square-zero extension.  The \'etale lifting argument of
Chapter~\ref{chap:smooth-spectral-geometry}, obtained from the classification
of spectral \'etale algebras over a connective base by ordinary \'etale
algebras over its ring of components and subsequent \'etale gluing, shows
that pullback across such a nilpotent closed immersion induces an equivalence
of the full \'etale $\infty$-categories.  Equivalently, ordinary \'etale
algebras are invariant under nilpotent extensions, and the spectral
classification and gluing carry this invariance to
\[
    j^*:\operatorname{Et}(T)\xrightarrow{\simeq}\operatorname{Et}(U).
\]
The affine part of the same classification shows that affine objects
correspond, giving the asserted restriction to
$\operatorname{Et}_{\mathrm{aff}}$.

It remains to compare the Nisnevich topologies.  By
Proposition~\ref{prop:infinitesimal-square-zero-properties}, $j$ is a
homeomorphism on underlying spaces and induces isomorphisms on corresponding
residue fields.  It also identifies open subsets.  Hence the pullback of a
spectral Nisnevich distinguished square over $T$ is distinguished over $U$,
and conversely every distinguished square over $U$ lifts uniquely, through a
contractible space of \'etale choices, to one over $T$.  Equivalently, the
pointwise criterion for a distinguished square is preserved and reflected.
Thus pullback identifies the generated Nisnevich topologies on
$\operatorname{Et}(T)$ and $\operatorname{Et}(U)$.  Applying the finite
affine basis theorem on both sides shows that finite affine basis covers and
their refinement data correspond.
\end{proof}

\begin{theorem}[Finite infinitesimal \'etale-site rigidity]
\label{thm:infinitesimal-finite-etale-rigidity}
Let
\[
    j:U\longrightarrow T
\]
be a finite spectral infinitesimal thickening.  Pullback induces an
equivalence
\[
    \boxed{
    j^*:\operatorname{Et}(T)
    \xrightarrow{\simeq}
    \operatorname{Et}(U).
    }
\]
It restricts to an equivalence
\[
    j^*:\operatorname{Et}_{\mathrm{aff}}(T)
    \xrightarrow{\simeq}
    \operatorname{Et}_{\mathrm{aff}}(U),
\]
and induces an equivalence of affine Nisnevich basis-cover categories
\[
    \boxed{
    j^*:
    \operatorname{Cov}^{\acute et}_{\Nis}(T)
    \xrightarrow{\simeq}
    \operatorname{Cov}^{\acute et}_{\Nis}(U).
    }
\]
The latter equivalence preserves refinement morphisms, finite fibre products,
and \v{C}ech nerves.
\end{theorem}

\begin{proof}
Let \(\mathcal E\) be the class of morphisms \(j:U\to T\) in \(\mathcal C\)
for which pullback induces the asserted equivalences on full and affine
\'etale categories and on affine Nisnevich basis-cover categories, with the
stated compatibility with refinements, finite fibre products, and
\v{C}ech nerves.  Equivalences belong to \(\mathcal E\).  Every spectral
square-zero immersion belongs to \(\mathcal E\) by
Lemma~\ref{lem:infinitesimal-one-stage-etale-rigidity}.  The class
\(\mathcal E\) is closed under finite composition because pullback functors
compose and equivalences of the indicated categories preserve the stated
finite-limit and covering structures under composition.  Therefore
\[
    \operatorname{Inf}_{\mathrm{fin}}\subseteq\mathcal E
\]
by Definition~\ref{def:infinitesimal-finite-thickening}.  This proves all
assertions for \(j\).
\end{proof}

\begin{corollary}[Source classification of infinitesimal covers]
\label{cor:infinitesimal-source-cover-classification}
For every $j:U\to T$ in $\mathcal C_{\mathrm{inf}}$, source evaluation induces
an equivalence
\[
    \boxed{
    \operatorname{Cov}_{\mathrm{inf}}(j)
    \xrightarrow{\simeq}
    \operatorname{Cov}^{\acute et}_{\Nis}(U),
    }
\]
where the left side is the $\infty$-category of target-pullback basis covers
of $j$ and their refinements.
\end{corollary}

\begin{proof}
A target-pullback basis cover of $j$ is determined by its target cover
$\{T_\alpha\to T\}$, while its source is the pullback family
$\{U\times_TT_\alpha\to U\}$.  Theorem~\ref{thm:infinitesimal-finite-etale-rigidity}
identifies target basis covers with source basis covers, including all
refinement and finite-intersection data.  This is exactly the displayed
equivalence.
\end{proof}

\section{The finite-infinitesimal spectral Nisnevich topos}
\label{sec:infinitesimal-topos}

Define
\[
    \mathcal X_{\Sp}^{\mathrm{inf}}
    :=
    \operatorname{Shv}_{\Nis}
    (\mathcal C_{\mathrm{inf}};\mathcal S)
\]
for the Grothendieck topology generated in
Proposition~\ref{prop:infinitesimal-target-pullback-pretopology}.

\begin{proposition}[Subcanonicity of the infinitesimal topology]
\label{prop:infinitesimal-subcanonical}
The target-pullback Nisnevich topology on $\mathcal C_{\mathrm{inf}}$ is
subcanonical.
\end{proposition}

\begin{proof}
Let $k:V\to S$ be a finite infinitesimal thickening.  For
$j:U\to T$ there is a canonical pullback formula
\[
\Map_{\mathcal C_{\mathrm{inf}}}(j,k)
\simeq
\Map_{\mathcal C}(U,V)
\times_{\Map_{\mathcal C}(U,S)}
\Map_{\mathcal C}(T,S),
\]
where the first map to $\Map(U,S)$ is postcomposition with $V\to S$ and the
second is precomposition with $U\to T$.

A target-pullback basis cover of $j$ is obtained from a Nisnevich basis cover
$\{T_\alpha\to T\}$, and its source family
$\{U_\alpha=U\times_TT_\alpha\to U\}$ is a Nisnevich cover.  In every degree
of the \v{C}ech nerve, target evaluation is the corresponding target
intersection and source evaluation is the corresponding source intersection.
The three representable mapping-space functors appearing in the displayed
pullback therefore satisfy descent by
Proposition~\ref{prop:infinitesimal-big-nis}.  Limits commute with limits in
spaces, so their pullback satisfies the same descent condition.  Hence every
arrow representable is a sheaf.
\end{proof}

Let
\[
    i^*:
    \operatorname{PSh}(\mathcal C_{\mathrm{inf}})
    \longrightarrow
    \operatorname{PSh}(\mathcal C)
\]
be restriction along $i^{\op}$.

\begin{proposition}[Kan-extension formulas]
\label{prop:infinitesimal-kan-formulas}
For every presheaf $F$ on $\mathcal C$, there are canonical equivalences
\[
    \operatorname{Lan}_{i^{\op}}F\simeq q^*F,
    \qquad
    \operatorname{Ran}_{i^{\op}}F\simeq p^*F.
\]
Equivalently,
\[
    (q^*F)(U\to T)=F(T),
    \qquad
    (p^*F)(U\to T)=F(U).
\]
\end{proposition}

\begin{proof}
The adjunctions $q\dashv i\dashv p$ of
Proposition~\ref{prop:infinitesimal-endpoint-adjunctions} become, after
passage to opposite categories, the adjunctions which compute the left and
right Kan extensions along $i^{\op}$.  Since $i^{\op}$ has right adjoint
$q^{\op}$ and left adjoint $p^{\op}$, the pointwise Kan-extension formulas
reduce to precomposition by $q^{\op}$ and $p^{\op}$, respectively.
\end{proof}

\begin{proposition}[The first three sheaf adjoints]
\label{prop:infinitesimal-first-three-adjoints}
Precomposition with $q$ and $p$ preserves Nisnevich sheaves.  Hence the
presheaf adjunctions restrict to sheaf adjunctions
\[
    i_!\dashv i^*\dashv i_*:
    \mathcal X_{\Sp}
    \rightleftarrows
    \mathcal X_{\Sp}^{\mathrm{inf}},
\]
with
\[
    i_!=q^*,
    \qquad
    i_*=p^*.
\]
Both $i_!$ and $i_*$ are fully faithful, and $i_!$ preserves all limits and
in particular finite products.
\end{proposition}

\begin{proof}
Let $F$ be a sheaf on $\mathcal C$.  On a target-pullback basis cover of
$j:U\to T$, the $q^*$-descent diagram is exactly the Nisnevich descent diagram
for the target cover of $T$.  Thus $q^*F$ is a sheaf.  The $p^*$-descent
diagram is the descent diagram for the source cover of $U$, so $p^*F$ is also
a sheaf.  Conversely, if $E$ is a sheaf on $\mathcal C_{\mathrm{inf}}$, then
$i^*E$ is a sheaf because $i$ carries every affine Nisnevich basis cover to a
target-pullback basis cover by
Proposition~\ref{prop:infinitesimal-endpoint-continuity}.

The Kan-extension formulas therefore restrict to sheaves and give the
adjunctions.  Since
\[
    qi=\id_{\mathcal C}=pi,
\]
one has
\[
    i^*i_!\simeq\id,
    \qquad
    i^*i_*\simeq\id,
\]
so both $i_!$ and $i_*$ are fully faithful.  Finally, $i_!=q^*$ is
precomposition, hence preserves all presheaf limits; limits of sheaves are
computed in presheaves.  Thus $i_!$ preserves all sheaf limits and in
particular finite products.
\end{proof}

\section{Source sheafification and the fourth adjoint}
\label{sec:infinitesimal-fourth-adjoint}

Let
\[
    a:\operatorname{PSh}(\mathcal C)\to\mathcal X_{\Sp},
    \qquad
    a_{\mathrm{inf}}:
    \operatorname{PSh}(\mathcal C_{\mathrm{inf}})
    \to\mathcal X_{\Sp}^{\mathrm{inf}}
\]
denote Nisnevich sheafification.

\begin{lemma}[Source pullback preserves local equivalences]
\label{lem:infinitesimal-source-local-equivalences}
Let
\[
    u:P\longrightarrow Q
\]
be a Nisnevich local equivalence in $\operatorname{PSh}(\mathcal C)$.  Then
\[
    p^*u:p^*P\longrightarrow p^*Q
\]
is a Nisnevich local equivalence in
$\operatorname{PSh}(\mathcal C_{\mathrm{inf}})$.
\end{lemma}

\begin{proof}
For a Grothendieck topology on an essentially small $\infty$-category, the
local-equivalence class of sheafification is the strongly saturated class
generated by the covering-sieve monomorphisms
\[
    R\hookrightarrow h_V.
\]
Let $\mathcal K$ be the class of morphisms $v$ in
$\operatorname{PSh}(\mathcal C)$ for which $p^*v$ is an infinitesimal
Nisnevich local equivalence.  Precomposition with $p$ preserves all small
colimits, equivalences, compositions, and retracts, while the infinitesimal
local-equivalence class is strongly saturated.  Hence $\mathcal K$ is
strongly saturated.  It is therefore enough to show that $\mathcal K$
contains every covering-sieve monomorphism.

Fix a Nisnevich covering sieve
\[
    R\hookrightarrow h_V
\]
on an affine spectral scheme $V$.  We prove that
\[
    p^*R\longrightarrow p^*h_V
\]
is an infinitesimal local equivalence.  Let
\[
    j:U\longrightarrow T
\]
be an object of $\mathcal C_{\mathrm{inf}}$ and let
\[
    x\in(p^*h_V)(j)=h_V(U)=\Map_{\mathcal C}(U,V)
\]
be a section.  Pulling back the covering sieve along $x$ gives a covering
sieve
\[
    x^*R\hookrightarrow h_U.
\]
By Proposition~\ref{prop:infinitesimal-affine-nis-basis}, this sieve contains
a finite affine Nisnevich basis cover
\[
    \{U_\alpha\longrightarrow U\}_{\alpha\in A}.
\]
By Corollary~\ref{cor:infinitesimal-source-cover-classification}, through a
contractible space of compatible choices this source cover is the source
family of a target-pullback basis cover
\[
    \{j_\alpha:U_\alpha\longrightarrow T_\alpha\}_{\alpha\in A}
    \longrightarrow
    j.
\]
Since each $U_\alpha\to U$ belongs to the pullback sieve $x^*R$, the
restriction
\[
    x|_{U_\alpha}:U_\alpha\longrightarrow V
\]
defines a section of $R(U_\alpha)$.  Equivalently, the restriction of $x$ to
$j_\alpha$ lifts through
\[
    p^*R\longrightarrow p^*h_V.
\]
Thus every section of $p^*h_V$ is infinitesimal Nisnevich-locally in the
image of $p^*R$.

Apply infinitesimal sheafification.  Since $R\to h_V$ is a monomorphism,
$p^*R\to p^*h_V$ is a monomorphism, and left exactness of
$a_{\mathrm{inf}}$ gives a monomorphism
\[
    a_{\mathrm{inf}}p^*R
    \longrightarrow
    a_{\mathrm{inf}}p^*h_V.
\]
By Proposition~\ref{prop:infinitesimal-first-three-adjoints},
$p^*h_V$ is already an infinitesimal Nisnevich sheaf, so
\[
    a_{\mathrm{inf}}p^*h_V\simeq p^*h_V.
\]
The local lifting property proved above says that the resulting morphism
\[
    a_{\mathrm{inf}}p^*R\longrightarrow p^*h_V
\]
is locally surjective, hence an effective epimorphism.  A morphism in an
$\infty$-topos which is both a monomorphism and an effective epimorphism is
an equivalence.  Therefore
\[
    a_{\mathrm{inf}}p^*R\xrightarrow{\simeq}p^*h_V,
\]
so $p^*(R\hookrightarrow h_V)$ is an infinitesimal Nisnevich local
equivalence.

Hence every covering-sieve generator belongs to $\mathcal K$.  Strong
saturation now shows that $p^*$ carries every Nisnevich local equivalence in
$\operatorname{PSh}(\mathcal C)$ to an infinitesimal Nisnevich local
equivalence.
\end{proof}

\begin{theorem}[Source pullback commutes with sheafification]
\label{thm:infinitesimal-sheafification-commutation}
There is a canonical equivalence of functors
\[
    \boxed{
    a_{\mathrm{inf}}p^*\xrightarrow{\simeq}p^*a.
    }
\]
\end{theorem}

\begin{proof}
Let $P\in\operatorname{PSh}(\mathcal C)$.  The sheafification unit
\[
    P\longrightarrow aP
\]
is a Nisnevich local equivalence.  By
Lemma~\ref{lem:infinitesimal-source-local-equivalences}, the induced map
\[
    p^*P\longrightarrow p^*aP
\]
is an infinitesimal Nisnevich local equivalence.  By
Proposition~\ref{prop:infinitesimal-first-three-adjoints}, the target
$p^*aP$ is already an infinitesimal Nisnevich sheaf.  Hence the displayed
local equivalence exhibits $p^*aP$ as the Nisnevich sheafification of
$p^*P$.  By the universal property of sheafification there is therefore a
canonical equivalence
\[
    a_{\mathrm{inf}}p^*P
    \xrightarrow{\simeq}
    p^*aP,
\]
natural in $P$.
\end{proof}

\begin{corollary}[Colimit preservation of the source embedding]
\label{cor:infinitesimal-i-star-colimits}
The sheaf functor
\[
    i_*=p^*:\mathcal X_{\Sp}\longrightarrow
    \mathcal X_{\Sp}^{\mathrm{inf}}
\]
preserves all small colimits.
\end{corollary}

\begin{proof}
Let $F_\lambda$ be a small diagram of sheaves on $\mathcal C$.  Then
\[
\begin{aligned}
    p^*\left(\colim_{\lambda}^{\mathcal X_{\Sp}}F_\lambda\right)
    &\simeq
    p^*a\left(\colim_{\lambda}^{\operatorname{PSh}(\mathcal C)}F_\lambda\right)\\
    &\simeq
    a_{\mathrm{inf}}p^*
    \left(\colim_{\lambda}^{\operatorname{PSh}(\mathcal C)}F_\lambda\right)\\
    &\simeq
    a_{\mathrm{inf}}
    \left(\colim_{\lambda}^{\operatorname{PSh}(\mathcal C_{\mathrm{inf}})}p^*F_\lambda\right)\\
    &\simeq
    \colim_{\lambda}^{\mathcal X_{\Sp}^{\mathrm{inf}}}p^*F_\lambda.
\end{aligned}
\]
The second equivalence is
Theorem~\ref{thm:infinitesimal-sheafification-commutation}.
\end{proof}

\begin{theorem}[The brave-new infinitesimal adjunction]
\label{thm:brave-new-infinitesimal-adjunction}
The functor $i_*$ admits a right adjoint $i^!$, and there is an adjoint
quadruple
\[
    \boxed{
    i_!\dashv i^*\dashv i_*\dashv i^!:
    \mathcal X_{\Sp}
    \rightleftarrows
    \mathcal X_{\Sp}^{\mathrm{inf}}.
    }
\]
For $E\in\mathcal X_{\Sp}^{\mathrm{inf}}$, the fourth adjoint is characterized
on affine tests by the canonical formula
\[
    \boxed{
    (i^!E)(T)
    \simeq
    \Map_{\mathcal X_{\Sp}^{\mathrm{inf}}}(i_*h_T,E).
    }
\]
The functors $i_!$ and $i_*$ are fully faithful, $i^*$ is left exact, and
$i_!$ preserves all limits and in particular finite products.  Hence the geometric embedding determined by $i^*\dashv i_*$ is essential,
the native topos is the infinitesimal level of
$\mathcal X_{\Sp}^{\mathrm{inf}}$, and $i_!\dashv i^*\dashv i_*$ is its
infinitesimal cylinder.
\end{theorem}

\begin{proof}
By Proposition~\ref{prop:infinitesimal-first-three-adjoints}, one already has
\[
    i_!\dashv i^*\dashv i_*.
\]
The sheaf categories are presentable and $i_*$ preserves all small colimits
by Corollary~\ref{cor:infinitesimal-i-star-colimits}.  The presentable adjoint
functor theorem therefore supplies a right adjoint
\[
    i_*\dashv i^!.
\]
For $T\in\mathcal C$, Yoneda and the adjunction give
\[
\begin{aligned}
    (i^!E)(T)
    &\simeq
    \Map_{\mathcal X_{\Sp}}(h_T,i^!E)\\
    &\simeq
    \Map_{\mathcal X_{\Sp}^{\mathrm{inf}}}(i_*h_T,E),
\end{aligned}
\]
which proves the pointwise formula.

Full faithfulness of $i_!$ and $i_*$ and preservation of all limits by $i_!$
were proved in Proposition~\ref{prop:infinitesimal-first-three-adjoints}.
Restriction along $i$ preserves finite limits, so $i^*$ is left exact.  Thus
$i^*\dashv i_*$ is the inverse/direct-image adjunction of a geometric
embedding, and the additional left adjoint makes it essential.  The remaining
terminology is the convention fixed in
Section~\ref{sec:infinitesimal-level-terminology}.
\end{proof}

\section{Reduction, crystallization, and crystalline coreflection}
\label{sec:infinitesimal-modal-triple}

\begin{definition}[Infinitesimal modalities]
\label{def:infinitesimal-modalities}
Define endofunctors of $\mathcal X_{\Sp}^{\mathrm{inf}}$ by
\[
    \RedInf:=i_!i^*,
    \qquad
    \CrystInf:=i_*i^*,
    \qquad
    \CrysCoref:=i_*i^!.
\]
We call $\RedInf$ the \emph{infinitesimal reduction}, $\CrystInf$ the
\emph{crystallization}, and $\CrysCoref$ the \emph{crystalline
coreflection}.
\end{definition}

\begin{theorem}[The crystalline modal triple]
\label{thm:infinitesimal-modal-triple}
There is an adjoint triple
\[
    \boxed{
    \RedInf\dashv\CrystInf\dashv\CrysCoref.
    }
\]
All three endofunctors are idempotent.  The counit of $i_!\dashv i^*$ and the
unit of $i^*\dashv i_*$ give canonical natural transformations
\[
    \boxed{
    \RedInf E\longrightarrow E\longrightarrow\CrystInf E.
    }
\]
For $E\in\mathcal X_{\Sp}^{\mathrm{inf}}$ and $j:U\to T$ one has canonical
equivalences
\[
    (\RedInf E)(j)\simeq E(\id_T),
    \qquad
    (\CrystInf E)(j)\simeq E(\id_U),
\]
and
\[
    (\CrysCoref E)(j)
    \simeq
    \Map_{\mathcal X_{\Sp}^{\mathrm{inf}}}(i_*h_U,E).
\]
The composite $\RedInf E\to E\to\CrystInf E$ evaluates at $j$ as the
restriction map induced by $U\to T$.
\end{theorem}

\begin{proof}
For $E,F\in\mathcal X_{\Sp}^{\mathrm{inf}}$,
\[
\begin{aligned}
    \Map(\RedInf E,F)
    &=\Map(i_!i^*E,F)\\
    &\simeq\Map(i^*E,i^*F)\\
    &\simeq\Map(E,i_*i^*F)\\
    &=\Map(E,\CrystInf F),
\end{aligned}
\]
so $\RedInf\dashv\CrystInf$.  Similarly,
\[
\begin{aligned}
    \Map(\CrystInf E,F)
    &=\Map(i_*i^*E,F)\\
    &\simeq\Map(i^*E,i^!F)\\
    &\simeq\Map(E,i_*i^!F)\\
    &=\Map(E,\CrysCoref F),
\end{aligned}
\]
so $\CrystInf\dashv\CrysCoref$.

Full faithfulness gives
\[
    i^*i_!\simeq\id,
    \qquad
    i^*i_*\simeq\id,
    \qquad
    i^!i_*\simeq\id.
\]
Substitution proves idempotence.  The first two pointwise formulas follow from
Proposition~\ref{prop:infinitesimal-kan-formulas}; the third follows from the
formula for $i^!$.  The counit and unit of the adjacent adjunctions give the
displayed natural transformations, and their composite evaluates as
$E(\id_T)\to E(\id_U)$.
\end{proof}

\begin{remark}[The word ``crystalline'']
\label{rem:infinitesimal-crystalline-terminology}
The terminology refers only to invariance along the finite spectral
infinitesimal arrows of the present site.  No divided-power structure is part
of $\mathcal C_{\mathrm{inf}}$, so the finite-infinitesimal topos itself is
not called a crystalline topos.  The word ``crystallization'' is used for the
modal reflection $i_*i^*$ because its essential image is the source-constant
copy of the native topos, and the native de Rham-local objects below are
exactly those whose target and source extensions agree along every finite
infinitesimal thickening.
\end{remark}

\section*{Part II: Infinitesimal lifting theory}
\addcontentsline{toc}{section}{Part II: Infinitesimal lifting theory}

\section{Formal smoothness and formal \texorpdfstring{\'etaleness}{etaleness}}
\label{sec:infinitesimal-formal-smooth-etale}

For a morphism $f:X\to S$ in $\mathcal X_{\Sp}$, the canonical natural
transformation $i_!\to i_*$ induces a comparison
\[
    \Theta_f:
    i_!X
    \longrightarrow
    i_!S\times_{i_*S}i_*X.
\]

\begin{proposition}[Pointwise lifting interpretation]
\label{prop:infinitesimal-theta-pointwise}
For every finite infinitesimal thickening $j:U\to T$, the comparison
$\Theta_f$ evaluates as
\[
    \boxed{
    \Theta_f(j):
    X(T)
    \longrightarrow
    S(T)\times_{S(U)}X(U).
    }
\]
Its homotopy fibre over a compatible pair
\[
    T\longrightarrow S,
    \qquad
    U\longrightarrow X
\]
is the space of extensions $T\to X$ together with the specified compatibility
homotopies.
\end{proposition}

\begin{proof}
By the pointwise formulas for the two fully faithful embeddings,
\[
    (i_!X)(j)=X(T),
    \qquad
    (i_*X)(j)=X(U),
\]
and similarly for $S$.  Limits of sheaves are computed objectwise, giving the
displayed formula.  The assertion about the homotopy fibre is the universal
property of the pullback.
\end{proof}

\begin{definition}[Formal \'etaleness and formal smoothness]
\label{def:infinitesimal-formal-smooth-etale}
A native morphism $f:X\to S$ is \emph{formally \'etale with respect to finite
spectral thickenings} if $\Theta_f$ is an equivalence.  It is
\emph{formally smooth with respect to finite spectral thickenings} if
$\Theta_f$ is an effective epimorphism in
$\mathcal X_{\Sp}^{\mathrm{inf}}$.
\end{definition}

\begin{proposition}[Base change of formal lifting properties]
\label{prop:infinitesimal-formal-base-change}
Let
\[
\begin{tikzcd}[column sep=large,row sep=large]
    X' \ar[r] \ar[d,"f'"'] & X \ar[d,"f"]\\
    S' \ar[r] & S
\end{tikzcd}
\]
be Cartesian in $\mathcal X_{\Sp}$.  If $f$ is formally \'etale with respect
to finite spectral thickenings, then so is $f'$.  If $f$ is formally smooth
with respect to finite spectral thickenings, then so is $f'$.
\end{proposition}

\begin{proof}
The comparison itself is stable under base change in the finite-infinitesimal
topos.  Since $i_!=q^*$ preserves all finite limits and $i_*$ is a right
adjoint, hence preserves all limits, the Cartesian identity
$X'\simeq X\times_SS'$ gives canonical equivalences
\[
    i_!X'\simeq i_!X\times_{i_!S}i_!S',
    \qquad
    i_*X'\simeq i_*X\times_{i_*S}i_*S'.
\]
Using naturality of the canonical transformation $i_!\to i_*$, one obtains
\[
\begin{aligned}
    i_!S'\times_{i_*S'}i_*X'
    &\simeq
    i_!S'\times_{i_*S'}
    \left(i_*X\times_{i_*S}i_*S'\right)\\
    &\simeq
    i_!S'\times_{i_*S}i_*X\\
    &\simeq
    \left(i_!S\times_{i_*S}i_*X\right)
    \times_{i_!S}i_!S'.
\end{aligned}
\]
Under these identifications, naturality of $i_!\to i_*$ identifies
$\Theta_{f'}$ with the pullback of $\Theta_f$ along
\[
    i_!S'\longrightarrow i_!S.
\]
Equivalences and effective epimorphisms in an $\infty$-topos are stable under
pullback.  Therefore formal \'etaleness and formal smoothness are preserved by
base change.
\end{proof}

\begin{theorem}[Spectral \'etale implies formally \'etale]
\label{thm:infinitesimal-etale-formally-etale}
Let
\[
    f:X\longrightarrow S
\]
be representable by spectral schemes and spectral \'etale.  Then $f$ is
formally \'etale with respect to finite spectral thickenings.
\end{theorem}

\begin{proof}
Let $j:U\to T$ be a finite infinitesimal thickening and fix a compatible pair
\[
    t:T\longrightarrow S,
    \qquad
    x_U:U\longrightarrow X.
\]
Base change $f$ along $t$ and put
\[
    Y:=X\times_ST.
\]
Then $Y\to T$ is spectral \'etale.  The compatibility of the pair is exactly a
section
\[
    \sigma_U:U\longrightarrow Y_U:=Y\times_TU
\]
of the pulled-back \'etale $U$-scheme.

By Theorem~\ref{thm:infinitesimal-finite-etale-rigidity}, pullback along $j$
is an equivalence
\[
    j^*:\operatorname{Et}(T)
    \xrightarrow{\simeq}
    \operatorname{Et}(U),
\]
carrying the terminal object $T\to T$ to $U\to U$ and $Y\to T$ to
$Y_U\to U$.  Full faithfulness therefore gives a canonical equivalence
\[
    \Map_{\operatorname{Et}(T)}(T,Y)
    \xrightarrow{\simeq}
    \Map_{\operatorname{Et}(U)}(U,Y_U).
\]
Thus $\sigma_U$ has a lift $T\to Y$ through a contractible space of choices.
Equivalently, every homotopy fibre of $\Theta_f(j)$ is contractible.  Hence
$\Theta_f$ is an equivalence.
\end{proof}

\begin{lemma}[Open locality of formal \'etaleness]
\label{lem:infinitesimal-formal-etale-open-local}
Let $f:X\to S$ be formally \'etale with respect to finite spectral
thickenings and let $v:V\hookrightarrow X$ be an open immersion of spectral
schemes.  Then the composite $V\to S$ is formally \'etale with respect to
finite spectral thickenings.
\end{lemma}

\begin{proof}
Let $j:U\to T$ be a finite infinitesimal thickening and fix a compatible pair
\[
    T\longrightarrow S,
    \qquad
    U\longrightarrow V.
\]
After composing with $V\to X$, formal \'etaleness of $X\to S$ gives a
contractibly unique extension
\[
    x:T\longrightarrow X.
\]
Form the pullback open immersion
\[
    T_V:=T\times_XV\longrightarrow T.
\]
Its restriction along $U\to T$ is all of $U$, because the given $U$-point
lands in $V$.  By Lemma~\ref{lem:infinitesimal-support-invariance},
$|U|\to|T|$ is a homeomorphism.  Therefore the open subset $|T_V|\subseteq|T|$
contains all of $|T|$.  Open spectral subschemes are detected on classical
truncations, so $T_V\to T$ is an equivalence.  Hence $x$ factors through
$V$.  The factorization space is contractible because $V\to X$ is a
monomorphism.  Thus the lifting comparison for $V\to S$ is an equivalence.
\end{proof}

\begin{theorem}[Recognition of spectral \'etale morphisms]
\label{thm:infinitesimal-etale-recognition}
Let $f:X\to S$ be representable by spectral schemes and locally of finite
presentation.  Then
\[
    \boxed{
    f\text{ is spectral \'etale}
    \quad\Longleftrightarrow\quad
    f\text{ is formally \'etale with respect to finite spectral thickenings}.
    }
\]
\end{theorem}

\begin{proof}
The forward implication is
Theorem~\ref{thm:infinitesimal-etale-formally-etale}.  For the converse,
base change on the target and work over an affine open
\[
    S=\Spec A.
\]
Let
\[
    V=\Spec B\subseteq X
\]
be an arbitrary affine open.  The morphism $V\to\Spec A$ is affine, hence
quasi-compact and quasi-separated.  Since it is also locally of finite
presentation as a restriction of $f$, it is of finite presentation.
Equivalently, $A\to B$ is of finite presentation.  By
Lemma~\ref{lem:infinitesimal-formal-etale-open-local}, $V\to\Spec A$ is
formally \'etale.

Let $M$ be any connective $B$-module.  The split square-zero extension
\[
    B\oplus M\longrightarrow B
\]
is a structured square-zero extension in $\CAlg_A$.  Formal \'etaleness,
applied to the identity point of $\Spec B$, identifies the space of
$A$-algebra lifts
\[
    B\longrightarrow B\oplus M
\]
restricting to $\id_B$ with a contractible space.  By the derivation universal
property this space is
\[
    \Omega^\infty\map_B(L_{B/A},M).
\]
Take $M=L_{B/A}$.  The relative cotangent complex is connective for a map of
connective commutative ring spectra.  Contractibility says that the identity
of $L_{B/A}$ is homotopic to zero.  Hence
\[
    L_{B/A}\simeq0.
\]
Together with finite presentation this is exactly the affine spectral
\'etale condition.  Since the conclusion holds on every affine open chart of
$X$ over every affine open of $S$, the morphism $f$ is spectral \'etale.
\end{proof}

\begin{lemma}[One-stage local lifting for smooth morphisms]
\label{lem:infinitesimal-smooth-one-stage-lifting}
Let $f:X\to S$ be a smooth morphism of spectral schemes, let
$U\to T$ be an affine spectral square-zero immersion over $S$, and let
$U\to X$ be an $S$-morphism.  Then there exists a finite affine Nisnevich
basis cover
\[
    \{T_\alpha\to T\}_{\alpha\in A}
\]
and $S$-morphisms
\[
    T_\alpha\longrightarrow X
\]
whose restrictions to $U_\alpha:=U\times_TT_\alpha$ agree with the given
map $U\to X$ through specified homotopies.
\end{lemma}

\begin{proof}
The assertion is invariant under equivalence of the square-zero arrow.  By
Definition~\ref{def:infinitesimal-square-zero-extension} it is therefore
enough to treat \(U\to T\) opposite to a constructed arrow
\[
    q_\eta:A^\eta\longrightarrow A
\]
of Construction~\ref{con:infinitesimal-square-zero-underlying-arrow}.
Replace $X$ by $X\times_ST$, so that the given point is a section over $U$.
Choose an affine open cover of $X$.  Its inverse image on the affine spectral
scheme $U$ is an open cover.  Since $U$ is quasi-compact, choose a finite
affine principal refinement
\[
    \{U_\alpha\to U\}_{\alpha\in A}
\]
such that the restricted section factors through an affine open
$V_\alpha\subseteq X$ for every $\alpha$.  The equivalence of small Zariski
sites across a square-zero thickening lifts this finite principal cover to a
finite affine open cover
\[
    \{T_\alpha\to T\}_{\alpha\in A}
\]
with $U_\alpha\simeq U\times_TT_\alpha$.  This is a finite affine Nisnevich
basis cover.

Fix $\alpha$ and write
\[
    T=\Spec R,
    \qquad
    T_\alpha=\Spec C',
    \qquad
    U_\alpha=\Spec C,
    \qquad
    V_\alpha=\Spec B.
\]
The base-change construction in the proof of
Proposition~\ref{prop:infinitesimal-square-zero-properties}, applied to
\(q_\eta\) along \(T_\alpha\to T\), canonically produces a constructed
absolute square-zero arrow
\[
    q_{\delta_\alpha}:C^{\delta_\alpha}\longrightarrow C
\]
and an equivalence of arrows
\[
    (C'\to C)\simeq
    (C^{\delta_\alpha}\to C)
\]
for a connective \(C\)-module \(J\) and derivation
\(\delta_\alpha:L_C\to J[1]\).  Transport the structural morphism
\(R\to C'\) across this canonical arrow equivalence.  Then
Lemma~\ref{lem:infinitesimal-relative-square-zero} constructs the corresponding
structured square-zero extension in \(\CAlg_R\) by \(J\).
The obstruction to lifting the restricted $R$-algebra map
$B\to C$ to $B\to C'$ lies in
\[
    \pi_0\map_C
    \left(
        L_{B/R}\otimes_BC,
        J[1]
    \right).
\]
Smoothness makes
\[
    P:=L_{B/R}\otimes_BC
\]
finite locally free.  Thus $P$ is a retract of a finite free $C$-module, and
the obstruction group is a retract of a finite product of copies of
\[
    \pi_0\map_C(C,J[1])
    \simeq
    \pi_{-1}(J)
    =0.
\]
The obstruction vanishes.  The fixed-extension lifting theorem supplies the
required lift $T_\alpha\to V_\alpha\subseteq X$ together with the specified
homotopy on $U_\alpha$.
\end{proof}

\begin{lemma}[Finite-stage local lifting for smooth morphisms]
\label{lem:infinitesimal-smooth-finite-stage-lifting}
Let $f:X\to S$ be a smooth morphism of spectral schemes and let
\[
    U=U_0\longrightarrow U_1\longrightarrow\cdots\longrightarrow U_n=T
\]
be a finite composable string of affine spectral square-zero immersions over
$S$.  For every $S$-morphism
$U\to X$, there exists a finite affine Nisnevich basis cover
\[
    \{T_\alpha\to T\}_{\alpha\in A}
\]
and $S$-morphisms
\[
    T_\alpha\longrightarrow X
\]
whose restrictions to $U\times_TT_\alpha$ agree with the given map
$U\to X$ through specified homotopies.
\end{lemma}

\begin{proof}
We construct, by induction on $r$, a finite affine Nisnevich basis cover
\[
    \mathcal V_r=\{V_{r,\alpha}\to U_r\}_{\alpha\in A_r}
\]
and maps $V_{r,\alpha}\to X$ whose restrictions to
$U\times_{U_r}V_{r,\alpha}$ agree with the given map $U\to X$.  For $r=0$,
take the singleton identity cover of $U$ and the given map.

Suppose the data have been constructed at stage $r<n$.  Pullback along the
square-zero immersion $U_r\to U_{r+1}$ induces, by
Theorem~\ref{thm:infinitesimal-finite-etale-rigidity}, an equivalence of
affine Nisnevich basis-cover categories.  Hence the cover
$\mathcal V_r$ is the pullback of a finite affine Nisnevich basis cover
\[
    \mathcal W_{r+1}
    =
    \{W_{r+1,\alpha}\to U_{r+1}\}_{\alpha\in A_r},
\]
through a contractible space of compatible choices, with
\[
    W_{r+1,\alpha}\times_{U_{r+1}}U_r
    \simeq
    V_{r,\alpha}.
\]
For every $\alpha$, the morphism
\[
    V_{r,\alpha}\longrightarrow W_{r+1,\alpha}
\]
is the affine base change of $U_r\to U_{r+1}$ and is therefore a spectral
square-zero immersion.  Apply
Lemma~\ref{lem:infinitesimal-smooth-one-stage-lifting} to the map
$V_{r,\alpha}\to X$.  It gives a finite affine Nisnevich basis cover
\[
    \{V_{r+1,\alpha\beta}\to W_{r+1,\alpha}\}_{\beta}
\]
and maps
\[
    V_{r+1,\alpha\beta}\longrightarrow X
\]
extending the given map on
$V_{r,\alpha}\times_{W_{r+1,\alpha}}V_{r+1,\alpha\beta}$ through specified
homotopies.

The composites
\[
    \{V_{r+1,\alpha\beta}\to U_{r+1}\}_{\alpha,\beta}
\]
form a finite affine Nisnevich basis cover by finite composition.  The
specified extension homotopies, together with the induction data, show that
the maps to $X$ restrict to the original map on
$U\times_{U_{r+1}}V_{r+1,\alpha\beta}$.  This completes the induction.

At $r=n$, put $T_\alpha:=V_{n,\alpha}$.  The resulting cover and maps have
the required property.
\end{proof}

\begin{theorem}[Spectral smooth implies formally smooth]
\label{thm:infinitesimal-smooth-formally-smooth}
Every smooth morphism of spectral schemes is formally smooth with respect to
finite spectral thickenings.
\end{theorem}

\begin{proof}
Let $j:U\to T$ be a finite infinitesimal thickening and fix a compatible pair
\[
    T\to S,
    \qquad
    U\to X.
\]
Lemma~\ref{lem:infinitesimal-smooth-finite-stage-lifting} proves the required
local lifting statement for every finite composite of affine spectral
square-zero immersions.  Equivalences satisfy the same lifting statement
trivially, and the statement is invariant under equivalence of the composite
arrow.  By Definition~\ref{def:infinitesimal-finite-thickening}, it therefore
holds for every \(j\in\operatorname{Inf}_{\mathrm{fin}}\).  Thus there is a finite affine
Nisnevich basis cover
\[
    \{T_\alpha\to T\}_{\alpha\in A}
\]
and maps $T_\alpha\to X$ extending the pullbacks of the given map $U\to X$.
The associated family in $\mathcal C_{\mathrm{inf}}$ is the target-pullback
basis cover
\[
    \{U\times_TT_\alpha\to T_\alpha\}_{\alpha\in A}
    \longrightarrow
    (U\to T).
\]
Hence every local section of the target of $\Theta_f$ is Nisnevich-locally in
the image of $\Theta_f$.  Effective epimorphisms of sheaves of spaces are
precisely the locally surjective morphisms.  Therefore $\Theta_f$ is an
effective epimorphism.
\end{proof}

\begin{remark}[No additional global finite-presentation condition]
\label{rem:infinitesimal-smooth-no-global-fp}
The proof uses the local cotangent-complex notion of smoothness from
Chapter~\ref{chap:smooth-spectral-geometry}: local finite presentation and a
finite locally free relative cotangent complex.  No additional global
quasi-compactness or quasi-separatedness condition is needed for the formal
smoothness statement itself.
\end{remark}

\section*{Part III: Native de Rham localization and finite spectral crystals}
\addcontentsline{toc}{section}{Part III: Native de Rham localization and finite spectral crystals}

\section{Square-zero localization and the de Rham reflector}
\label{sec:infinitesimal-dr-localization}

\begin{definition}[Square-zero generators]
\label{def:infinitesimal-dr-generators}
Let $\Sigma_{\mathrm{sqz}}$ be a set of representatives, in the fixed
universe, of all morphisms of representable Nisnevich sheaves
\[
    h_{\Spec A}\longrightarrow h_{\Spec A'}
\]
induced by spectral square-zero extensions $A'\to A$.
\end{definition}

\begin{lemma}[Affine density in slices]
\label{lem:infinitesimal-affine-density}
Let $S\in\mathcal X_{\Sp}$ be arbitrary and let $Y\to S$ be an object of
$\mathcal X_{\Sp,/S}$.  Write $\operatorname{AffPt}_S(Y)$ for the
$\infty$-category whose objects are pairs $(T,u)$ with $T\in\mathcal C$ and
a morphism
\[
    u:h_T\longrightarrow Y
\]
in the slice over $S$, and whose morphisms are the evident commuting triangles.
Then the canonical morphism
\[
    \colim_{(T,u)\in\operatorname{AffPt}_S(Y)}h_T
    \xrightarrow{\simeq}
    Y
\]
is an equivalence in $\mathcal X_{\Sp,/S}$.
\end{lemma}

\begin{proof}
In the presheaf category, an object over $S$ is the canonical colimit of the
representables mapping to it in the corresponding slice; its indexing category
is precisely $\operatorname{AffPt}_S(Y)$ after restricting to the chosen affine
skeleton.  Sheafification is a left adjoint and therefore preserves this
colimit.  Since $Y$ is already a sheaf, the sheafified colimit is $Y$.  The
argument does not require $S$ to be representable.
\end{proof}

\begin{definition}[Finite-infinitesimal congruence]
\label{def:infinitesimal-dr-congruence}
Let
\[
    \mathcal W_{\dR}
\]
be the smallest strongly saturated class of morphisms in
$\mathcal X_{\Sp}$ which contains $\Sigma_{\mathrm{sqz}}$ and is stable under
pullback.
\end{definition}

\begin{lemma}[The generated square-zero congruence]
\label{lem:infinitesimal-dr-generated-congruence}
The class $\mathcal W_{\dR}$ is the congruence generated by
$\Sigma_{\mathrm{sqz}}$.  In particular, if
$f:A\to B$ belongs to $\mathcal W_{\dR}$, then every iterated diagonal of
$f$ also belongs to $\mathcal W_{\dR}$.
\end{lemma}

\begin{proof}
Let
\[
    \mathcal W_0:=\operatorname{Sat}(\Sigma_{\mathrm{sqz}})
\]
be the ordinary strongly saturated class generated by
$\Sigma_{\mathrm{sqz}}$.  Since the generating family is a set in the fixed
universe, the accessible localization of $\mathcal X_{\Sp}$ at
$\Sigma_{\mathrm{sqz}}$ exists; denote it temporarily by $L_\Sigma$.  Its
local-equivalence class is $\mathcal W_0$.

We first show that every pullback of every square-zero generator belongs to
$\mathcal W_0$.  Let
\[
    w:h_U\longrightarrow h_T
\]
be a generator and let $Y\to h_T$ be arbitrary.  In the slice over $h_T$,
the canonical density presentation writes $Y$ as a colimit of affine
representables
\[
    Y\simeq\colim_{(V,u)\in\operatorname{AffPt}_{h_T}(Y)}h_V.
\]
Colimits in an $\infty$-topos are universal, so
\[
    Y\times_{h_T}h_U
    \simeq
    \colim_{(V,u)\in\operatorname{AffPt}_{h_T}(Y)}
    \left(h_V\times_{h_T}h_U\right).
\]
For every affine $V\to T$, the pullback
\[
    V\times_TU\longrightarrow V
\]
is again a spectral square-zero immersion by
Proposition~\ref{prop:infinitesimal-square-zero-properties}.  Hence, up to
the chosen representatives in $\Sigma_{\mathrm{sqz}}$, every morphism
\[
    h_{V\times_TU}\longrightarrow h_V
\]
is a square-zero generator.  The functor $L_\Sigma$ preserves colimits and
inverts each of these stage maps.  Therefore it inverts their colimit
\[
    Y\times_{h_T}h_U\longrightarrow Y,
\]
which proves that every pullback of a generator lies in $\mathcal W_0$.

Let $\mathcal W_0^{\mathrm{univ}}$ be the class of morphisms whose every
pullback belongs to $\mathcal W_0$.  Base-change functors in an
$\infty$-topos preserve all colimits, hence in particular pushouts and
transfinite compositions, and they preserve equivalences, compositions, and
retracts.  Therefore $\mathcal W_0^{\mathrm{univ}}$ is strongly
saturated.  The preceding
paragraph shows that
\[
    \Sigma_{\mathrm{sqz}}\subseteq\mathcal W_0^{\mathrm{univ}}.
\]
By strong saturation,
\[
    \mathcal W_0
    \subseteq
    \mathcal W_0^{\mathrm{univ}}.
\]
Thus $\mathcal W_0$ is stable under pullback.

It follows that $\mathcal W_0$ is a strongly saturated pullback-stable class
containing $\Sigma_{\mathrm{sqz}}$, so by the defining minimality of
$\mathcal W_{\dR}$,
\[
    \mathcal W_{\dR}\subseteq\mathcal W_0.
\]
Conversely, $\mathcal W_{\dR}$ is strongly saturated and contains
$\Sigma_{\mathrm{sqz}}$, hence
\[
    \mathcal W_0\subseteq\mathcal W_{\dR}.
\]
Therefore
\[
    \mathcal W_{\dR}
    =
    \operatorname{Sat}(\Sigma_{\mathrm{sqz}}),
\]
and this class is stable under pullback.

Finally, let $f:A\to B$ belong to $\mathcal W_{\dR}$.  The first projection
\[
    \operatorname{pr}_1:A\times_BA\longrightarrow A
\]
is a pullback of $f$, hence belongs to $\mathcal W_{\dR}$.  Its diagonal
section satisfies
\[
    \operatorname{pr}_1\circ\Delta_f\simeq\id_A.
\]
Two-out-of-three gives
\[
    \Delta_f\in\mathcal W_{\dR}.
\]
Iteration proves the same assertion for every iterated diagonal.  Hence
$\mathcal W_{\dR}$ is the congruence generated by
$\Sigma_{\mathrm{sqz}}$.
\end{proof}

\begin{theorem}[The intrinsic spectral de Rham localization]
\label{thm:infinitesimal-dr-localization}
There is an accessible left-exact localization
\[
\begin{tikzcd}[column sep=large]
    \mathcal X_{\Sp}
        \ar[r,"L_{\dR}",bend left=12]
    &
    \mathcal X_{\Sp}^{\dR}
        \ar[l,hook',"\iota_{\dR}"',bend left=12]
\end{tikzcd}
\]
whose local-equivalence class is $\mathcal W_{\dR}$.  It is initial among
accessible left-exact localizations of $\mathcal X_{\Sp}$ which invert every
spectral square-zero immersion.
\end{theorem}

\begin{proof}
Let
\[
    L_\Sigma:
    \mathcal X_{\Sp}\longrightarrow\mathcal X_\Sigma
\]
be the accessible localization generated by the set
$\Sigma_{\mathrm{sqz}}$.  Its local-equivalence class is the ordinary
strongly saturated class
\[
    \operatorname{Sat}(\Sigma_{\mathrm{sqz}})
    =
    \mathcal W_{\dR}
\]
by Lemma~\ref{lem:infinitesimal-dr-generated-congruence}.  The same lemma
shows that this class is stable under pullback.  The left-exactness criterion
for accessible localizations of an $\infty$-topos therefore implies that
$L_\Sigma$ is left exact.  Put
\[
    L_{\dR}:=L_\Sigma,
    \qquad
    \mathcal X_{\Sp}^{\dR}:=\mathcal X_\Sigma.
\]

If another accessible left-exact localization inverts every spectral
square-zero immersion, then its local-equivalence class is strongly saturated,
stable under pullback, and contains $\Sigma_{\mathrm{sqz}}$.  It therefore
contains $\mathcal W_{\dR}$ by minimality, so the localization factors
uniquely through $L_{\dR}$.  This is the asserted initiality.
\end{proof}

\begin{definition}[Spectral de Rham object]
\label{def:infinitesimal-absolute-dr-object}
For $X\in\mathcal X_{\Sp}$ define
\[
    X_{\dR}:=L_{\dR}X,
\]
with unit
\[
    \eta_X:X\longrightarrow X_{\dR}.
\]
\end{definition}

\begin{theorem}[De Rham locality criterion]
\label{thm:infinitesimal-dr-locality}
For $F\in\mathcal X_{\Sp}$ the following conditions are equivalent:
\begin{enumerate}
\item $F$ is $L_{\dR}$-local;
\item for every spectral square-zero extension $A'\to A$, restriction is an
equivalence
\[
    F(\Spec A')\xrightarrow{\simeq}F(\Spec A);
\]
\item the same restriction map is an equivalence for every finite spectral
infinitesimal thickening.
\end{enumerate}
\end{theorem}

\begin{proof}
By Theorem~\ref{thm:infinitesimal-dr-localization}, $L_{\dR}$ is the
accessible localization generated by the set
$\Sigma_{\mathrm{sqz}}$.  Thus $F$ is $L_{\dR}$-local if and only if it is
right orthogonal to every generator
\[
    h_{\Spec A}\longrightarrow h_{\Spec A'}
\]
coming from a spectral square-zero extension $A'\to A$.  By Yoneda, this
orthogonality condition is exactly the requirement that
\[
    F(\Spec A')\longrightarrow F(\Spec A)
\]
be an equivalence.  Hence (1) and (2) are equivalent.  For a fixed \(F\), the class of affine
arrows along which restriction is an equivalence contains equivalences and
the square-zero immersions of (2) and is closed under finite composition.
Definition~\ref{def:infinitesimal-finite-thickening} therefore gives (3).
Conversely, (3) implies (2) because every square-zero immersion is a finite
spectral infinitesimal thickening.
\end{proof}

\section{Crystals and recognition by the infinitesimal cylinder}
\label{sec:infinitesimal-dr-recognition}

For every $F\in\mathcal X_{\Sp}$, the canonical natural transformation
$i_!\to i_*$ gives
\[
    \delta_F:i_!F\longrightarrow i_*F.
\]
At a thickening $j:U\to T$ it is the restriction map
\[
    \delta_F(j):F(T)\longrightarrow F(U).
\]

\begin{definition}[Finite spectral crystal]
\label{def:infinitesimal-finite-crystal}
A native spectral Nisnevich sheaf $F\in\mathcal X_{\Sp}$ is a
\emph{finite spectral crystal} if the canonical comparison
\[
    \delta_F:i_!F\longrightarrow i_*F
\]
is an equivalence.  Equivalently, $F$ is invariant under every finite spectral
infinitesimal thickening.
\end{definition}

\begin{theorem}[Crystalline recognition of de Rham locality]
\label{thm:infinitesimal-dr-recognition}
For $F\in\mathcal X_{\Sp}$, the following are equivalent:
\begin{enumerate}
\item $F$ is $L_{\dR}$-local;
\item $F$ is a finite spectral crystal;
\item the canonical morphism
\[
    \boxed{
    \delta_F:i_!F\xrightarrow{\simeq}i_*F
    }
\]
is an equivalence in $\mathcal X_{\Sp}^{\mathrm{inf}}$.
\end{enumerate}
Thus $\mathcal X_{\Sp}^{\dR}$ is canonically the reflective subcategory of
finite spectral crystals in the native spectral Nisnevich topos.
\end{theorem}

\begin{proof}
By the pointwise formulas for $i_!$ and $i_*$, the third condition says that
\[
    F(T)\longrightarrow F(U)
\]
is an equivalence for every finite spectral infinitesimal thickening
$U\to T$.  This is condition (3) of
Theorem~\ref{thm:infinitesimal-dr-locality}, and it is exactly the crystal
condition of Definition~\ref{def:infinitesimal-finite-crystal}.
\end{proof}

\begin{corollary}[Crystals and formal \'etaleness]
\label{cor:infinitesimal-native-formally-etale}
For $F\in\mathcal X_{\Sp}$, the following conditions are equivalent:
\begin{enumerate}
\item $F$ is $L_{\dR}$-local;
\item the terminal morphism $F\to1$ is formally \'etale with respect to finite
spectral thickenings;
\item the reduction and crystallization extensions of $i_*F$ agree:
\[
    \RedInf(i_*F)
    \simeq
    i_!F
    \xrightarrow{\simeq}
    i_*F
    \simeq
    \CrystInf(i_*F).
\]
\end{enumerate}
\end{corollary}

\begin{proof}
For the terminal morphism $F\to1$, one has $i_!1\simeq1\simeq i_*1$, and the
comparison $\Theta_{F\to1}$ is exactly
$\delta_F:i_!F\to i_*F$.  Hence (1) and (2) are equivalent by
Theorem~\ref{thm:infinitesimal-dr-recognition}.  The identifications in (3)
follow from $i^*i_*\simeq\id$.
\end{proof}

\begin{remark}[De Rham localization is not crystallization]
\label{rem:infinitesimal-dr-not-crystallization}
The native reflector $L_{\dR}$ must not be identified with a composite of the
ambient modal endofunctors.  For example,
\[
    i^*\CrystInf i_!
    =i^*i_*i^*i_!
    \simeq\id_{\mathcal X_{\Sp}}.
\]
The meaningful relation is instead the crystal-recognition theorem
\[
    F\text{ de Rham-local}
    \quad\Longleftrightarrow\quad
    i_!F\simeq i_*F.
\]
\end{remark}

\section{Relative spectral de Rham objects}
\label{sec:infinitesimal-relative-dr}

\begin{definition}[Relative de Rham object]
\label{def:infinitesimal-relative-dr-object}
For a morphism $X\to S$ in $\mathcal X_{\Sp}$ define
\[
    \boxed{
    X_{\dR/S}
    :=
    S\times_{S_{\dR}}X_{\dR}.
    }
\]
The induced morphism
\[
    \eta_{X/S}:X\longrightarrow X_{\dR/S}
\]
is the relative de Rham unit.
\end{definition}

\begin{theorem}[Slice localization]
\label{thm:infinitesimal-slice-localization}
For every $S\in\mathcal X_{\Sp}$, the assignment
\[
    (X\to S)
    \longmapsto
    (X_{\dR/S}\to S)
\]
defines an accessible left-exact localization
\[
    L_{\dR,/S}:
    \mathcal X_{\Sp,/S}
    \longrightarrow
    \mathcal X_{\Sp,/S}^{\dR}.
\]
Its local objects are precisely those $F\to S$ for which, for every affine
spectral square-zero immersion $U\to T$ over $S$, restriction induces an
equivalence
\[
    \Map_S(T,F)\xrightarrow{\simeq}\Map_S(U,F).
\]
\end{theorem}

\begin{proof}
Applying $L_{\dR}$ to structure morphisms gives
\[
    \Lambda_S:
    \mathcal X_{\Sp,/S}
    \longrightarrow
    \mathcal X_{\Sp}^{\dR}{}_{/L_{\dR}S}.
\]
Its right adjoint sends $E\to L_{\dR}S$ to
\[
    S\times_{L_{\dR}S}\iota_{\dR}E\longrightarrow S.
\]
The right adjoint is fully faithful because $L_{\dR}$ is left exact and
idempotent.  The associated reflector is therefore
\[
    X\longmapsto S\times_{L_{\dR}S}L_{\dR}X=X_{\dR/S}.
\]
Accessibility and left exactness follow from those of $L_{\dR}$ and from
finite limits.  The locality criterion is proved in the slice exactly as in
Theorem~\ref{thm:infinitesimal-dr-locality}, using
Lemma~\ref{lem:infinitesimal-affine-density} and affine base-change stability
of square-zero immersions.
\end{proof}

\begin{theorem}[Base change for relative de Rham objects]
\label{thm:infinitesimal-relative-dr-base-change}
For every pullback square
\[
\begin{tikzcd}[column sep=large,row sep=large]
    X' \ar[r] \ar[d]
        & X \ar[d]\\
    S' \ar[r]
        & S,
\end{tikzcd}
\]
there is a canonical equivalence
\[
    \boxed{
    X'_{\dR/S'}
    \xrightarrow{\simeq}
    X_{\dR/S}\times_SS'.
    }
\]
These equivalences satisfy identity, composition, horizontal and vertical
pasting, and all higher base-change coherences.
\end{theorem}

\begin{proof}
Left exactness gives
\[
    L_{\dR}X'
    \simeq
    L_{\dR}X\times_{L_{\dR}S}L_{\dR}S'.
\]
Hence
\[
\begin{aligned}
    X'_{\dR/S'}
    &=S'\times_{L_{\dR}S'}L_{\dR}X'\\
    &\simeq
    S'\times_{L_{\dR}S'}
    \left(L_{\dR}X\times_{L_{\dR}S}L_{\dR}S'\right)\\
    &\simeq
    S'\times_{L_{\dR}S}L_{\dR}X\\
    &\simeq
    X_{\dR/S}\times_SS'.
\end{aligned}
\]
All coherence statements are inherited from functoriality of finite limits
and the naturality and idempotence coherences of $L_{\dR}$.
\end{proof}

\section{\'Etale rigidity and smooth effectivity of relative de Rham objects}
\label{sec:infinitesimal-dr-smooth-etale}

\begin{theorem}[Relative de Rham rigidity of \'etale morphisms]
\label{thm:infinitesimal-etale-dr-rigidity}
If $f:X\to S$ is spectral \'etale, then
\[
    \eta_{X/S}:X\xrightarrow{\simeq}X_{\dR/S}.
\]
\end{theorem}

\begin{proof}
By Theorem~\ref{thm:infinitesimal-etale-formally-etale}, for every square-zero
thickening $U\to T$ over $S$ the restriction map
\[
    \Map_S(T,X)\longrightarrow\Map_S(U,X)
\]
is an equivalence.  The slice locality criterion of
Theorem~\ref{thm:infinitesimal-slice-localization} therefore shows that
$X\to S$ is already local.  Its localization unit is consequently an
equivalence.
\end{proof}

\begin{lemma}[Local splitting of effective epimorphisms over affines]
\label{lem:infinitesimal-effective-epi-local-splitting}
Let $E\to h_T$ be an effective epimorphism in $\mathcal X_{\Sp}$ with $T$
affine.  Then there exists a finite affine spectral Nisnevich basis cover
\[
    \{T_\alpha\to T\}
\]
such that every pullback $E\times_{h_T}h_{T_\alpha}\to h_{T_\alpha}$ admits
a section.
\end{lemma}

\begin{proof}
An effective epimorphism in a sheaf $\infty$-topos is locally surjective.  The
sieve of maps $T'\to T$ over which the pullback admits a section is therefore
a covering sieve.  By Proposition~\ref{prop:infinitesimal-affine-nis-basis},
the sieve contains a finite affine basis cover, which supplies the required
local sections.
\end{proof}

\begin{lemma}[Image-locality criterion]
\label{lem:infinitesimal-image-locality}
Let $X\to S$ be any morphism and factor its relative de Rham unit in the slice
as
\[
    X\xrightarrow{e}Q\xrightarrow{m}Y,
    \qquad
    Y:=X_{\dR/S},
\]
with $e$ an effective epimorphism and $m$ a monomorphism.  Suppose that for
every affine spectral square-zero immersion $U\to T$ over $S$ and every
$q_0:U\to Q$, the unique extension
\[
    a:T\longrightarrow Y
\]
of $mq_0$ supplied by locality of $Y$ factors Nisnevich-locally through $Q$.
Then $Q$ is $L_{\dR,/S}$-local and $m:Q\to Y$ is an equivalence.
\end{lemma}

\begin{proof}
Fix an affine square-zero immersion $U\to T$ over $S$.  Since $Y$ is local,
restriction is an equivalence
\[
    \Map_S(T,Y)\xrightarrow{\simeq}\Map_S(U,Y).
\]
For $a:T\to Y$, let $a_0$ denote its restriction and define
\[
    \Phi_T(a)
    :=
    \fib_a\bigl(\Map_S(T,Q)\to\Map_S(T,Y)\bigr),
\]
\[
    \Phi_U(a_0)
    :=
    \fib_{a_0}\bigl(\Map_S(U,Q)\to\Map_S(U,Y)\bigr).
\]
Because $m$ is a monomorphism, both fibres are $(-1)$-truncated spaces and are
therefore either empty or contractible.  Restriction gives
\[
    \Phi_T(a)\longrightarrow\Phi_U(a_0).
\]
If the source is nonempty, then the target is nonempty.  Conversely, if
$\Phi_U(a_0)$ is nonempty, choose the corresponding factorization
$q_0:U\to Q$.  By the hypothesis, $a$ factors Nisnevich-locally through
$Q$.  On overlaps, two local factorizations have the same composite to $Y$;
since $Q\to Y$ is a monomorphism, the space of identifications between them
is contractible.  The local factorizations therefore descend to a global
factorization $T\to Q$.  Hence $\Phi_T(a)$ is nonempty.

Thus $\Phi_T(a)\to\Phi_U(a_0)$ is an equivalence for every $a$.  Together
with the equivalence on the $Y$-mapping spaces, this shows that
\[
\begin{tikzcd}[column sep=large,row sep=large]
    \Map_S(T,Q) \ar[r] \ar[d]
        & \Map_S(T,Y) \ar[d,"\simeq"]\\
    \Map_S(U,Q) \ar[r]
        & \Map_S(U,Y)
\end{tikzcd}
\]
is a pullback square.  Hence
\[
    \Map_S(T,Q)\xrightarrow{\simeq}\Map_S(U,Q).
\]
The slice locality criterion shows that $Q$ is $L_{\dR,/S}$-local.

Apply the left-exact localization $L_{\dR,/S}$ to
\[
    X\xrightarrow{e}Q\xrightarrow{m}Y.
\]
The localization preserves the geometric realization of the \v{C}ech nerve
of $e$, because it is a left adjoint, and preserves that \v{C}ech nerve
itself, because it is left exact.  Hence $L_{\dR,/S}e$ remains an effective
epimorphism.  Left exactness preserves monomorphisms, so
$L_{\dR,/S}m$ is a monomorphism.  Their composite is the localization of the
relative de Rham unit and is an equivalence.  Therefore $L_{\dR,/S}m$ is both
an effective epimorphism and a monomorphism and is an equivalence.  Both $Q$
and $Y$ are local, so $m$ itself is an equivalence.
\end{proof}

\begin{theorem}[Smooth effectivity of the relative de Rham unit]
\label{thm:infinitesimal-smooth-dr-effectivity}
If $f:X\to S$ is a smooth morphism of spectral schemes, then
\[
    \boxed{
    \eta_{X/S}:X\longrightarrow X_{\dR/S}
    }
\]
is an effective epimorphism in the slice $\mathcal X_{\Sp,/S}$.
\end{theorem}

\begin{proof}
Factor the relative de Rham unit as
\[
    X\xrightarrow{e}Q\xrightarrow{m}Y,
    \qquad
    Y:=X_{\dR/S},
\]
with $e$ an effective epimorphism and $m$ a monomorphism.  We verify the
hypothesis of Lemma~\ref{lem:infinitesimal-image-locality}.

Let $U\to T$ be an affine spectral square-zero immersion over $S$, and let
$q_0:U\to Q$.  Pulling back $e$ gives an effective epimorphism
\[
    U\times_QX\longrightarrow U.
\]
By Lemma~\ref{lem:infinitesimal-effective-epi-local-splitting}, after a finite
affine Nisnevich basis cover $\{U_\alpha\to U\}$ the map $q_0$ lifts to maps
\[
    x_{0,\alpha}:U_\alpha\longrightarrow X.
\]
One-stage \'etale rigidity transports this cover to a finite affine
Nisnevich basis cover $\{T_\alpha\to T\}$ with
\[
    U_\alpha\simeq U\times_TT_\alpha.
\]
Applying Lemma~\ref{lem:infinitesimal-smooth-one-stage-lifting} to the maps
$x_{0,\alpha}$ and replacing the family by a finite composite refinement
gives extensions
\[
    x_\alpha:T_\alpha\longrightarrow X
\]
with specified homotopies on the closed stages.

Since $Y$ is local, the composite $mq_0:U\to Y$ has an extension
\[
    a:T\longrightarrow Y
\]
through a contractible space of choices.  On $U_\alpha$, the restriction of
$a$ agrees, through the retained homotopies, with the composite
$U_\alpha\to X\to Y$.  Locality of $Y$ gives an equivalence
\[
    \Map_S(T_\alpha,Y)
    \xrightarrow{\simeq}
    \Map_S(U_\alpha,Y),
\]
so agreement extends through a contractible space of homotopies to
$T_\alpha$.  Therefore $a$ factors Nisnevich-locally through $X$, hence
through $Q$.

Lemma~\ref{lem:infinitesimal-image-locality} now gives $Q\simeq Y$.  Thus the
effective-image factorization of $\eta_{X/S}$ has monomorphic factor an
equivalence, so $\eta_{X/S}$ is an effective epimorphism.
\end{proof}

\section{Local linear models and tubular linearization}
\label{sec:infinitesimal-local-linear-models}

\begin{proposition}[Brave-new affine spaces]
\label{prop:infinitesimal-affine-spaces}
For every spectral scheme $S$ and $n\geq0$, the brave-new affine space
\[
    \mathbb A^n_S\longrightarrow S
\]
is smooth and
\[
    L_{\mathbb A^n_S/S}
    \simeq
    \mathcal O_{\mathbb A^n_S}^{\oplus n}.
\]
\end{proposition}

\begin{proof}
Affine-locally this is Corollary~2.3.2: for
\[
    A\{t_1,\ldots,t_n\}=\Sym_A(A^{\oplus n})
\]
one has
\[
    L_{A\{t_1,\ldots,t_n\}/A}
    \simeq
    A\{t_1,\ldots,t_n\}^{\oplus n}.
\]
The free algebra is a finite cell algebra, hence of finite presentation, so
the affine-space projection is smooth.  The affine calculation is compatible
with base change and therefore glues over an arbitrary spectral scheme $S$.
\end{proof}

\begin{proposition}[Relative spectral vector spaces]
\label{prop:infinitesimal-vector-spaces}
Let $E$ be a finite locally free quasi-coherent module on a spectral scheme
$S$.  The relative linear space
\[
    \pi:V_S(E)\longrightarrow S
\]
is smooth and its relative cotangent complex is canonically
\[
    L_{V_S(E)/S}\simeq\pi^*E.
\]
\end{proposition}

\begin{proof}
The assertion is Zariski-local on $S$.  After trivializing $E$ it reduces to
Proposition~\ref{prop:infinitesimal-affine-spaces}; the local identifications
glue by pullback compatibility of the relative symmetric algebra and the
cotangent complex.
\end{proof}

\begin{theorem}[Local standard smooth charts]
\label{thm:infinitesimal-local-coordinates}
Let $f:X\to S$ be a smooth morphism of spectral schemes and let $x\in|X|$.
After shrinking $X$ and $S$ Zariski-locally around $x$ and $f(x)$, there is a
factorization
\[
    X\longrightarrow\mathbb A^n_S\longrightarrow S
\]
in which the first morphism is spectral \'etale.
\end{theorem}

\begin{proof}
This is the local standard-form theorem of
Chapter~\ref{chap:smooth-spectral-geometry}.
\end{proof}

\begin{theorem}[Tubular linearization along a smooth affine section]
\label{thm:infinitesimal-tubular-linearization}
Let
\[
    p:X\longrightarrow S
\]
be a smooth affine morphism of spectral schemes and let
\[
    s:S\longrightarrow X
\]
be a section.  Write
\[
    N^*_{S/X}:=L_{S/X}[-1]
\]
for its conormal bundle.  There exists a Zariski cover
\[
    \{S_\alpha\longrightarrow S\}_{\alpha\in A}
\]
such that, with $X_\alpha:=X\times_SS_\alpha$, there is for every $\alpha$ an
open neighbourhood
\[
    X_\alpha^\circ\subseteq X_\alpha
\]
of the image of $s|_{S_\alpha}$ and an \'etale morphism
\[
    \phi_\alpha:
    X_\alpha^\circ
    \longrightarrow
    V_{S_\alpha}(N^*_{S/X}|_{S_\alpha})
\]
carrying the section to the zero section and inducing the identity on
conormal bundles.  For every finite spectral infinitesimal thickening
$U\to T$ over $S_\alpha$, the canonical map
\[
\begin{aligned}
    \Map_{S_\alpha}(T,X_\alpha^\circ)
    \longrightarrow{}
    &\Map_{S_\alpha}
      \left(T,V_{S_\alpha}(N^*_{S/X}|_{S_\alpha})\right)\\
    &\times_{
      \Map_{S_\alpha}
      \left(U,V_{S_\alpha}(N^*_{S/X}|_{S_\alpha})\right)}
      \Map_{S_\alpha}(U,X_\alpha^\circ)
\end{aligned}
\]
is an equivalence.
\end{theorem}

\begin{proof}
The existence of the local \'etale tubular comparison, including the
conormal identification and the smooth-affine scope, is the spectral tubular
neighbourhood theorem of Chapter~\ref{chap:smooth-spectral-geometry}.  Apply
the pointwise formal-\'etale comparison of
Theorem~\ref{thm:infinitesimal-etale-formally-etale} to $\phi_\alpha$.  The
displayed relative mapping-space comparison is the homotopy fibre of that
absolute comparison over the fixed structural map $T\to S_\alpha$, and is
therefore an equivalence.
\end{proof}

\section{The reduced-classical shadow and spectral refinement}
\label{sec:infinitesimal-reduced-comparison}

\begin{definition}[Reduced-classical test functor]
\label{def:infinitesimal-reduced-test}
For a connective commutative ring spectrum $A$, put
\[
    A_{\mathrm{red}}^{\mathrm{cl}}
    :=
    H\bigl(\pi_0(A)_{\mathrm{red}}\bigr).
\]
Let
\[
    r:\mathcal C\longrightarrow\mathcal C,
    \qquad
    r(\Spec A)=\Spec A_{\mathrm{red}}^{\mathrm{cl}}.
\]
For $F\in\mathcal X_{\Sp}$ define
\[
    (D_{\mathrm{red}}F)(\Spec A)
    :=
    F(\Spec A_{\mathrm{red}}^{\mathrm{cl}}).
\]
\end{definition}

\begin{lemma}[Reduction preserves affine Nisnevich basis covers]
\label{lem:infinitesimal-reduction-nis}
Reduced classical truncation carries affine spectral \'etale morphisms to
ordinary \'etale morphisms, affine spectral open immersions to the
corresponding ordinary open immersions, and finite affine spectral Nisnevich
basis covers to finite ordinary affine Nisnevich covers regarded as discrete
spectral families.
\end{lemma}

\begin{proof}
If $A\to B$ is spectral \'etale, then $\pi_0(B)$ is \'etale over
$\pi_0(A)$ and
\[
    \pi_n(B)
    \simeq
    \pi_n(A)\otimes_{\pi_0(A)}\pi_0(B).
\]
Base change to the discrete reduced algebra $H(\pi_0(A)_{\mathrm{red}})$ is
therefore discrete, with degree-zero ring
\[
    \pi_0(B)\otimes_{\pi_0(A)}\pi_0(A)_{\mathrm{red}}.
\]
An \'etale algebra over a reduced ordinary ring is reduced, so the last ring
is canonically $\pi_0(B)_{\mathrm{red}}$.  Thus
\[
    B\otimes_AH(\pi_0(A)_{\mathrm{red}})
    \simeq
    H(\pi_0(B)_{\mathrm{red}}).
\]
Reduction commutes with ordinary localization, hence with affine open
immersions, and it does not change underlying topological spaces or residue
fields.

It follows that reduced classical truncation carries every spectral
Nisnevich distinguished square to the corresponding ordinary distinguished
square and therefore carries every finite cover in the pretopology
$\mathcal P_T$ of Proposition~\ref{prop:infinitesimal-affine-nis-basis} to an
ordinary Nisnevich cover.  Now let $\{T_\alpha\to T\}$ be an affine basis
cover in the sense of Definition~\ref{def:infinitesimal-affine-nis-basis}.
Its generated covering sieve contains a finite $\mathcal P_T$-cover, hence
that cover is a refinement of the given family.  After reduction the
refinement remains a Nisnevich cover, so the sieve generated by the reduced
family is covering as well.  Thus reduction carries every finite affine basis
cover to an ordinary affine Nisnevich covering family.
\end{proof}

\begin{lemma}[Reduced \'etale base change and \v{C}ech compatibility]
\label{lem:infinitesimal-reduction-cech}
Let $A\to B$ be a spectral \'etale morphism of connective commutative ring
spectra.  There is a canonical equivalence
\[
    B\otimes_AH(\pi_0(A)_{\mathrm{red}})
    \xrightarrow{\simeq}
    H(\pi_0(B)_{\mathrm{red}}).
\]
Equivalently, the square
\[
\begin{tikzcd}[column sep=large,row sep=large]
    r(\Spec B) \ar[r] \ar[d]
        & \Spec B \ar[d]\\
    r(\Spec A) \ar[r]
        & \Spec A
\end{tikzcd}
\]
is Cartesian.  Consequently, for affine spectral \'etale $A$-algebras
$B_0,\ldots,B_n$, there is a canonical equivalence
\[
\begin{aligned}
    r\!\left(
        \Spec(B_0\otimes_A\cdots\otimes_AB_n)
    \right)
    \simeq
    r(\Spec B_0)
    \times_{r(\Spec A)}
    \cdots
    \times_{r(\Spec A)}
    r(\Spec B_n).
\end{aligned}
\]
Thus reduced classical truncation carries the \v{C}ech nerve of every finite
affine spectral \'etale family to the \v{C}ech nerve of its reduced family.
\end{lemma}

\begin{proof}
Put
\[
    R:=\pi_0(A),
    \qquad
    N:=\sqrt{0_R}.
\]
Spectral \'etaleness and the flat homotopy-group formula give
\[
    \pi_k(B)
    \simeq
    \pi_k(A)\otimes_R\pi_0(B).
\]
After base change along $A\to H(R/N)$, the resulting algebra is spectral
\'etale over the discrete ring spectrum $H(R/N)$ and has no nonzero higher
homotopy groups.  Hence
\[
    B\otimes_AH(R/N)
    \simeq
    H\!\left(\pi_0(B)\otimes_RR/N\right)
    =
    H\!\left(\pi_0(B)/N\pi_0(B)\right).
\]

The quotient $\pi_0(B)/N\pi_0(B)$ is \'etale over the reduced ring $R/N$ and
is therefore reduced.  Hence
\[
    \sqrt{0_{\pi_0(B)}}\subseteq N\pi_0(B).
\]
Conversely, every element of $N$ maps to a nilpotent element of $\pi_0(B)$,
and the nilradical of a commutative ring is an ideal.  Therefore
\[
    N\pi_0(B)\subseteq\sqrt{0_{\pi_0(B)}}.
\]
Thus
\[
    N\pi_0(B)=\sqrt{0_{\pi_0(B)}},
\]
which proves the first equivalence and the Cartesian square.

Every iterated tensor product
\[
    B_0\otimes_A\cdots\otimes_AB_n
\]
is again spectral \'etale over $A$.  Applying the first equivalence and
associativity of base change gives
\[
\begin{aligned}
&H(\pi_0(B_0\otimes_A\cdots\otimes_AB_n)_{\mathrm{red}})\\
&\qquad\simeq
(B_0\otimes_A\cdots\otimes_AB_n)
\otimes_AH(\pi_0(A)_{\mathrm{red}})\\
&\qquad\simeq
H(\pi_0(B_0)_{\mathrm{red}})
\otimes_{H(\pi_0(A)_{\mathrm{red}})}
\cdots
\otimes_{H(\pi_0(A)_{\mathrm{red}})}
H(\pi_0(B_n)_{\mathrm{red}}).
\end{aligned}
\]
Passing to affine spectra gives the displayed fibre-product equivalence in
every simplicial degree.  The face and degeneracy maps are inherited from
the same base-change functoriality, so the entire \v{C}ech nerves correspond.
\end{proof}

\begin{proposition}[Reduced-classical localization]
\label{prop:infinitesimal-reduced-localization}
The functor $D_{\mathrm{red}}$ preserves spectral Nisnevich sheaves and defines
an accessible left-exact localization
\[
    D_{\mathrm{red}}:
    \mathcal X_{\Sp}\longrightarrow\mathcal X_{\Sp}^{\mathrm{red}}.
\]
Its unit is induced by
\[
    \Spec A_{\mathrm{red}}^{\mathrm{cl}}
    \longrightarrow
    \Spec A,
\]
and $F$ is $D_{\mathrm{red}}$-local if and only if
\[
    F(\Spec A)
    \longrightarrow
    F(\Spec A_{\mathrm{red}}^{\mathrm{cl}})
\]
is an equivalence for every connective $A$.
\end{proposition}

\begin{proof}
Let $F$ be a spectral Nisnevich sheaf and let
\[
    \{T_\alpha\to T\}_{\alpha\in A}
\]
be a finite affine Nisnevich basis cover.  By
Lemma~\ref{lem:infinitesimal-reduction-nis}, the reduced family
\[
    \{r(T_\alpha)\to r(T)\}_{\alpha\in A}
\]
is a Nisnevich covering family.  By
Lemma~\ref{lem:infinitesimal-reduction-cech}, reduction identifies every
finite intersection in the \v{C}ech nerve:
\[
    r(T_{\alpha_0}\times_T\cdots\times_TT_{\alpha_n})
    \simeq
    r(T_{\alpha_0})
    \times_{r(T)}
    \cdots
    \times_{r(T)}
    r(T_{\alpha_n}).
\]
Therefore the descent diagram for $D_{\mathrm{red}}F$ on the original cover
is exactly the descent diagram for $F$ on the reduced cover.  Since $F$ is a
sheaf, this diagram is a limit diagram.  Thus $D_{\mathrm{red}}F$ is a
spectral Nisnevich sheaf.

If
\[
    \iota:\mathcal X_{\Sp}
    \hookrightarrow
    \operatorname{PSh}(\mathcal C)
\]
is the inclusion and $a_{\Nis}$ is sheafification, then
\[
    D_{\mathrm{red}}
    \simeq
    a_{\Nis}\circ r^*\circ\iota.
\]
The sheafification is redundant on the image by the preceding paragraph, but
this expression proves accessibility.  Since $D_{\mathrm{red}}$ is
precomposition by $r$ on sheaves and limits of sheaves are computed
objectwise, it preserves all limits and in particular is left exact.

The natural maps
$\Spec A_{\mathrm{red}}^{\mathrm{cl}}\to\Spec A$ give a coaugmentation
$\id\to D_{\mathrm{red}}$.  Applying reduction twice changes nothing, so both
composites
\[
    D_{\mathrm{red}}
    \rightrightarrows
    D_{\mathrm{red}}D_{\mathrm{red}}
\]
are equivalences.  If $Y$ is fixed by $D_{\mathrm{red}}$, precomposition with
the unit gives
\[
    \Map(D_{\mathrm{red}}X,Y)
    \longrightarrow
    \Map(X,Y).
\]
Applying $D_{\mathrm{red}}$ to a map $X\to Y$ and using the fixed-object
equivalence $D_{\mathrm{red}}Y\simeq Y$ constructs an inverse.  Thus
$D_{\mathrm{red}}$ is the reflector onto its fixed objects.  The pointwise
criterion is precisely the condition that the unit be an equivalence.
\end{proof}

\begin{lemma}[Discrete square-zero recognition]
\label{lem:infinitesimal-discrete-square-zero-recognition}
Let $R'\to R$ be a surjection of ordinary commutative rings with square-zero
kernel.  Then
\[
    HR'\longrightarrow HR
\]
is a structured spectral square-zero extension.
\end{lemma}

\begin{proof}
Let $I=\ker(R'\to R)$.  The fibre of
\[
    HR'\longrightarrow HR
\]
is the discrete $HR$-module $HI$, concentrated in degree zero, and the
multiplication on this fibre is null because $I^2=0$.  In the terminology of
the bounded small-extension recognition theorem recalled in
Chapter~\ref{chap:cotangent-spectral-deformation}, the morphism
$HR'\to HR$ is therefore a $0$-small extension.  That recognition theorem
supplies a classifying derivation
\[
    \delta:L_{HR}\longrightarrow HI[1]
\]
and a specified equivalence over $HR$
\[
    HR'
    \xrightarrow{\simeq}
    HR\times_{HR\oplus HI[1]}HR,
\]
where the two maps from $HR$ to $HR\oplus HI[1]$ are the derivation
classified by $\delta$ and the zero derivation.  This is exactly a structured
spectral square-zero presentation in the sense of
Chapter~\ref{chap:cotangent-spectral-deformation}.
\end{proof}

\begin{theorem}[Comparison of de Rham and reduced-classical modalities]
\label{thm:infinitesimal-reduced-comparison}
Every spectral square-zero immersion is inverted by $D_{\mathrm{red}}$.
Hence the initiality of $L_{\dR}$ supplies, after viewing the de Rham reflector
through its fully faithful inclusion $\iota_{\dR}$, a canonical transformation
of left-exact endofunctors of $\mathcal X_{\Sp}$
\[
    \boxed{
    \kappa:\iota_{\dR}L_{\dR}\longrightarrow D_{\mathrm{red}}.
    }
\]
We suppress $\iota_{\dR}$ from the notation below and write this comparison
as $L_{\dR}\to D_{\mathrm{red}}$.  It induces, for every $X\to S$, a map
compatible with pullback base change
\[
    X_{\dR/S}
    \longrightarrow
    S\times_{D_{\mathrm{red}}S}D_{\mathrm{red}}X.
\]
\end{theorem}

\begin{proof}
Let $A'\to A$ be a spectral square-zero extension.  To show that
$D_{\mathrm{red}}$ inverts the corresponding generator
\[
    h_{\Spec A}\longrightarrow h_{\Spec A'},
\]
evaluate after $D_{\mathrm{red}}$ on an arbitrary affine test
$T=\Spec C$ and put
\[
    R:=\pi_0(C)_{\mathrm{red}}.
\]
Since $HR$ is discrete, there are canonical equivalences
\[
\begin{aligned}
    (D_{\mathrm{red}}h_{\Spec A})(T)
    &\simeq
    \Map_{\CAlg(\Sp)}(A,HR)
    \simeq
    \operatorname{Hom}_{\operatorname{CRing}}(\pi_0(A),R),\\
    (D_{\mathrm{red}}h_{\Spec A'})(T)
    &\simeq
    \Map_{\CAlg(\Sp)}(A',HR)
    \simeq
    \operatorname{Hom}_{\operatorname{CRing}}(\pi_0(A'),R).
\end{aligned}
\]
The map
\[
    \pi_0(A')\longrightarrow\pi_0(A)
\]
is surjective with square-zero kernel $I$.  Every ring map
$\pi_0(A')\to R$ kills $I$, because the image of every $x\in I$ satisfies
$f(x)^2=0$ and $R$ is reduced.  Hence every such map factors uniquely through
$\pi_0(A)$.  Restriction therefore induces an equivalence
\[
    \operatorname{Hom}_{\operatorname{CRing}}(\pi_0(A),R)
    \xrightarrow{\simeq}
    \operatorname{Hom}_{\operatorname{CRing}}(\pi_0(A'),R).
\]
Thus
\[
    D_{\mathrm{red}}h_{\Spec A}
    \xrightarrow{\simeq}
    D_{\mathrm{red}}h_{\Spec A'}
\]
objectwise on affine tests, and $D_{\mathrm{red}}$ inverts every square-zero
generator.

The universal property of
Theorem~\ref{thm:infinitesimal-dr-localization} now gives $\kappa$.  The
relative comparison is obtained by applying the two localizations to the
cospan defining the relative de Rham object and taking the induced map between
pullbacks.  Base-change coherence follows from left exactness and functoriality
of finite limits.
\end{proof}

\begin{proposition}[Finite Postnikov and nilpotent collapse]
\label{prop:infinitesimal-finite-collapse}
Let $A$ be an eventually coconnective connective commutative ring spectrum and
suppose that the nilradical of $\pi_0(A)$ is nilpotent.  Then the canonical
map
\[
    A\longrightarrow H\bigl(\pi_0(A)_{\mathrm{red}}\bigr)
\]
admits a finite factorization by structured square-zero extensions.
Consequently
\[
    L_{\dR}h_{\Spec A}
    \simeq
    L_{\dR}h_{\Spec H(\pi_0(A)_{\mathrm{red}})},
\]
and every $L_{\dR}$-local $F$ satisfies
\[
    F(\Spec A)
    \simeq
    F\bigl(\Spec H(\pi_0(A)_{\mathrm{red}})\bigr).
\]
\end{proposition}

\begin{proof}
Choose $N$ with $A\simeq\tau_{\leq N}A$.  The multiplicative Postnikov tower
\[
    A\to\tau_{\leq N-1}A\to\cdots\to\tau_{\leq0}A\simeq H\pi_0(A)
\]
is finite, and every transition is a structured square-zero extension by the
multiplicative Postnikov theorem of
Chapter~\ref{chap:cotangent-spectral-deformation}.

Let $R=\pi_0(A)$, let $I\subset R$ be the nilradical, and choose $m$ with
$I^m=0$.  The quotient tower
\[
    R=R/I^m\to R/I^{m-1}\to\cdots\to R/I=R_{\mathrm{red}}
\]
has square-zero successive kernels because
\[
    (I^{j-1})^2=I^{2j-2}\subseteq I^j
    \qquad(j\geq2).
\]
By Lemma~\ref{lem:infinitesimal-discrete-square-zero-recognition}, applying
$H$ makes every stage a structured spectral square-zero extension.
Concatenating the two finite towers proves the first assertion.  The final
two equivalences follow because $L_{\dR}$ inverts every stage and local
objects send de Rham equivalences to equivalences of values.
\end{proof}

\begin{remark}[Infinite Postnikov towers]
\label{rem:infinitesimal-infinite-postnikov}
For a general connective $E_\infty$-ring, the passage from $A$ to $H\pi_0(A)$
uses an infinite inverse Postnikov tower.  The left adjoint $L_{\dR}$ does not
automatically commute with that inverse limit.  The intrinsic modality is
therefore defined from finite square-zero generators rather than from the
pointwise formula $A\mapsto H(\pi_0(A)_{\mathrm{red}})$.
\end{remark}

\begin{definition}[Postnikov completeness on affine tests]
\label{def:infinitesimal-postnikov-complete}
An $L_{\dR}$-local object $F$ is \emph{Postnikov complete on affine tests} if
for every connective commutative ring spectrum $A$ the canonical map
\[
    F(\Spec A)
    \longrightarrow
    \lim_nF(\Spec\tau_{\leq n}A)
\]
is an equivalence.
\end{definition}

\begin{definition}[Nilradical continuity on affine tests]
\label{def:infinitesimal-nilradical-continuity}
An $L_{\dR}$-local object $F$ is \emph{nilradical-continuous on affine tests}
if, for every ordinary commutative ring $R$ with nilradical $N$, the canonical
map
\[
    \colim_{\substack{J\subseteq N\\J\text{ finitely generated}}}
    F(\Spec H(R/J))
    \longrightarrow
    F(\Spec H(R/N))
\]
is an equivalence.
\end{definition}

\begin{theorem}[Convergence criterion for reduced-classical locality]
\label{thm:infinitesimal-reduced-convergence}
Let $F\in\mathcal X_{\Sp}$ be $L_{\dR}$-local.  Then the following are
equivalent:
\begin{enumerate}
\item $F$ is $D_{\mathrm{red}}$-local;
\item $F$ is Postnikov complete and nilradical-continuous on affine tests.
\end{enumerate}
Consequently $D_{\mathrm{red}}$ is the further left-exact localization of the
intrinsic de Rham topos imposing exactly these two convergence conditions.
\end{theorem}

\begin{proof}
Assume first that $F$ is $D_{\mathrm{red}}$-local.  The rings $A$,
$\tau_{\leq n}A$, and $H(\pi_0(A)_{\mathrm{red}})$ have the same reduced
classical truncation.  The entire Postnikov evaluation diagram is therefore
naturally equivalent to the constant diagram with value
\[
    F\bigl(\Spec H(\pi_0(A)_{\mathrm{red}})\bigr),
\]
so the inverse-limit comparison is an equivalence.  Similarly, if $J$ is a
finitely generated ideal contained in the nilradical $N$ of an ordinary ring
$R$, then $R/J$ and $R/N$ have the same reduction.  The
nilradical-continuity diagram is a constant diagram indexed by the filtered
poset of finitely generated ideals contained in $N$, whose nerve is
contractible.  Hence its colimit is the same value.

Conversely suppose that the two continuity properties hold.  For every $n$,
the map
\[
    \tau_{\leq n}A\longrightarrow H\pi_0(A)
\]
is a finite composite of multiplicative Postnikov square-zero extensions.
Locality gives compatible equivalences
\[
    F(\Spec\tau_{\leq n}A)
    \simeq
    F(\Spec H\pi_0(A)).
\]
Postnikov completeness therefore yields
\[
    F(\Spec A)
    \simeq
    F(\Spec H\pi_0(A)).
\]

Put $R=\pi_0(A)$ and let $N$ be its nilradical.  Every finitely generated
ideal $J\subseteq N$ is nilpotent: if $J=(x_1,\ldots,x_r)$ and
$x_i^{n_i}=0$, then
\[
    J^{n_1+\cdots+n_r}=0.
\]
The quotient $R\to R/J$ therefore factors through a finite tower with
square-zero successive kernels.  By
Lemma~\ref{lem:infinitesimal-discrete-square-zero-recognition} and de Rham
locality,
\[
    F(\Spec HR)
    \xrightarrow{\simeq}
    F(\Spec H(R/J))
\]
compatibly as $J$ varies.  Nilradical continuity gives
\[
    F(\Spec HR)
    \simeq
    F(\Spec H(R/N)).
\]
Combining the two equivalences gives
\[
    F(\Spec A)
    \simeq
    F(\Spec A_{\mathrm{red}}^{\mathrm{cl}}),
\]
so $F$ is $D_{\mathrm{red}}$-local.
\end{proof}

\begin{corollary}[Invertibility of the modality comparison]
\label{cor:infinitesimal-reduced-comparison-invertibility}
For every $G\in\mathcal X_{\Sp}$, the comparison
\[
    \kappa_G:L_{\dR}G\longrightarrow D_{\mathrm{red}}G
\]
is an equivalence if and only if $L_{\dR}G$ is Postnikov complete and
nilradical-continuous on affine tests.
\end{corollary}

\begin{proof}
The comparison identifies $D_{\mathrm{red}}G$ with the
$D_{\mathrm{red}}$-localization of the already $L_{\dR}$-local object
$L_{\dR}G$.  Apply
Theorem~\ref{thm:infinitesimal-reduced-convergence}.
\end{proof}

\begin{remark}[Precise scope of finite collapse]
\label{rem:infinitesimal-finite-collapse-scope}
Proposition~\ref{prop:infinitesimal-finite-collapse} says that the map of
representables
\[
    h_{\Spec H(\pi_0(A)_{\mathrm{red}})}
    \longrightarrow
    h_{\Spec A}
\]
is an $L_{\dR}$-equivalence under the stated boundedness and nilpotence
conditions.  Equivalently, every $L_{\dR}$-local object evaluates the two
tests identically.  It does not assert that the natural comparison
$L_{\dR}G\to D_{\mathrm{red}}G$ is an equivalence for arbitrary $G$ without
the convergence criterion above.
\end{remark}


\section{Stable intrinsic de Rham coefficients on the affine smooth basis}
\label{sec:infinitesimal-stable-dr-interface}

Fix a spectral base $S$ in the standing motivic scope.  Let
\[
    \operatorname{AffSm}^{\fp}_{/S}
    \subseteq
    \operatorname{Sm}^{\fp}_{/S}
\]
be the full subcategory spanned by smooth finite-presentation spectral schemes
whose underlying spectral scheme is affine.

\begin{lemma}[Affine smooth basis]
\label{lem:infinitesimal-affine-smooth-basis}
The full subcategory $\operatorname{AffSm}^{\fp}_{/S}$ is a basis for the
spectral Nisnevich topology on $\operatorname{Sm}^{\fp}_{/S}$.  Restriction
therefore induces an equivalence of stable $\infty$-categories
\[
    \boxed{
    \operatorname{Shv}_{\Nis}
    (\operatorname{Sm}^{\fp}_{/S};\Sp)
    \xrightarrow{\simeq}
    \operatorname{Shv}_{\Nis}
    (\operatorname{AffSm}^{\fp}_{/S};\Sp).
    }
\]
\end{lemma}

\begin{proof}
Every object $X\to S$ of $\operatorname{Sm}^{\fp}_{/S}$ is quasi-compact and
quasi-separated in the standing scope.  Chapter~\ref{chap:smooth-spectral-geometry}
constructs finite affine Zariski covers of such objects, and finite
intersections admit finite affine refinements.  A finite affine Zariski cover
is a spectral Nisnevich cover.  Hence every object of the smooth site admits a
cover by objects of $\operatorname{AffSm}^{\fp}_{/S}$, and the same assertion
holds after pullback to members and finite intersections of such covers.  The
comparison lemma for a basis of a Grothendieck topology therefore identifies
space-valued sheaves on the two sites.  Applying the same argument objectwise
to spectrum-valued presheaves, or equivalently applying
$\Omega^\infty\Sigma^n$ for all $n\in\mathbb Z$, gives the displayed stable
equivalence.
\end{proof}

For later use write
\[
    \mathcal N_S^{\Sp}
    :=
    \operatorname{Shv}_{\Nis}
    (\operatorname{Sm}^{\fp}_{/S};\Sp).
\]

\begin{definition}[Stable intrinsic de Rham coefficients]
\label{def:infinitesimal-stable-dr-coefficients}
Let
\[
    \mathcal D_{\dR}^{\mathrm{st}}(S)
    :=
    \operatorname{Sp}
    \bigl((\mathcal X_{\Sp,/S}^{\dR})_*\bigr),
\]
the stabilization of the pointed relative de Rham-local slice topos of
Theorem~\ref{thm:infinitesimal-slice-localization}.  Its objects are called
\emph{stable intrinsic de Rham coefficients over $S$}.
\end{definition}

Because $L_{\dR,/S}$ is accessible and left exact, pointed localization and
stabilization give an exact reflective localization
\[
    \operatorname{Sp}\bigl((\mathcal X_{\Sp,/S})_*\bigr)
    \longrightarrow
    \mathcal D_{\dR}^{\mathrm{st}}(S).
\]

\begin{lemma}[Pointed dependent sum in slices]
\label{lem:infinitesimal-pointed-dependent-sum}
Let $\mathcal X$ be an $\infty$-topos and let
\[
    f:T\longrightarrow S
\]
be a morphism.  Pullback on pointed slices admits a canonical left adjoint
\[
    \Sigma_{f,+}:
    (\mathcal X_{/T})_*
    \rightleftarrows
    (\mathcal X_{/S})_*:
    f^*.
\]
For a pointed object represented by a section
\[
    s:T\longrightarrow Y\longrightarrow T,
\]
one has
\[
    \boxed{
    \Sigma_{f,+}(Y,s)
    \simeq
    \left(
      Y\amalg_T S\longrightarrow S,\;
      S\longrightarrow Y\amalg_T S
    \right).
    }
\]
If $Y\to T$ is unpointed and
\[
    Y_+:=Y\amalg T
\]
denotes the freely pointed object of the slice, then
\[
    \boxed{
    \Sigma_{f,+}(Y_+)
    \simeq
    (\Sigma_fY)_+,
    }
\]
where $\Sigma_fY$ is $Y$ regarded over $S$ by composition with $f$.
For composable morphisms
\[
    U\xrightarrow{g}T\xrightarrow{f}S
\]
there are canonical coherent equivalences
\[
    \Sigma_{f,+}\Sigma_{g,+}
    \xrightarrow{\simeq}
    \Sigma_{fg,+}.
\]
After stabilization these give an adjunction
\[
    \Sigma_{f,+}^{\mathrm{st}}
    \dashv
    f^{*,\mathrm{st}}
\]
with the same free-pointing and composition compatibilities.
\end{lemma}

\begin{proof}
Let $(Z,z)$ be a pointed object of the slice over $S$.  A pointed morphism
\[
    \left(Y\amalg_TS,\;S\to Y\amalg_TS\right)
    \longrightarrow
    (Z,z)
\]
over $S$ is, by the universal property of the pushout, exactly a morphism
$Y\to Z$ over $S$ whose restriction along $s:T\to Y$ is the specified
composite $T\xrightarrow{f}S\xrightarrow{z}Z$.  This is precisely a pointed
morphism
\[
    (Y,s)\longrightarrow f^*(Z,z)
\]
in the slice over $T$.  The resulting mapping-space equivalence is natural
in both variables and proves $\Sigma_{f,+}\dashv f^*$.

For a freely pointed object,
\[
\begin{aligned}
    \Sigma_{f,+}(Y_+)
    &\simeq
    (Y\amalg T)\amalg_T S\\
    &\simeq
    Y\amalg S
    =
    (\Sigma_fY)_+.
\end{aligned}
\]
For a composable pair $U\xrightarrow gT\xrightarrow fS$, associativity of
pushouts gives
\[
    (Y\amalg_UT)\amalg_TS
    \simeq
    Y\amalg_US,
\]
naturally in the pointed object $Y$.  Functoriality of iterated pushouts
supplies the identity, composition, and higher associativity coherences.
The pullback functor $f^*$ is left exact and accessible, so it induces
$f^{*,\mathrm{st}}$ on spectrum objects.  Its stabilized left adjoint is
characterized by
\[
    \Sigma_{f,+}^{\mathrm{st}}\Sigma^\infty(Y,s)
    \simeq
    \Sigma^\infty\Sigma_{f,+}(Y,s),
\]
and the adjunction just proved identifies it as the left adjoint of
$f^{*,\mathrm{st}}$.  The free-pointing and composition equivalences pass
through this construction with the same coherences.
\end{proof}

\begin{proposition}[Intrinsic de Rham coefficient sheaf]
\label{prop:infinitesimal-stable-dr-coefficient-sheaf}
For $K\in\mathcal D_{\dR}^{\mathrm{st}}(S)$ and
$X\in\operatorname{AffSm}^{\fp}_{/S}$ define
\[
    \Gamma_{\dR,S}^{\mathrm{aff}}(K)(X)
    :=
    \map_{\mathcal D_{\dR}^{\mathrm{st}}(S)}
    \left(
      \Sigma^\infty_+L_{\dR,/S}h_X,
      K
    \right).
\]
This presheaf satisfies affine spectral Nisnevich descent and therefore extends
uniquely, through the equivalence of
Lemma~\ref{lem:infinitesimal-affine-smooth-basis}, to a spectrum-valued
Nisnevich sheaf
\[
    \boxed{
    \Gamma_{\dR,S}(K)\in\mathcal N_S^{\Sp}.
    }
\]
The assignment
\[
    \Gamma_{\dR,S}:
    \mathcal D_{\dR}^{\mathrm{st}}(S)
    \longrightarrow
    \mathcal N_S^{\Sp}
\]
is exact.

If $f:T\to S$ is smooth and of finite presentation, pullback on native slices
preserves relative de Rham-local objects and induces an exact functor
\[
    f_{\dR}^{*,\mathrm{st}}:
    \mathcal D_{\dR}^{\mathrm{st}}(S)
    \longrightarrow
    \mathcal D_{\dR}^{\mathrm{st}}(T).
\]
There is a canonical coherent equivalence
\[
    \boxed{
    f^*_{\Nis}\Gamma_{\dR,S}
    \xrightarrow{\simeq}
    \Gamma_{\dR,T}f_{\dR}^{*,\mathrm{st}}.
    }
\]
\end{proposition}

\begin{proof}
Let
\[
    \{X_\alpha\to X\}_{\alpha\in A}
\]
be a finite affine spectral Nisnevich basis cover with $X$ affine and smooth
over $S$, and let $X_\bullet$ be its \v{C}ech nerve.  In the native affine
Nisnevich topos, subcanonicity and effective descent give a colimit diagram in
the slice over $S$,
\[
    \operatorname*{colim}_{[n]\in\Delta^{\op}}h_{X_n}
    \xrightarrow{\simeq}
    h_X.
\]
The functors $L_{\dR,/S}$, adjoining a disjoint basepoint, and suspension
spectrum are left adjoints and preserve this colimit.  Mapping the resulting
colimit into $K$ gives
\[
\begin{aligned}
    \Gamma_{\dR,S}^{\mathrm{aff}}(K)(X)
    &\simeq
    \lim_{[n]\in\Delta}
    \Gamma_{\dR,S}^{\mathrm{aff}}(K)(X_n).
\end{aligned}
\]
Thus the affine presheaf is a spectrum-valued Nisnevich sheaf.  The affine
smooth basis lemma gives its unique extension to $\mathcal N_S^{\Sp}$.
Mapping spectra out of a fixed object are exact in the second variable, and
the basis equivalence is exact, so $\Gamma_{\dR,S}$ is exact.

Now let $f:T\to S$ be smooth of finite presentation.  The relative de Rham
base-change equivalence of
Theorem~\ref{thm:infinitesimal-relative-dr-base-change} is equivalently a
canonical natural equivalence
\[
    f^*L_{\dR,/S}
    \xrightarrow{\simeq}
    L_{\dR,/T}f^*
\]
on native slices.  Hence slice pullback carries relative de Rham-local objects
to relative de Rham-local objects and descends to a colimit-preserving
left-exact functor between the local slice topoi.  After pointing and
stabilization it induces the exact functor $f_{\dR}^{*,\mathrm{st}}$.

Write
\[
    \iota_{\dR,S}^{\mathrm{st}}:
    \mathcal D_{\dR}^{\mathrm{st}}(S)
    \hookrightarrow
    \operatorname{Sp}((\mathcal X_{\Sp,/S})_*)
\]
and similarly $\iota_{\dR,T}^{\mathrm{st}}$ for the fully faithful right
adjoints of the stabilized relative de Rham localizations.  The preceding
base-change equivalence gives a canonical coherent identification
\[
    \iota_{\dR,T}^{\mathrm{st}}f_{\dR}^{*,\mathrm{st}}
    \simeq
    f^{*,\mathrm{st}}\iota_{\dR,S}^{\mathrm{st}}.
\]

Let $X\to T$ be affine, smooth, and of finite presentation.  Since $f$ is
smooth, the composite $X\to S$ is again a member of the affine smooth basis.
Using the two localization adjunctions and the stabilized pointed
dependent-sum adjunction of
Lemma~\ref{lem:infinitesimal-pointed-dependent-sum} gives canonical
equivalences
\[
\begin{aligned}
&\map_{\mathcal D_{\dR}^{\mathrm{st}}(T)}
 \left(
   \Sigma^\infty_+L_{\dR,/T}h_X,
   f_{\dR}^{*,\mathrm{st}}K
 \right)
\\
&\qquad\simeq
 \map_{\operatorname{Sp}((\mathcal X_{\Sp,/T})_*)}
 \left(
   \Sigma^\infty_+h_X,
   \iota_{\dR,T}^{\mathrm{st}}f_{\dR}^{*,\mathrm{st}}K
 \right)
\\
&\qquad\simeq
 \map_{\operatorname{Sp}((\mathcal X_{\Sp,/T})_*)}
 \left(
   \Sigma^\infty_+h_X,
   f^{*,\mathrm{st}}\iota_{\dR,S}^{\mathrm{st}}K
 \right)
\\
&\qquad\simeq
 \map_{\operatorname{Sp}((\mathcal X_{\Sp,/S})_*)}
 \left(
   \Sigma_{f,+}^{\mathrm{st}}\Sigma^\infty_+h_X,
   \iota_{\dR,S}^{\mathrm{st}}K
 \right)
\\
&\qquad\simeq
 \map_{\operatorname{Sp}((\mathcal X_{\Sp,/S})_*)}
 \left(
   \Sigma^\infty_+h_X,
   \iota_{\dR,S}^{\mathrm{st}}K
 \right)
\\
&\qquad\simeq
 \map_{\mathcal D_{\dR}^{\mathrm{st}}(S)}
 \left(
   \Sigma^\infty_+L_{\dR,/S}h_X,
   K
 \right).
\end{aligned}
\]
In the fourth line $h_X$ is the representable over $T$, while in the fifth it
is the same spectral scheme regarded over $S$ by composition; the equivalence
between those two lines is exactly the free-pointing formula of
Lemma~\ref{lem:infinitesimal-pointed-dependent-sum}.

These equivalences are natural in $X$ and $K$.  They identify the two sheaves
on the affine smooth basis, and the basis equivalence gives the displayed
smooth pullback comparison.  Identity, composition, and higher coherence are
inherited from the coherent pointed dependent sums and the base-change
coherence of
Theorem~\ref{thm:infinitesimal-relative-dr-base-change}.
\end{proof}

\section*{Part IV: Stable BNV differential coefficients and reconstruction}
\addcontentsline{toc}{section}{Part IV: Stable BNV differential coefficients and reconstruction}

\section{Stable spectral Nisnevich sheaves and the raw Thom spectrum}
\label{sec:infinitesimal-stable-nis-raw-thom}

Fix a spectral base $S$ in the standing motivic scope and write
\[
    \mathcal D_S:=\operatorname{Sm}^{\fp}_{/S}.
\]
Let
\[
    \mathcal P_S^{\mathrm{st}}
    :=
    \Fun(\mathcal D_S^{\op},\Sp)
\]
be the stable spectral presheaf category.  Equip it with the presentably
symmetric monoidal Day-convolution structure induced by fibre product over
$S$ and the smash product of spectra.  Thus the suspension-spectrum Yoneda
objects satisfy
\[
    \Sigma^\infty_+h_X\otimes\Sigma^\infty_+h_Y
    \simeq
    \Sigma^\infty_+h_{X\times_SY}.
\]

For a spectral Nisnevich distinguished square
\[
\begin{tikzcd}[column sep=large,row sep=large]
    W \ar[r] \ar[d] & V \ar[d]\\
    U \ar[r] & X,
\end{tikzcd}
\]
let
\[
    e_Q^{\mathrm{st}}:
    \Sigma^\infty_+h_U
    \amalg_{\Sigma^\infty_+h_W}
    \Sigma^\infty_+h_V
    \longrightarrow
    \Sigma^\infty_+h_X
\]
be the stable excision morphism, and let
\[
    e_\varnothing^{\mathrm{st}}:
    0\longrightarrow\Sigma^\infty_+h_\varnothing
\]
be the empty-object morphism.  Let $\mathcal E^{\mathrm{st}}_{\Nis,S}$ be a
set containing all integer suspensions of these maps.

\begin{proposition}[Stable spectral Nisnevich sheafification]
\label{prop:infinitesimal-stable-nisnevich}
The full subcategory
\[
    \mathcal N_S^{\Sp}
    :=
    \operatorname{Shv}_{\Nis}
    \left(
        \operatorname{Sm}^{\fp}_{/S};
        \Sp
    \right)
    \subseteq
    \mathcal P_S^{\mathrm{st}}
\]
is the local subcategory for $\mathcal E^{\mathrm{st}}_{\Nis,S}$.  Its
inclusion admits an accessible exact left adjoint
\[
    L^{\mathrm{st}}_{\Nis,S}:
    \mathcal P_S^{\mathrm{st}}
    \longrightarrow
    \mathcal N_S^{\Sp}.
\]
The class of stable Nisnevich equivalences is a tensor ideal, so
$L^{\mathrm{st}}_{\Nis,S}$ is strong symmetric monoidal and $\mathcal N_S^{\Sp}$
is a stable presentably symmetric monoidal $\infty$-category.

For $X\in\mathcal D_S$, write again
\[
    \Sigma^\infty_+h_X
    :=
    L^{\mathrm{st}}_{\Nis,S}\Sigma^\infty_+h_X
    \in\mathcal N_S^{\Sp}.
\]
Then
\[
    \mathcal G_S
    :=
    \left\{
        \Sigma^a\Sigma^\infty_+h_X
        \ \middle|\
        X\in\mathcal D_S,\ a\in\mathbb Z
    \right\}
\]
is a tensor-closed set of compact generators of $\mathcal N_S^{\Sp}$ and contains
the tensor unit.
\end{proposition}

\begin{proof}
For $E\in\mathcal P_S^{\mathrm{st}}$, mapping a suspension of
$e_Q^{\mathrm{st}}$ into $E$ gives the corresponding suspension of the map of
mapping spectra
\[
    E(X)
    \longrightarrow
    E(U)\times_{E(W)}E(V).
\]
Locality for all integer suspensions is therefore equivalent to this being an
equivalence of spectra for every distinguished square, together with
$E(\varnothing)\simeq0$.  Applying $\Omega^\infty\Sigma^n$ for all
$n\in\mathbb Z$ identifies these conditions with spectral Nisnevich descent.
The accessible localization theorem gives the displayed localization.  Since
the generating set is stable under suspension, the local subcategory is
stable and the localization is exact.

Tensoring $e_Q^{\mathrm{st}}$ with $\Sigma^\infty_+h_Y$ gives the stable
excision morphism associated with the base-changed distinguished square over
$X\times_SY$; tensoring $e_\varnothing^{\mathrm{st}}$ with a representable
again gives the empty-object generator.  Hence the generating equivalences
are stable under tensoring with suspension representables.  These objects
generate $\mathcal P_S^{\mathrm{st}}$ under colimits and tensor product
preserves colimits separately, so the strongly saturated class of stable
Nisnevich equivalences is a tensor ideal.  The symmetric monoidal
localization theorem makes $L^{\mathrm{st}}_{\Nis,S}$ strong symmetric
monoidal.

The stable Nisnevich sheaf conditions are finite-limit conditions and are
preserved by filtered colimits of spectra.  Hence the inclusion
$\mathcal N_S^{\Sp}\hookrightarrow\mathcal P_S^{\mathrm{st}}$ preserves filtered
colimits.  Stable Yoneda and the localization adjunction give
\[
    \map_{\mathcal N_S^{\Sp}}
    (\Sigma^\infty_+h_X,E)
    \simeq
    E(X).
\]
Thus every $\Sigma^a\Sigma^\infty_+h_X$ is compact.  If all mapping spectra
from these objects into $E$ vanish, then every $E(X)$ vanishes, so
$E\simeq0$.  Finally,
\[
    \Sigma^a\Sigma^\infty_+h_X
    \otimes
    \Sigma^b\Sigma^\infty_+h_Y
    \simeq
    \Sigma^{a+b}\Sigma^\infty_+h_{X\times_SY},
\]
so the set is tensor-closed, and the object attached to $X=S$ is the tensor
unit.
\end{proof}

\begin{definition}[Raw Thom spectrum]
\label{def:infinitesimal-raw-thom}
Let
\[
    \mathbb G_{m,S}
    :=
    \mathbb A^1_S\setminus0_S.
\]
Define the \emph{raw Thom spectrum of the trivial line} by
\[
    \boxed{
    \mathbb T_S^{\mathrm{raw}}
    :=
    \cofib\left(
        \Sigma^\infty_+h_{\mathbb G_{m,S}}
        \longrightarrow
        \Sigma^\infty_+h_{\mathbb A^1_S}
    \right)
    \in\mathcal N_S^{\Sp}.
    }
\]
The superscript ``raw'' records that no affine-line localization has yet been
imposed.
\end{definition}

\begin{proposition}[Compactness and tensor formula for the raw Thom spectrum]
\label{prop:infinitesimal-raw-thom-properties}
The object $\mathbb T_S^{\mathrm{raw}}$ is compact.  For every
$X\in\operatorname{Sm}^{\fp}_{/S}$ and $a\in\mathbb Z$, there is a canonical
cofiber identification
\[
\begin{aligned}
    \mathbb T_S^{\mathrm{raw}}
    \otimes
    \Sigma^a\Sigma^\infty_+h_X
    \simeq
    \Sigma^a\cofib\Bigl(
        \Sigma^\infty_+h_{X\times_S\mathbb G_{m,S}}
        \longrightarrow
        \Sigma^\infty_+h_{X\times_S\mathbb A^1_S}
    \Bigr).
\end{aligned}
\]
Consequently tensoring with $\mathbb T_S^{\mathrm{raw}}$ preserves compact
objects of $\mathcal N_S^{\Sp}$.
\end{proposition}

\begin{proof}
Both suspension representables in the defining cofiber are compact by
Proposition~\ref{prop:infinitesimal-stable-nisnevich}, and compact objects in
a stable presentable category are closed under finite colimits.  Thus
$\mathbb T_S^{\mathrm{raw}}$ is compact.

The tensor product preserves cofibers and suspension representables multiply
by fibre product over $S$.  Tensoring the defining cofiber with
$\Sigma^a\Sigma^\infty_+h_X$ gives the displayed formula.  Its two terms are
compact, so tensoring with $\mathbb T_S^{\mathrm{raw}}$ carries every
generator in $\mathcal G_S$ to a compact object.  Compact objects are the
idempotent-complete finite-colimit closure of the compact generators, and the
tensor functor is exact and preserves retracts.  Hence it preserves all
compact objects.
\end{proof}

\section{Raw symmetric Thom prespectra}
\label{sec:infinitesimal-raw-stable-coefficients}

\begin{definition}[Symmetric sequences]
\label{def:infinitesimal-symmetric-sequences}
Define
\[
    \SymSeq(\mathcal N_S^{\Sp})
    :=
    \Fun
    \left(
        \coprod_{m\geq0}B\Sigma_m,
        \mathcal N_S^{\Sp}
    \right)
\]
with additive Day convolution.  Let
\[
    \mathbb S_{\mathbb T_S^{\mathrm{raw}}}^{\mathrm{sym}}(m)
    :=
    (\mathbb T_S^{\mathrm{raw}})^{\otimes m}
\]
with its canonical $\Sigma_m$-action.  This is a commutative algebra object in
$\SymSeq(\mathcal N_S^{\Sp})$.
\end{definition}

\begin{definition}[Raw brave-new differential prespectra]
\label{def:infinitesimal-raw-prespectra}
Define
\[
    \mathcal P_S^{\mathrm{pre}}
    :=
    \DiffSp^{\mathrm{pre}}(S)
    :=
    \Mod_{\mathbb S_{\mathbb T_S^{\mathrm{raw}}}^{\mathrm{sym}}}
    \bigl(\SymSeq(\mathcal N_S^{\Sp})\bigr).
\]
An object of $\mathcal P_S^{\mathrm{pre}}$ is called a \emph{raw symmetric Thom
prespectrum}.  The category $\mathcal P_S^{\mathrm{pre}}$ is stable,
presentable, and symmetric monoidal under relative tensor product; limits and
colimits are created in underlying symmetric sequences.
\end{definition}

\begin{definition}[Free levels]
\label{def:infinitesimal-free-levels}
For $m\geq0$, let
\[
    F_{m,S}:\mathcal N_S^{\Sp}\longrightarrow\mathcal P_S^{\mathrm{pre}}
\]
be the left adjoint to level-$m$ evaluation followed by forgetting the
$\Sigma_m$-action.  Thus
\[
    \map_{\mathcal P_S^{\mathrm{pre}}}(F_{m,S}A,E)
    \simeq
    \map_{\mathcal N_S^{\Sp}}(A,E_m).
\]
\end{definition}

\begin{proposition}[Free-level tensor formula]
\label{prop:infinitesimal-free-level-tensor}
For all $m,n\geq0$ and $A,B\in\mathcal N_S^{\Sp}$, there is a canonical equivalence
\[
    \boxed{
    F_{m,S}(A)\otimes F_{n,S}(B)
    \simeq
    F_{m+n,S}(A\otimes B).
    }
\]
These equivalences are associative, symmetric, unital, and coherent in all
variables.  In particular, $F_{0,S}$ is colimit-preserving, fully faithful,
and strong symmetric monoidal.
\end{proposition}

\begin{proof}
Let $J_m(A)$ denote the symmetric sequence concentrated in level $m$ with the
free $\Sigma_m$-action on $A$.  Additive Day convolution gives
\[
    J_m(A)\otimes_{\mathrm{Day}}J_n(B)
    \simeq
    J_{m+n}(A\otimes B),
\]
because induction of the free $\Sigma_m\times\Sigma_n$-action along
$\Sigma_m\times\Sigma_n\subseteq\Sigma_{m+n}$ is the free
$\Sigma_{m+n}$-action.  The free-level object is
\[
    F_{m,S}(A)
    =
    \mathbb S_{\mathbb T_S^{\mathrm{raw}}}^{\mathrm{sym}}
    \otimes_{\mathrm{Day}}J_m(A).
\]
Relative tensor product of free modules therefore gives
\[
\begin{aligned}
    F_{m,S}(A)\otimes F_{n,S}(B)
    &\simeq
    \mathbb S_{\mathbb T_S^{\mathrm{raw}}}^{\mathrm{sym}}
    \otimes_{\mathrm{Day}}
    \bigl(J_m(A)\otimes_{\mathrm{Day}}J_n(B)\bigr)\\
    &\simeq
    F_{m+n,S}(A\otimes B).
\end{aligned}
\]
All coherence is inherited from Day convolution and the free-module
construction.  At level zero the free module has degree-zero term $A$, so the
unit $A\to\operatorname{ev}_0F_{0,S}A$ is an equivalence, proving full
faithfulness.
\end{proof}

\begin{proposition}[Compact generators of the raw prespectrum category]
\label{prop:infinitesimal-raw-compact-generators}
The objects
\[
    F_{m,S}G,
    \qquad
    m\geq0,
    \quad
    G\in\mathcal G_S,
\]
are compact generators of $\mathcal P_S^{\mathrm{pre}}$.
\end{proposition}

\begin{proof}
Evaluation at each prespectrum level creates filtered colimits.  Therefore the
adjunction $F_{m,S}\dashv\operatorname{ev}_m$ carries every compact
$G\in\mathcal G_S$ to a compact object.  If a raw prespectrum receives only
null mapping spectra from all $F_{m,S}G$, then every level receives only null
maps from all generators of $\mathcal N_S^{\Sp}$ and hence every level vanishes.
Thus the object itself vanishes.
\end{proof}

\begin{definition}[Raw Thom stabilization morphism]
\label{def:infinitesimal-raw-thom-stabilization}
For $m\geq0$ and $A\in\mathcal N_S^{\Sp}$, let
\[
    u_{m,A}:A\longrightarrow(F_{m,S}A)_m
\]
be the unit corresponding to $\id_{F_{m,S}A}$ under
$F_{m,S}\dashv\operatorname{ev}_m$.  The module structure of $F_{m,S}A$
gives a canonical composite
\[
\begin{aligned}
    \mathbb T_S^{\mathrm{raw}}\otimes A
    &\xrightarrow{\id\otimes u_{m,A}}
    \mathbb T_S^{\mathrm{raw}}\otimes(F_{m,S}A)_m\\
    &\longrightarrow
    (F_{m,S}A)_{m+1}.
\end{aligned}
\]
Define
\[
    \boxed{
    \lambda_{m,A}:
    F_{m+1,S}(\mathbb T_S^{\mathrm{raw}}\otimes A)
    \longrightarrow
    F_{m,S}(A)
    }
\]
to be its adjoint under
$F_{m+1,S}\dashv\operatorname{ev}_{m+1}$.  For
$G\in\mathcal G_S$ we write $\lambda_{m,G}$ for this morphism.
\end{definition}

\begin{lemma}[Tensor compatibility of raw Thom stabilization]
\label{lem:infinitesimal-raw-thom-stabilization-tensor}
For $m,n\geq0$ and $A,B\in\mathcal N_S^{\Sp}$, under the canonical free-level
equivalences
\[
\begin{aligned}
    F_{m+1,S}(\mathbb T_S^{\mathrm{raw}}\otimes A)
    \otimes F_{n,S}(B)
    &\simeq
    F_{m+n+1,S}(\mathbb T_S^{\mathrm{raw}}\otimes A\otimes B),\\
    F_{m,S}(A)\otimes F_{n,S}(B)
    &\simeq
    F_{m+n,S}(A\otimes B),
\end{aligned}
\]
the morphism
\[
    \lambda_{m,A}\otimes\id_{F_{n,S}(B)}
\]
identifies canonically with
\[
    \boxed{
    \lambda_{m+n,A\otimes B}.
    }
\]
These identifications are compatible with associativity, symmetry, and
iteration in the free levels.
\end{lemma}

\begin{proof}
By Definition~\ref{def:infinitesimal-raw-thom-stabilization},
$\lambda_{m,A}$ is adjoint to the composite
\[
    \mathbb T_S^{\mathrm{raw}}\otimes A
    \xrightarrow{\id\otimes u_{m,A}}
    \mathbb T_S^{\mathrm{raw}}\otimes(F_{m,S}A)_m
    \longrightarrow
    (F_{m,S}A)_{m+1},
\]
where the second arrow is the module action.  Tensor this construction with
$F_{n,S}(B)$.  Under the free-level tensor equivalence of
Proposition~\ref{prop:infinitesimal-free-level-tensor}, the level-$(m+n)$
unit for
\[
    F_{m+n,S}(A\otimes B)
\]
is the tensor product of the level-$m$ and level-$n$ free units, with the
canonical induction of symmetric-group actions.  Associativity of the
$\mathbb S_{\mathbb T_S^{\mathrm{raw}}}^{\mathrm{sym}}$-module action then
identifies the resulting adjoint map with
\[
    \mathbb T_S^{\mathrm{raw}}\otimes A\otimes B
    \longrightarrow
    \bigl(F_{m+n,S}(A\otimes B)\bigr)_{m+n+1}
\]
used to define $\lambda_{m+n,A\otimes B}$.  Applying
$F_{m+n+1,S}\dashv\operatorname{ev}_{m+n+1}$ gives the asserted
identification.  The coherence statements are inherited from Day convolution,
the free-module tensor formula, and the associative symmetric module action.
\end{proof}

\section{Thom stabilization before affine-line homotopification}
\label{sec:infinitesimal-thom-stable-differential-category}

The next localization is deliberately performed before affine-line
homotopification.  No symmetry statement for
\(\mathbb T_S^{\mathrm{raw}}\) is used in its construction.

\begin{definition}[Thom-stable differential spectra]
\label{def:infinitesimal-thom-stable-differential-category}
Let
\[
  L_{\mathrm{Th},S}:
  \mathcal P_S^{\mathrm{pre}}
  \longrightarrow
  \mathcal R_S^0
\]
be the accessible localization generated by the morphisms
\[
  \lambda_{m,G}:
  F_{m+1,S}(\mathbb T_S^{\mathrm{raw}}\otimes G)
  \longrightarrow
  F_{m,S}(G),
  \qquad
  m\geq0,
  \quad
  G\in\mathcal G_S.
\]
An object of \(\mathcal R_S^0\) is called a \emph{Thom-stable brave-new
differential coefficient}.  Write
\[
  Q_S:=L_{\mathrm{Th},S}F_{0,S}:
  \mathcal N_S^{\Sp}\longrightarrow\mathcal R_S^0
\]
for the level-zero realization and put
\[
  \mathbb T_S^0:=Q_S(\mathbb T_S^{\mathrm{raw}}).
\]
\end{definition}

\begin{theorem}[Raw Thom stabilization without affine-line localization]
\label{thm:infinitesimal-raw-thom-stabilization}
The localization \(L_{\mathrm{Th},S}\) is exact and strong symmetric monoidal.
Its local objects are precisely the symmetric raw Thom prespectra whose adjoint
structure maps
\[
  E_m
  \longrightarrow
  \operatorname{Hom}_{\mathcal N_S^{\Sp}}
  (\mathbb T_S^{\mathrm{raw}},E_{m+1})
\]
are equivalences for every \(m\geq0\).  The object
\(\mathbb T_S^0\) is tensor-invertible, with inverse canonically represented by
\[
  L_{\mathrm{Th},S}F_{1,S}(\mathbf1_S).
\]
The objects
\[
  L_{\mathrm{Th},S}F_{m,S}(G),
  \qquad m\geq0,
  \quad G\in\mathcal G_S,
\]
form a set of compact generators of \(\mathcal R_S^0\).
\end{theorem}

\begin{proof}
By Lemma~\ref{lem:infinitesimal-raw-thom-stabilization-tensor}, tensoring a
generating stabilization morphism by a free compact generator gives another
stabilization morphism:
\[
  \lambda_{m,G}\otimes F_{n,S}(H)
  \simeq
  \lambda_{m+n,G\otimes H}.
\]
The free compact generators generate \(\mathcal P_S^{\mathrm{pre}}\) under
colimits and tensor product preserves colimits separately.  Hence the strongly
saturated class generated by the \(\lambda_{m,G}\) is a tensor ideal.  The
symmetric monoidal localization theorem makes \(L_{\mathrm{Th},S}\) strong
symmetric monoidal.  The generating set is closed under suspension, so the
local subcategory is stable and the localization is exact.

For a prespectrum \(E\), the free-level adjunction identifies mapping a
generator into \(E\) with
\[
  \map_{\mathcal N_S^{\Sp}}(G,E_m)
  \longrightarrow
  \map_{\mathcal N_S^{\Sp}}
  (\mathbb T_S^{\mathrm{raw}}\otimes G,E_{m+1}).
\]
Tensor--Hom adjunction turns this into
\[
  \map_{\mathcal N_S^{\Sp}}(G,E_m)
  \longrightarrow
  \map_{\mathcal N_S^{\Sp}}
  \bigl(G,\operatorname{Hom}(\mathbb T_S^{\mathrm{raw}},E_{m+1})\bigr).
\]
Since \(\mathcal G_S\) detects equivalences, locality is equivalent to the
stated \(\Omega\)-spectrum condition.

At \(m=0\), the stabilization morphism and the free-level tensor formula give
\[
\begin{aligned}
  \mathbb T_S^0\otimes L_{\mathrm{Th},S}F_{1,S}(\mathbf1_S)
  &\simeq
  L_{\mathrm{Th},S}F_{1,S}(\mathbb T_S^{\mathrm{raw}})\\
  &\simeq
  L_{\mathrm{Th},S}F_{0,S}(\mathbf1_S)
   =\mathbf1_{\mathcal R_S^0}.
\end{aligned}
\]
Thus \(\mathbb T_S^0\) is invertible.  Notice that this argument never uses a
cyclic-permutation homotopy for \(\mathbb T_S^{\mathrm{raw}}\).

Finally, the sources and targets of the \(\lambda_{m,G}\) are compact by
Proposition~\ref{prop:infinitesimal-raw-compact-generators} and
Proposition~\ref{prop:infinitesimal-raw-thom-properties}.  Locality with respect
to maps between compact objects is preserved by filtered colimits, so the
inclusion \(\mathcal R_S^0\hookrightarrow\mathcal P_S^{\mathrm{pre}}\)
preserves filtered colimits.  The localization adjunction therefore shows that
each \(L_{\mathrm{Th},S}F_{m,S}(G)\) is compact.  If an object receives only
null maps from all of them, its image under the fully faithful inclusion
receives only null maps from every free compact generator and is zero.  Hence
they generate.
\end{proof}


\begin{theorem}[Comparison with formal raw Thom periodization]
\label{thm:infinitesimal-raw-thom-periodization}
There are canonical colimit-preserving strong symmetric monoidal functors
\[
    \boxed{
    \mathcal N_S^{\Sp}
    [ (\mathbb T_S^{\mathrm{raw}})^{-1} ]
    \xrightarrow{\ \Psi_S\ }
    \mathcal R_S^0
    \xrightarrow{\ \Phi_S\ }
    \mathcal N_S^{\Sp}
    [ (\mathbb T_S^{\mathrm{raw}})^{-1} ]
    }
\]
such that
\[
    \Psi_S\operatorname{per}_{\mathbb T_S^{\mathrm{raw}}}
    \simeq
    Q_S,
    \qquad
    \Phi_SQ_S
    \simeq
    \operatorname{per}_{\mathbb T_S^{\mathrm{raw}}},
\]
and
\[
    \boxed{
    \Phi_S\Psi_S\xrightarrow{\simeq}\id.
    }
\]
Thus formal symmetric monoidal inversion of the raw Thom coordinate is a
canonical symmetric monoidal retract of the Thom-stable differential category.
No equivalence in the opposite composite is asserted before affine-line
localization, and no symmetry hypothesis on
$\mathbb T_S^{\mathrm{raw}}$ is used in this comparison.
\end{theorem}

\begin{proof}
Put
\[
    \mathcal C:=\mathcal N_S^{\Sp},
    \qquad
    T:=\mathbb T_S^{\mathrm{raw}},
    \qquad
    \mathcal P:=\mathcal P_S^{\mathrm{pre}},
    \qquad
    L:=L_{\mathrm{Th},S},
\]
and write
\[
    Q:=LF_0:\mathcal C\longrightarrow\mathcal R_S^0.
\]
By Theorem~\ref{thm:infinitesimal-raw-thom-stabilization}, $Q(T)$ is
tensor-invertible.  The universal property of presentable symmetric monoidal
periodization from Chapter~6 therefore gives a canonical colimit-preserving
strong symmetric monoidal functor
\[
    \Psi:
    \mathcal C[T^{-1}]
    \longrightarrow
    \mathcal R_S^0
\]
with
\[
    \Psi\operatorname{per}_T\simeq Q.
\]

We construct the comparison in the other direction.  First note that
$\mathcal C[T^{-1}]$ is stable without any symmetry assumption on $T$.
Indeed, $\mathcal C[T^{-1}]$ is pointed by the formal periodization
construction.  In the stable category $\mathcal C$, the suspension
$\Sigma\mathbf1_{\mathcal C}$ of the tensor unit is tensor-invertible, with
inverse $\Sigma^{-1}\mathbf1_{\mathcal C}$.  Since
$\operatorname{per}_T$ preserves colimits and is strong symmetric monoidal,
\[
    \operatorname{per}_T(\Sigma\mathbf1_{\mathcal C})
    \simeq
    \Sigma\mathbf1_{\mathcal C[T^{-1}]}
\]
is tensor-invertible.  Tensor product preserves colimits separately, so
\[
    \Sigma\mathbf1_{\mathcal C[T^{-1}]}\otimes A
    \simeq
    \Sigma A
\]
for every $A$.  Suspension is therefore an equivalence, and the pointed
presentable category $\mathcal C[T^{-1}]$ is stable.

Write
\[
    \overline T:=\operatorname{per}_T(T)
    \in\Pic(\mathcal C[T^{-1}]).
\]
For a stable presentably symmetric monoidal category $\mathcal D$ and an
object $U\in\mathcal D$, write
\[
    \operatorname{PreSp}^{\Sigma}_U(\mathcal D)
    :=
    \Mod_{\mathbb S_U^{\mathrm{sym}}}
    (\SymSeq(\mathcal D)).
\]
Levelwise application of $\operatorname{per}_T$ gives a colimit-preserving
strong symmetric monoidal functor
\[
    \operatorname{per}_T^{\mathrm{pre}}:
    \operatorname{PreSp}^{\Sigma}_T(\mathcal C)
    \longrightarrow
    \operatorname{PreSp}^{\Sigma}_{\overline T}
    (\mathcal C[T^{-1}]).
\]
Choose a small tensor-closed generating family of the stable presentable
category $\mathcal C[T^{-1}]$, containing the tensor unit and closed under
integer suspension, and localize the target symmetric prespectra at the
corresponding stabilization morphisms.  The free-level tensor formula makes
this a strong symmetric monoidal localization, and the free-level adjunction
identifies its local objects with the symmetric $\overline T$-$\Omega$-spectra.
Since $\overline T$ is tensor-invertible, the free level-zero prespectrum
$F_0(A)$ is local for every $A$: its $m$-th level is
$\overline T^{\otimes m}\otimes A$ and all adjoint structure maps are
equivalences.  Moreover the adjunction
\[
    \overline T\otimes-
    \dashv
    \operatorname{Hom}(\overline T,-)
\]
is an adjoint equivalence.  Hence for every local prespectrum $E$ its adjoint
structure equivalences are equivalent to structure equivalences
\[
    \overline T\otimes E_m\xrightarrow{\simeq}E_{m+1}.
\]
Iteration gives
\[
    E_m\simeq\overline T^{\otimes m}\otimes E_0,
\]
and the map
\[
    F_0(E_0)\longrightarrow E
\]
adjoint to $\id_{E_0}$ is a levelwise equivalence.  Thus zeroth evaluation
identifies the local target category symmetric monoidally with
$\mathcal C[T^{-1}]$.  This description is independent of the chosen
generating family because in each case the local subcategory is exactly the
full subcategory defined by the displayed $\Omega$-spectrum condition.

Composing levelwise periodization, this target stabilization, and zeroth
evaluation gives a colimit-preserving strong symmetric monoidal functor
\[
    \Phi_0:\mathcal P\longrightarrow\mathcal C[T^{-1}].
\]
For every $m$ and $A$, strong monoidality of $\operatorname{per}_T$ carries
the defining free-level unit and module action for $\lambda_{m,A}$ to those
for the corresponding $\overline T$-stabilization morphism.  Consequently
$\Phi_0$ inverts every $\lambda_{m,G}$ and factors uniquely through a
colimit-preserving strong symmetric monoidal functor
\[
    \Phi:\mathcal R_S^0\longrightarrow\mathcal C[T^{-1}].
\]
On level-zero free objects the construction gives
\[
    \Phi Q(A)\simeq\operatorname{per}_T(A).
\]
Therefore
\[
    \Phi\Psi\operatorname{per}_T
    \simeq
    \operatorname{per}_T.
\]
By the universal property of formal periodization, restriction along
$\operatorname{per}_T$ is fully faithful on colimit-preserving strong
symmetric monoidal functors out of $\mathcal C[T^{-1}]$.  Hence the preceding
identification determines a canonical equivalence
\[
    \Phi\Psi\xrightarrow{\simeq}\id_{\mathcal C[T^{-1}]}.
\]
This proves all assertions.
\end{proof}

\begin{remark}[Why symmetry is postponed]
\label{rem:infinitesimal-raw-thom-symmetry-postponed}
Chapter~6 proves symmetry of the motivic Thom coordinate by an
\(\mathbb A^1\)-homotopy between the identity and a cyclic permutation.  That
argument is not available before affine-line localization.  The raw symmetric
prespectrum construction and
Theorem~\ref{thm:infinitesimal-raw-thom-stabilization} therefore use no
symmetry assertion: Thom stabilization is defined directly and makes
\(\mathbb T_S^0\) tensor-invertible.  The formal periodization comparison of
Theorem~\ref{thm:infinitesimal-raw-thom-periodization} supplies only the
canonical split comparison before affine-line localization; it does not
identify the two categories.  After affine-line localization the coordinate
becomes the symmetric motivic Thom sphere, and the symmetric-prespectrum
periodization theorem below then identifies the simultaneous
Thom-and-affine-line localization with the periodized model of \(\SH(S)\).
\end{remark}

\begin{definition}[Stable affine-line localization]
\label{def:infinitesimal-stable-a1-localization}
For $X\in\operatorname{Sm}^{\fp}_{/S}$ and $a\in\mathbb Z$, let
\[
    \beta_{X,a}:
    \Sigma^a\Sigma^\infty_+h_{X\times_S\mathbb A^1_S}
    \longrightarrow
    \Sigma^a\Sigma^\infty_+h_X
\]
be induced by affine-line projection.  Let
\[
    L^{\mathrm{st}}_{\mathbb A^1,S}:
    \mathcal N_S^{\Sp}
    \longrightarrow
    \mathcal N_S^{\Sp,\mathbb A^1}
\]
be the accessible localization generated by the set of all $\beta_{X,a}$,
and put
\[
    \overline{\mathbb T}_S
    :=
    L^{\mathrm{st}}_{\mathbb A^1,S}\mathbb T_S^{\mathrm{raw}}.
\]
\end{definition}

\begin{lemma}[Spectrum objects and spectral Nisnevich sheaves]
\label{lem:infinitesimal-spectrum-sheaf-interchange}
Put
\[
    \mathcal N_S
    :=
    \operatorname{Shv}_{\Nis}(\mathcal D_S;\mathcal S),
    \qquad
    \mathcal N_{S,\bullet}:=(\mathcal N_S)_*.
\]
There is a canonical symmetric monoidal equivalence
\[
    \boxed{
    \operatorname{Sp}_{S^1}(\mathcal N_{S,\bullet})
    \xrightarrow{\simeq}
    \mathcal N_S^{\Sp}.
    }
\]
Under this equivalence, the full subcategory of spectrum-valued sheaves which
are stable affine-line local identifies canonically with
\[
    \operatorname{Sp}_{S^1}\bigl(\mathcal H_\bullet(S)\bigr).
\]
\end{lemma}

\begin{proof}
Objectwise stabilization gives a canonical equivalence
\[
    \operatorname{Sp}_{S^1}
    \left(
        \operatorname{Fun}(\mathcal D_S^{\op},\mathcal S)_*
    \right)
    \simeq
    \operatorname{Fun}(\mathcal D_S^{\op},\Sp).
\]
A spectrum-valued presheaf $E$ satisfies spectral Nisnevich descent if and
only if, for every $n\in\mathbb Z$, the space-valued presheaf
\[
    \Omega^\infty\Sigma^nE
\]
satisfies spectral Nisnevich descent.  Indeed, the Nisnevich descent diagrams
are limits, and the family of functors
$\{\Omega^\infty\Sigma^n\}_{n\in\mathbb Z}$ detects equivalences of spectra.
Consequently the preceding objectwise stabilization equivalence restricts to
\[
    \operatorname{Sp}_{S^1}(\mathcal N_{S,\bullet})
    \xrightarrow{\simeq}
    \operatorname{Shv}_{\Nis}(\mathcal D_S;\Sp)
    =
    \mathcal N_S^{\Sp}.
\]

The equivalence is symmetric monoidal.  On suspension-spectrum
representables it sends the tensor product to fibre product over $S$:
\[
    \Sigma^\infty_+h_X\otimes\Sigma^\infty_+h_Y
    \simeq
    \Sigma^\infty_+h_{X\times_SY}.
\]
These objects generate under colimits, and both tensor products preserve
colimits separately, so the comparison extends uniquely as a symmetric
monoidal equivalence.

Finally, the pointed motivic localization
\[
    \mathcal N_{S,\bullet}
    \longrightarrow
    \mathcal H_\bullet(S)
\]
has fully faithful right adjoint the inclusion
\[
    \mathcal H_\bullet(S)
    \hookrightarrow
    \mathcal N_{S,\bullet}.
\]
The inclusion creates limits.  Consequently an $S^1$-spectrum object of
$\mathcal N_{S,\bullet}$ belongs to
$\operatorname{Sp}_{S^1}(\mathcal H_\bullet(S))$ if and only if every one of
its pointed space-valued levels is affine-line local.  Under the objectwise
stabilization equivalence, this levelwise condition is equivalent to locality
for every stable affine-line projection map.  Hence the stable
affine-line-local subcategory is
\[
    \operatorname{Sp}_{S^1}\bigl(\mathcal H_\bullet(S)\bigr).
\]
\end{proof}

\begin{proposition}[Stable affine-line comparison]
\label{prop:infinitesimal-stable-a1-comparison}
The localization $L^{\mathrm{st}}_{\mathbb A^1,S}$ is exact and strong
symmetric monoidal.  The images of $\mathcal G_S$ are compact generators of
$\mathcal N_S^{\Sp,\mathbb A^1}$.  With $\mathcal N_S$ and
$\mathcal N_{S,\bullet}$ as in
Lemma~\ref{lem:infinitesimal-spectrum-sheaf-interchange}, there is a canonical
symmetric monoidal equivalence
\[
    \boxed{
    \mathcal N_S^{\Sp,\mathbb A^1}
    \simeq
    \operatorname{Sp}_{S^1}\bigl(\mathcal H_\bullet(S)\bigr),
    }
\]
under which
\[
    \overline{\mathbb T}_S
    \simeq
    \Sigma^\infty_{S^1}T_S,
\]
where $T_S=\operatorname{Th}_S(\mathcal O_S)$ is the motivic Thom sphere of
Chapter~6.  In particular, $\overline{\mathbb T}_S$ is a symmetric tensor
object.
\end{proposition}

\begin{proof}
The set $\{\beta_{X,a}\}$ is stable under suspension, so its localization in
the stable category is exact.  Tensoring $\beta_{X,a}$ with a generator
$\Sigma^b\Sigma^\infty_+h_Y$ gives $\beta_{X\times_SY,a+b}$.  Hence the
stable affine-line equivalences form a tensor ideal and the localization is
strong symmetric monoidal.

Affine-line locality is preserved by filtered colimits because it is an
objectwise equivalence condition in spectra.  Therefore the inclusion of the
local subcategory preserves filtered colimits, and the localization of every
compact generator remains compact.  The localized generators detect zero
objects by the localization adjunction, so they generate.

Lemma~\ref{lem:infinitesimal-spectrum-sheaf-interchange} identifies the
stable affine-line-local objects with
\[
    \operatorname{Sp}_{S^1}\bigl(\mathcal H_\bullet(S)\bigr)
\]
through a canonical symmetric monoidal equivalence.

Exact affine-line localization carries the raw Thom cofiber to
\[
    \cofib\left(
      \Sigma^\infty_{S^1}\operatorname{Mot}_{S,+}(\mathbb G_{m,S})
      \longrightarrow
      \Sigma^\infty_{S^1}\operatorname{Mot}_{S,+}(\mathbb A^1_S)
    \right).
\]
The pointed cofiber inside $\mathcal H_\bullet(S)$ is exactly
\[
    T_S=\mathbb A^1_S/(\mathbb A^1_S\setminus0),
\]
so the displayed object is $\Sigma^\infty_{S^1}T_S$.  Chapter~6 proves that
the cyclic permutation on $T_S^{\otimes3}$ is $\mathbb A^1$-homotopic to the
identity.  Thus $T_S$ is symmetric, and strong monoidality of
$\Sigma^\infty_{S^1}$ makes $\overline{\mathbb T}_S$ symmetric as well.
\end{proof}

\begin{lemma}[Levelwise monoidal localization of symmetric prespectra]
\label{lem:infinitesimal-levelwise-module-localization}
Let
\[
    L:\mathcal C\rightleftarrows\mathcal C':j
\]
be an accessible strong symmetric monoidal localization of stable presentably
symmetric monoidal $\infty$-categories, and let $T\in\mathcal C$.  Then
levelwise application of $L$ induces an accessible strong symmetric monoidal
localization
\[
    \overline L:
    \Mod_{\mathbb S_T^{\mathrm{sym}}}
    (\SymSeq(\mathcal C))
    \longrightarrow
    \Mod_{\mathbb S_{LT}^{\mathrm{sym}}}
    (\SymSeq(\mathcal C')).
\]
Its fully faithful right adjoint is obtained by applying $j$ levelwise with
the module structure induced by the canonical lax symmetric monoidal
structure on $j$.  A symmetric $T$-prespectrum is $\overline L$-local if and
only if every underlying level is $L$-local.

Writing
\[
    F_m^T:\mathcal C
    \longrightarrow
    \Mod_{\mathbb S_T^{\mathrm{sym}}}(\SymSeq(\mathcal C))
\]
and
\[
    F_m^{LT}:\mathcal C'
    \longrightarrow
    \Mod_{\mathbb S_{LT}^{\mathrm{sym}}}(\SymSeq(\mathcal C'))
\]
for the corresponding free level-$m$ functors, there are canonical
equivalences
\[
    \boxed{
    \overline L\,F_m^T(A)
    \xrightarrow{\simeq}
    F_m^{LT}(LA).
    }
\]
If
\[
    \lambda^T_{m,A}:
    F_{m+1}^T(T\otimes A)
    \longrightarrow
    F_m^T(A)
\]
and
\[
    \lambda^{LT}_{m,LA}:
    F_{m+1}^{LT}(LT\otimes LA)
    \longrightarrow
    F_m^{LT}(LA)
\]
denote the stabilization morphisms adjoint to the canonical module-action
composites, then under the preceding free-level equivalences and the strong
monoidal equivalence
\[
    L(T\otimes A)\simeq LT\otimes LA
\]
one has
\[
    \boxed{
    \overline L(\lambda^T_{m,A})
    \simeq
    \lambda^{LT}_{m,LA}.
    }
\]
\end{lemma}

\begin{proof}
Levelwise application of $L$ gives a colimit-preserving strong symmetric
monoidal functor
\[
    L^{\Sigma}:\SymSeq(\mathcal C)\longrightarrow\SymSeq(\mathcal C')
\]
for additive Day convolution.  Strong monoidality gives
\[
    L^{\Sigma}(\mathbb S_T^{\mathrm{sym}})
    \simeq
    \mathbb S_{LT}^{\mathrm{sym}}.
\]
Hence extension of scalars defines $\overline L$ on the module categories.
The right adjoint is levelwise $j$; the lax symmetric monoidal structure of
$j$ supplies the action of $\mathbb S_T^{\mathrm{sym}}$ on the underlying
symmetric sequence.

The module adjunction is detected on underlying symmetric sequences, so full
faithfulness of $j$ implies full faithfulness of the module right adjoint.
The unit
\[
    E\longrightarrow\overline j\,\overline L E
\]
is an equivalence precisely when its underlying level maps
\[
    E_m\longrightarrow jLE_m
\]
are equivalences for every $m$.  This proves the characterization of local
objects.

Let $J_m(A)$ denote the symmetric sequence concentrated in level $m$ with the
free $\Sigma_m$-action on $A$.  The free level object is
\[
    F_m^T(A)
    \simeq
    \mathbb S_T^{\mathrm{sym}}
    \otimes_{\mathrm{Day}}
    J_m(A).
\]
Applying extension of scalars and using strong monoidality gives
\[
\begin{aligned}
    \overline L\,F_m^T(A)
    &\simeq
    \mathbb S_{LT}^{\mathrm{sym}}
    \otimes_{\mathrm{Day}}
    L^\Sigma J_m(A)\\
    &\simeq
    \mathbb S_{LT}^{\mathrm{sym}}
    \otimes_{\mathrm{Day}}
    J_m(LA)\\
    &\simeq
    F_m^{LT}(LA).
\end{aligned}
\]
These equivalences are natural in $m$ and $A$.

The morphism $\lambda^T_{m,A}$ is adjoint to the composite
\[
    T\otimes A
    \longrightarrow
    T\otimes(F_m^T A)_m
    \longrightarrow
    (F_m^T A)_{m+1}
\]
formed from the free-level unit and the
$\mathbb S_T^{\mathrm{sym}}$-module action.  The strong symmetric monoidal
functor $L^\Sigma$, together with extension of scalars, carries the free-level
unit to the free-level unit, the module action to the
$\mathbb S_{LT}^{\mathrm{sym}}$-module action, and identifies
$L(T\otimes A)$ with $LT\otimes LA$.  Hence it carries the defining adjoint
composite for $\lambda^T_{m,A}$ to the defining adjoint composite for
$\lambda^{LT}_{m,LA}$.  This proves
\[
    \overline L(\lambda^T_{m,A})
    \simeq
    \lambda^{LT}_{m,LA}.
\]

Finally, relative tensor products of modules are geometric realizations of
two-sided bar constructions.  Since $L^{\Sigma}$ is strong symmetric
monoidal and preserves colimits, it carries the bar construction computing
$E\otimes_{\mathbb S_T}F$ to the bar construction computing
$\overline LE\otimes_{\mathbb S_{LT}}\overline LF$.  Hence $\overline L$ is
strong symmetric monoidal.
\end{proof}

\begin{lemma}[Symmetric spectra for an invertible coordinate]
\label{lem:infinitesimal-invertible-coordinate-spectra}
Let $\mathcal C$ be a stable presentably symmetric monoidal category and let
$T\in\mathcal C$ be tensor-invertible.  Let
\[
    \Omega\operatorname{Sp}^{\Sigma}_T(\mathcal C)
\]
be the full subcategory of symmetric $T$-prespectra $E$ for which every
adjoint structure morphism
\[
    E_m\longrightarrow\operatorname{Hom}(T,E_{m+1})
\]
is an equivalence.  This full subcategory is closed under the symmetric
monoidal structure on symmetric $T$-prespectra, and zeroth evaluation induces
a symmetric monoidal equivalence
\[
    \operatorname{ev}_0:
    \Omega\operatorname{Sp}^{\Sigma}_T(\mathcal C)
    \xrightarrow{\simeq}
    \mathcal C.
\]
Its inverse sends $A$ to the symmetric $T$-spectrum with $m$-th level
$T^{\otimes m}\otimes A$ and the canonical symmetric-group actions.
\end{lemma}

\begin{proof}
Let $A\in\mathcal C$.  The free level-zero module
\[
    F_0(A)=\mathbb S_T^{\mathrm{sym}}\otimes J_0(A)
\]
has $m$-th level $T^{\otimes m}\otimes A$.  Since $T$ is invertible, the
adjunction
\[
    T\otimes-\dashv\operatorname{Hom}(T,-)
\]
is an adjoint equivalence.  Hence every adjoint structure map of $F_0(A)$ is
an equivalence, so $F_0(A)$ is a symmetric $T$-$\Omega$-spectrum and
$\operatorname{ev}_0F_0(A)\simeq A$.

Conversely, let $E$ be a symmetric $T$-$\Omega$-spectrum.  Since tensoring
with $T$ is an equivalence, the adjoint structure equivalences are equivalent
to structure equivalences
\[
    T\otimes E_m\xrightarrow{\simeq}E_{m+1}.
\]
By the adjunction $F_0\dashv\operatorname{ev}_0$, the identity of $E_0$
determines a canonical morphism of symmetric $T$-prespectra
\[
    F_0(E_0)\longrightarrow E.
\]
Its level-$m$ component is the iterated action morphism
\[
    T^{\otimes m}\otimes E_0\longrightarrow E_m.
\]
Each of these is an equivalence by iteration of the preceding structure
equivalences.  Equivalences of modules over
$\mathbb S_T^{\mathrm{sym}}$ are detected on their underlying symmetric
sequences, hence levelwise.  Therefore
\[
    F_0(E_0)\xrightarrow{\simeq}E.
\]
Thus $F_0$ and zeroth evaluation are mutually inverse equivalences.

The free-module Day-convolution calculation gives
\[
    F_0(A)\otimes F_0(B)
    \simeq
    F_0(A\otimes B),
    \qquad
    F_0(\mathbf 1)\simeq\mathbf 1.
\]
Hence $F_0$ is strong symmetric monoidal.  Since every
$T$-$\Omega$-spectrum is canonically equivalent to an object of the form
$F_0(A)$, the full subcategory of $T$-$\Omega$-spectra is closed under tensor
products and contains the tensor unit.  The mutually inverse equivalences
$F_0$ and $\operatorname{ev}_0$ are therefore symmetric monoidal.
\end{proof}

\begin{lemma}[Stability of periodization of a stable category]
\label{lem:infinitesimal-stable-periodization}
Let $\mathcal C$ be a stable presentably symmetric monoidal category and let
$T\in\mathcal C$ be arbitrary.  Then the presentable periodization
\[
    \mathcal C[T^{-1}]
\]
is stable.  No symmetry hypothesis on $T$ is required.
\end{lemma}

\begin{proof}
Put
\[
    \mathcal D:=\mathcal C[T^{-1}]
\]
and let
\[
    \operatorname{per}_T:\mathcal C\longrightarrow\mathcal D
\]
be the formal periodization functor of Chapter~6.  The periodized category is
pointed by construction, and $\operatorname{per}_T$ preserves colimits and is
strong symmetric monoidal.  In the stable category $\mathcal C$, the
suspension of the tensor unit is tensor-invertible, with inverse the
desuspension of the tensor unit.  Colimit preservation gives a canonical equivalence
\[
    \operatorname{per}_T(\Sigma\mathbf1_{\mathcal C})
    \simeq
    \Sigma\mathbf1_{\mathcal D}.
\]
Strong monoidality then shows that this object is tensor-invertible: an inverse
is
\[
    \operatorname{per}_T(\Sigma^{-1}\mathbf1_{\mathcal C}),
\]
since the tensor product of the two images is the image of the tensor unit.

Because the tensor product of $\mathcal D$ preserves colimits separately, for
every $A\in\mathcal D$ there is a canonical equivalence
\[
    \Sigma\mathbf1_{\mathcal D}\otimes A
    \simeq
    \Sigma A.
\]
Thus the suspension functor of the pointed presentable category $\mathcal D$
is tensoring with an invertible object and is therefore an equivalence.
A pointed presentable $\infty$-category whose suspension functor is an
equivalence is stable.  Hence $\mathcal C[T^{-1}]$ is stable.
\end{proof}

\begin{lemma}[Symmetric-prespectrum periodization]
\label{lem:infinitesimal-symmetric-periodization}
Let $\mathcal C$ be a stable presentably symmetric monoidal category, let
$T\in\mathcal C$ be symmetric, and let $\mathcal G$ be a tensor-closed set of
compact generators containing the tensor unit and closed under suspension.
Put
\[
    \operatorname{PreSp}^{\Sigma}_T(\mathcal C)
    :=
    \Mod_{\mathbb S_T^{\mathrm{sym}}}
    \bigl(\SymSeq(\mathcal C)\bigr).
\]
Let $L_{\mathrm{st}}$ be the accessible localization generated by
\[
    F_{m+1}(T\otimes G)\longrightarrow F_m(G),
    \qquad
    m\geq0,
    \quad
    G\in\mathcal G.
\]
Then there is a canonical symmetric monoidal equivalence
\[
    \boxed{
    \operatorname{PreSp}^{\Sigma}_T(\mathcal C)_{\mathrm{st}}
    \simeq
    \mathcal C[T^{-1}].
    }
\]
Under this equivalence, $L_{\mathrm{st}}F_0$ identifies with the
periodization functor.
\end{lemma}

\begin{proof}
Write
\[
    \mathcal P:=\operatorname{PreSp}^{\Sigma}_T(\mathcal C),
    \qquad
    \mathcal P_{\mathrm{st}}:=L_{\mathrm{st}}\mathcal P.
\]
The free-level tensor formula gives
\[
    F_m(A)\otimes F_n(B)
    \simeq
    F_{m+n}(A\otimes B).
\]
Tensoring a stabilization generator by $F_n(H)$, with $H\in\mathcal G$,
gives
\[
    F_{m+n+1}(T\otimes G\otimes H)
    \longrightarrow
    F_{m+n}(G\otimes H).
\]
Since $\mathcal G$ is tensor-closed, the generating equivalences are stable
under tensoring with the free compact generators.  Those objects generate
$\mathcal P$ under colimits, and tensor product preserves colimits separately.
Thus the strongly saturated stabilization class is a tensor ideal and
$L_{\mathrm{st}}$ is a strong symmetric monoidal localization.

For a symmetric $T$-prespectrum $E$, the free-level adjunction identifies
locality for the generator indexed by $(m,G)$ with invertibility of
\[
    \map_{\mathcal C}(G,E_m)
    \longrightarrow
    \map_{\mathcal C}(T\otimes G,E_{m+1})
    \simeq
    \map_{\mathcal C}
    \bigl(G,\operatorname{Hom}(T,E_{m+1})\bigr).
\]
Since $\mathcal G$ detects equivalences, the local objects are exactly the
symmetric $T$-$\Omega$-spectra, namely those for which every adjoint structure
map
\[
    E_m\xrightarrow{\simeq}\operatorname{Hom}(T,E_{m+1})
\]
is an equivalence.

It remains to check that localizing only on $\mathcal G$ gives the full
symmetric-spectrum stabilization.  For every $m\geq0$, let
$\mathcal K_m\subseteq\mathcal C$ be the full subcategory of objects $A$ for
which
\[
    L_{\mathrm{st}}
    \bigl(F_{m+1}(T\otimes A)\to F_m(A)\bigr)
\]
is an equivalence.  The source and target depend exactly and
colimit-preservingly on $A$, and $L_{\mathrm{st}}$ is exact and preserves
colimits.  Hence $\mathcal K_m$ is a localizing subcategory.  It contains
$\mathcal G$ by construction, so $\mathcal K_m=\mathcal C$.  Therefore
$L_{\mathrm{st}}$ is equivalently the localization at the stabilization
morphisms for all objects $A\in\mathcal C$.

For a symmetric coordinate $T$, the symmetric-spectrum comparison theorem
identifies this localization of symmetric $T$-prespectra with the formal
symmetric monoidal inversion of $T$; this is the comparison used in Chapter~6,
cf.~\cite{RobaloNoncommutativeMotivesI}.  Consequently there is a canonical
symmetric monoidal equivalence
\[
    \mathcal P_{\mathrm{st}}
    \xrightarrow{\simeq}
    \mathcal C[T^{-1}].
\]
The comparison sends the free level-zero stabilization functor to the
canonical functor from $\mathcal C$ which makes $T$ invertible.  By the
universal property of Chapter~6 this canonical functor is
$\operatorname{per}_T$.  Hence
\[
    L_{\mathrm{st}}F_0
    \simeq
    \operatorname{per}_T
\]
under the displayed equivalence.  Contractible uniqueness in the formal
periodization universal property supplies the stated symmetric monoidal and
higher coherences.
\end{proof}

\section{Motivic homotopification after Thom stabilization}
\label{sec:infinitesimal-bnv-adjoint-triple}

For \(m\geq0\), write
\[
  F^0_{m,S}:=L_{\mathrm{Th},S}F_{m,S}:
  \mathcal N_S^{\Sp}\longrightarrow\mathcal R_S^0.
\]
For \(X\in\operatorname{Sm}^{\fp}_{/S}\) and \(a\in\mathbb Z\), define
\[
  \alpha^0_{m,X,a}:
  F^0_{m,S}
  \bigl(\Sigma^a\Sigma^\infty_+h_{X\times_S\mathbb A^1_S}\bigr)
  \longrightarrow
  F^0_{m,S}
  \bigl(\Sigma^a\Sigma^\infty_+h_X\bigr)
\]
to be the image of affine-line projection.

\begin{definition}[Motivic-local differential coefficient]
\label{def:infinitesimal-motivic-local-raw}
An object \(E\in\mathcal R_S^0\) is \emph{motivic local} if it is local with
respect to every \(\alpha^0_{m,X,a}\).  Write
\[
  (\mathcal R_S^0)_{\mathbb A^1}
  \subseteq
  \mathcal R_S^0
\]
for the full local subcategory.
\end{definition}

\begin{lemma}[Successive-localization comparison]
\label{lem:infinitesimal-successive-localizations}
Let \(\mathcal C\) be a presentable \(\infty\)-category and let \(A\) and \(B\)
be sets of morphisms.  If \(L_A:\mathcal C\to\mathcal C[A^{-1}]\) is the
localization at \(A\), then localization of \(\mathcal C[A^{-1}]\) at the images
of \(B\) has the same local objects as localization of \(\mathcal C\) at
\(A\cup B\).  Hence there is a canonical equivalence
\[
  \mathcal C[A^{-1}][(L_AB)^{-1}]
  \simeq
  \mathcal C[(A\cup B)^{-1}].
\]
The equivalence is compatible with symmetric monoidal structures whenever the
three localizations are symmetric monoidal.
\end{lemma}

\begin{proof}
An object of \(\mathcal C[A^{-1}]\), viewed through the fully faithful right
adjoint, is already \(A\)-local.  It is local for every image \(L_A(b)\) if and
only if it is also right orthogonal to every \(b\in B\).  Thus both sides have
the same fully faithful local subcategory of \(\mathcal C\).  The assertion
follows from uniqueness of reflective localization.  The monoidal statement
follows from the same uniqueness in the symmetric monoidal localization
category.
\end{proof}

\begin{theorem}[Symmetric-prespectrum comparison in the corrected order]
\label{thm:infinitesimal-symmetric-prespectrum-comparison}
There is a canonical symmetric monoidal equivalence
\[
  \boxed{
  (\mathcal R_S^0)_{\mathbb A^1}
  \simeq
  \SH(S).
  }
\]
Equivalently, the composite
\[
  \mathcal P_S^{\mathrm{pre}}
  \xrightarrow{L_{\mathrm{Th},S}}
  \mathcal R_S^0
  \longrightarrow
  (\mathcal R_S^0)_{\mathbb A^1}
\]
is the localization of raw symmetric Thom prespectra obtained by imposing both
raw Thom stabilization and stable affine-line invariance, but the BNV
homotopification on \(\mathcal R_S^0\) itself imposes only the latter.
\end{theorem}

\begin{proof}
By Lemma~\ref{lem:infinitesimal-successive-localizations}, the left-hand side is
the localization of \(\mathcal P_S^{\mathrm{pre}}\) at the union of the
stabilization morphisms \(\lambda_{m,G}\) and the affine-line projection
morphisms at every free level.

Compute this common localization in the opposite order.  First impose the
affine-line maps.  By the free-level adjunction, locality is equivalent to
stable affine-line locality of every prespectrum level.  Applying
Lemma~\ref{lem:infinitesimal-levelwise-module-localization} to
\(L^{\mathrm{st}}_{\mathbb A^1,S}\) identifies the intermediate category with
\[
  \Mod_{\mathbb S_{\overline{\mathbb T}_S}^{\mathrm{sym}}}
  \bigl(\SymSeq(\mathcal N_S^{\Sp,\mathbb A^1})\bigr).
\]
The image of \(\lambda_{m,G}\) is the corresponding stabilization morphism
\[
  F^{\mathbb A^1}_{m+1,S}
  \bigl(\overline{\mathbb T}_S\otimes
  L^{\mathrm{st}}_{\mathbb A^1,S}G\bigr)
  \longrightarrow
  F^{\mathbb A^1}_{m,S}
  \bigl(L^{\mathrm{st}}_{\mathbb A^1,S}G\bigr).
\]
Proposition~\ref{prop:infinitesimal-stable-a1-comparison} shows that
\(\overline{\mathbb T}_S\) is symmetric, so
Lemma~\ref{lem:infinitesimal-symmetric-periodization} identifies the remaining
localization with
\[
  \mathcal N_S^{\Sp,\mathbb A^1}
  [\overline{\mathbb T}_S^{-1}].
\]
Under
\[
  \mathcal N_S^{\Sp,\mathbb A^1}
  \simeq
  \operatorname{Sp}_{S^1}(\mathcal H_\bullet(S)),
\]
the object \(\overline{\mathbb T}_S\) is
\(\Sigma^\infty_{S^1}T_S\).  We now compare this periodization with the
one used in Chapter~6.  There is a canonical motivic factorization
\[
  T_S\simeq S^1_S\wedge G_{m,S}.
\]
If a tensor product \(A\otimes B\) is tensor-invertible in a symmetric
monoidal category, then each factor is tensor-invertible: an inverse for
\(A\) is \(B\otimes(A\otimes B)^{-1}\), and similarly for \(B\).
Consequently every colimit-preserving strong symmetric monoidal functor out of
\(\mathcal H_\bullet(S)\) which inverts \(T_S\) already inverts \(S^1_S\).
By the universal property of \(S^1\)-stabilization it therefore factors,
through a contractible space of choices, through
\[
  \mathcal H_\bullet(S)
  \longrightarrow
  \operatorname{Sp}_{S^1}(\mathcal H_\bullet(S)).
\]
The induced functor sends \(\Sigma^\infty_{S^1}T_S\) to an invertible object,
so it factors uniquely through
\[
  \operatorname{Sp}_{S^1}(\mathcal H_\bullet(S))
  [ (\Sigma^\infty_{S^1}T_S)^{-1}].
\]
Conversely, restriction along the stabilization functor sends every functor
out of this last periodization to a functor on \(\mathcal H_\bullet(S)\)
which inverts \(T_S\).  Thus the two presentable symmetric monoidal
categories have the same universal property, and there is a canonical
symmetric monoidal equivalence
\[
  \operatorname{Sp}_{S^1}(\mathcal H_\bullet(S))
  [ (\Sigma^\infty_{S^1}T_S)^{-1}]
  \xrightarrow{\simeq}
  \mathcal H_\bullet(S)[T_S^{-1}]
  =\SH(S).
\]
All localizations used above are tensor-ideal localizations, so the comparison
is symmetric monoidal.
\end{proof}

Under Theorem~\ref{thm:infinitesimal-symmetric-prespectrum-comparison}, write
\[
  \Disc_S:\SH(S)\hookrightarrow\mathcal R_S^0
\]
for the inclusion of motivic-local objects and
\[
  \mathbb H_S:\mathcal R_S^0\longrightarrow\SH(S)
\]
for the reflective affine-line homotopification.

\begin{lemma}[Bireflectivity of compactly generated stable localizations]
\label{lem:infinitesimal-compact-stable-bireflection}
Let \(\mathcal C\) be a stable presentable \(\infty\)-category and let
\(\mathcal W_0\) be a set of morphisms between compact objects, closed under
integer suspension.  Let
\[
  L:\mathcal C\rightleftarrows\mathcal D:j
\]
be the accessible localization generated by \(\mathcal W_0\).  Then the local
subcategory \(\mathcal D\) is closed under all small colimits computed in
\(\mathcal C\).  Consequently \(j\) preserves all small limits and colimits
and admits a right adjoint
\[
  j\dashv S.
\]
All three functors in \(L\dashv j\dashv S\) are exact.
\end{lemma}

\begin{proof}
For \(w\in\mathcal W_0\), put \(K_w:=\cofib(w)\).  The object \(K_w\) is
compact.  Since the generating set is closed under suspension, an object is
local precisely when
\[
  \map_{\mathcal C}(K_w,E)\simeq0
\]
for every \(w\).  This condition is stable under limits, fibres, cofibres and
arbitrary coproducts; compactness gives
\[
  \map_{\mathcal C}
  \left(K_w,\bigoplus_iE_i\right)
  \simeq
  \bigoplus_i\map_{\mathcal C}(K_w,E_i).
\]
Arbitrary colimits in a stable presentable category are generated by coproducts
and cofibres, so the local subcategory is closed under all colimits.  The
accessible inclusion therefore admits a right adjoint by the presentable
adjoint functor theorem.  Exactness follows from stability.
\end{proof}

\begin{theorem}[Stable BNV adjoint triple]
\label{thm:infinitesimal-bnv-adjoint-triple}
The inclusion \(\Disc_S\) occurs in an exact adjoint triple
\[
  \boxed{
  \mathbb H_S
  \dashv
  \Disc_S
  \dashv
  \mathbb S_S.
  }
\]
The functor \(\Disc_S\) is fully faithful and preserves all small limits and
colimits.  The left adjoint \(\mathbb H_S\) is stable affine-line
homotopification of an already Thom-stable coefficient, and the right adjoint
\(\mathbb S_S\) is the stable BNV secondary coreflection.
\end{theorem}

\begin{proof}
The sources and targets of every \(\alpha^0_{m,X,a}\) are compact by
Theorem~\ref{thm:infinitesimal-raw-thom-stabilization}.  The generating set is
closed under integer suspension.  Apply
Lemma~\ref{lem:infinitesimal-compact-stable-bireflection} to the
\(\mathbb A^1\)-localization of \(\mathcal R_S^0\).
\end{proof}

\begin{proposition}[Monoidality of BNV homotopification]
\label{prop:infinitesimal-bnv-monoidality}
The functor \(\mathbb H_S\) is colimit-preserving and strong symmetric
monoidal.  The fully faithful right adjoint \(\Disc_S\) carries its canonical
lax symmetric monoidal structure.  Consequently
\[
  \mathcal H_S:=\Disc_S\mathbb H_S
\]
is an exact idempotent localization monad with its canonical lax symmetric
monoidal structure.
\end{proposition}

\begin{proof}
Tensoring \(\alpha^0_{m,X,a}\) by a compact generator
\(F^0_{n,S}(\Sigma^b\Sigma^\infty_+h_Y)\) gives
\[
  \alpha^0_{m+n,X\times_SY,a+b}.
\]
The compact generators generate under colimits and tensor product preserves
colimits separately, so the affine-line equivalences form a tensor ideal in
\(\mathcal R_S^0\).  The symmetric monoidal localization theorem gives the
first assertion.  The second is the canonical lax structure on the right
adjoint of a strong symmetric monoidal functor.
\end{proof}

\begin{definition}[Stable BNV secondary coreflection]
\label{def:infinitesimal-bnv-secondary-coreflection}
Define
\[
    \mathcal S_S
    :=
    \Disc_S\mathbb S_S:
    \mathcal R_S^0\longrightarrow\mathcal R_S^0.
\]
We call $\mathcal S_S$ the \emph{stable BNV secondary coreflection}.
\end{definition}

\begin{proposition}[Idempotence of the secondary coreflection]
\label{prop:infinitesimal-bnv-secondary-idempotent}
There is a canonical equivalence
\[
    \mathcal S_S^2\xrightarrow{\simeq}\mathcal S_S.
\]
\end{proposition}

\begin{proof}
Full faithfulness of $\Disc_S$ gives
\[
    \mathbb S_S\Disc_S\simeq\id_{\SH(S)}.
\]
Therefore
\[
    \mathcal S_S^2
    \simeq
    \Disc_S\mathbb S_S\Disc_S\mathbb S_S
    \simeq
    \Disc_S\mathbb S_S
    =
    \mathcal S_S.
\]
\end{proof}

\section{The stable BNV differential-cohomology reconstruction}
The pullback theorem proved in this section is the brave-new counterpart of
the BNV reconstruction square
\cite[Proposition~3.5]{BunkeNikolausVoelkl}.  Here it is derived directly
from the exact adjacent adjunctions just constructed.
\label{sec:infinitesimal-bnv-reconstruction}

Put
\[
    \mathcal H_S:=\Disc_S\mathbb H_S.
\]
Recall that $\mathcal S_S=\Disc_S\mathbb S_S$ is the stable BNV secondary
coreflection.
Write
\[
    \eta_E:E\longrightarrow\mathcal H_SE
\]
for the localization unit and
\[
    \epsilon_E:\mathcal S_SE\longrightarrow E
\]
for the coreflection counit.

\begin{lemma}[Adjacent-mates comparison]
\label{lem:infinitesimal-bnv-adjacent-mates}
Under the canonical equivalences
\[
    \mathbb H_S\Disc_S\simeq\id_{\SH(S)},
    \qquad
    \mathbb S_S\Disc_S\simeq\id_{\SH(S)},
\]
the two natural morphisms
\[
    \mathbb H_S(\epsilon_E),\
    \mathbb S_S(\eta_E):
    \mathbb S_SE\longrightarrow\mathbb H_SE
\]
are canonically equivalent.  Denote their common value by
\[
    \chi_E:\mathbb S_SE\longrightarrow\mathbb H_SE.
\]
\end{lemma}

\begin{proof}
Apply the fully faithful functor $\Disc_S$.  Naturality of the localization
unit with respect to
\[
    \epsilon_E:\Disc_S\mathbb S_SE\longrightarrow E
\]
gives
\[
    \Disc_S\mathbb H_S(\epsilon_E)\circ
    \eta_{\Disc_S\mathbb S_SE}
    \simeq
    \eta_E\circ\epsilon_E.
\]
Since $\Disc_S\mathbb S_SE$ is local, its localization unit is an
equivalence, and therefore
\[
    \Disc_S\mathbb H_S(\epsilon_E)
    \simeq
    \eta_E\circ\epsilon_E
\]
under the canonical identification of the source.

Naturality of the coreflection counit with respect to
\[
    \eta_E:E\longrightarrow\Disc_S\mathbb H_SE
\]
gives
\[
    \eta_E\circ\epsilon_E
    \simeq
    \epsilon_{\Disc_S\mathbb H_SE}
    \circ
    \Disc_S\mathbb S_S(\eta_E).
\]
Since $\Disc_S\mathbb H_SE$ is in the essential image of $\Disc_S$, its
coreflection counit is an equivalence.  Hence
\[
    \Disc_S\mathbb S_S(\eta_E)
    \simeq
    \eta_E\circ\epsilon_E.
\]
Thus the images under $\Disc_S$ of
$\mathbb H_S(\epsilon_E)$ and $\mathbb S_S(\eta_E)$ agree canonically.
Full faithfulness of $\Disc_S$ supplies the asserted canonical equivalence.
\end{proof}

\begin{definition}[Underlying, secondary, deformation, and cycle sectors]
\label{def:infinitesimal-bnv-invariants}
For $E\in\mathcal R_S^0$ define
\[
    \mathsf U_S(E):=\mathbb H_SE,
    \qquad
    \mathsf S_S(E):=\mathbb S_SE,
\]
\[
    \mathcal A_S(E)
    :=
    \fib\bigl(E\xrightarrow{\eta_E}\mathcal H_SE\bigr),
\]
\[
    \mathcal Z_S(E)
    :=
    \cofib\bigl(\mathcal S_SE\xrightarrow{\epsilon_E}E\bigr),
\]
and
\[
    \mathsf Z_S(E)
    :=
    \mathbb H_S\mathcal Z_S(E).
\]
The objects $\mathsf U_S(E)$ and $\mathsf S_S(E)$ are the underlying and
secondary motivic theories; $\mathcal A_S(E)$ is the deformation sector;
$\mathcal Z_S(E)$ is the BNV-cycle sector; and $\mathsf Z_S(E)$ is the
localized cycle theory.
\end{definition}

There are canonical fibre and cofiber sequences
\[
    \mathcal A_S(E)
    \longrightarrow
    E
    \xrightarrow{I_E}
    \mathcal H_SE
\]
and
\[
    \mathcal S_SE
    \longrightarrow
    E
    \xrightarrow{R_E}
    \mathcal Z_S(E).
\]
Applying $\mathbb H_S$ to the second gives
\[
    \mathsf S_S(E)
    \longrightarrow
    \mathsf U_S(E)
    \xrightarrow{\phi_E}
    \mathsf Z_S(E),
\]
where $\phi_E$ is the BNV characteristic morphism.

\begin{proposition}[BNV-pure cycle objects and the cycle reflector]
\label{prop:infinitesimal-bnv-pure-cycles}
Define
\[
    \mathcal R_S^{\mathrm{BNV}}
    :=
    \{P\in\mathcal R_S^0\mid\mathbb S_SP\simeq0\}.
\]
Then \(\mathcal R_S^{\mathrm{BNV}}\) is a stable full subcategory of
\(\mathcal R_S^0\), and
\[
    \mathcal Z_S:\mathcal R_S^0
    \longrightarrow
    \mathcal R_S^{\mathrm{BNV}}
\]
is an exact reflector left adjoint to the inclusion.  In particular,
\[
    \mathbb S_S\mathcal Z_S(E)\simeq0,
    \qquad
    \mathcal Z_S^2(E)\xrightarrow{\simeq}\mathcal Z_S(E).
\]
The deformation and localization sectors satisfy the canonical equivalences
\[
    \boxed{
    \mathbb H_S\mathcal A_S(E)\simeq0,
    \qquad
    \mathcal Z_S\mathcal A_S(E)\xrightarrow{\simeq}\mathcal Z_S(E),
    \qquad
    \mathcal Z_S\mathcal H_S(E)\simeq0.
    }
\]
Thus the terminology ``BNV-pure'' used in this chapter is independent of the
coefficientwise closed-purity condition of Chapter~8.
\end{proposition}

\begin{proof}
The functor \(\mathbb S_S\) is exact, so its kernel
\(\mathcal R_S^{\mathrm{BNV}}\) is a stable full subcategory.  Apply
\(\mathbb S_S\) to
\[
    \mathcal S_SE\longrightarrow E\longrightarrow\mathcal Z_S(E).
\]
Since \(\mathbb S_S\Disc_S\simeq\id\), the first induced map is the identity
of \(\mathbb S_SE\).  Hence
\[
    \mathbb S_S\mathcal Z_S(E)\simeq0.
\]

Let \(P\in\mathcal R_S^{\mathrm{BNV}}\).  Applying \(\map(-,P)\) to the same
cofiber sequence and using \(\Disc_S\dashv\mathbb S_S\) gives
\[
\begin{aligned}
    \map(\mathcal S_SE,P)
    &\simeq
    \map_{\SH(S)}(\mathbb S_SE,\mathbb S_SP)\\
    &\simeq0.
\end{aligned}
\]
Therefore
\[
    \map(\mathcal Z_S(E),P)
    \xrightarrow{\simeq}
    \map(E,P),
\]
naturally in \(E\) and \(P\).  This proves the reflector adjunction.  The
functor \(\mathcal Z_S\) is exact because it is the cofiber of the natural
transformation between the exact functors \(\mathcal S_S\) and the identity.
Idempotence follows by applying the reflector to an already BNV-pure object.

Apply \(\mathbb H_S\) to the localization fibre sequence
\[
    \mathcal A_S(E)\longrightarrow E
    \xrightarrow{\eta_E}\mathcal H_SE.
\]
The morphism
\[
    \mathbb H_SE
    \longrightarrow
    \mathbb H_S\mathcal H_SE
    \simeq
    \mathbb H_SE
\]
is the localization of the unit and is an equivalence.  Hence
\[
    \mathbb H_S\mathcal A_S(E)\simeq0.
\]

Because \(\mathcal H_SE\) lies in the essential image of \(\Disc_S\), its
coreflection counit is an equivalence, and therefore
\[
    \mathcal Z_S\mathcal H_S(E)\simeq0.
\]
Finally apply the exact functor \(\mathcal Z_S\) to
\[
    \mathcal A_S(E)\longrightarrow E\longrightarrow\mathcal H_SE.
\]
The last term is sent to zero, so
\[
    \mathcal Z_S\mathcal A_S(E)
    \xrightarrow{\simeq}
    \mathcal Z_S(E).
\]
\end{proof}

\begin{proposition}[Presentability and colimits of the BNV-pure sector]
\label{prop:infinitesimal-bnv-pure-presentable}
Write
\[
    j_S:
    \mathcal R_S^{\mathrm{BNV}}
    \hookrightarrow
    \mathcal R_S^0
\]
for the inclusion.  The stable category
\(\mathcal R_S^{\mathrm{BNV}}\) is presentable, and
\[
    \mathcal Z_S:
    \mathcal R_S^0
    \longrightarrow
    \mathcal R_S^{\mathrm{BNV}}
\]
is an accessible exact localization.  For every small diagram
\[
    k\longmapsto P_k
    \qquad
    (P_k\in\mathcal R_S^{\mathrm{BNV}}),
\]
there is a canonical equivalence
\[
    \boxed{
    \colim_k^{\mathcal R_S^{\mathrm{BNV}}}P_k
    \simeq
    \mathcal Z_S
    \left(
      \colim_k^{\mathcal R_S^0}j_SP_k
    \right).
    }
\]
\end{proposition}

\begin{proof}
The categories \(\mathcal R_S^0\) and \(\SH(S)\) are presentable, and all
functors in the adjoint triple
\[
    \mathbb H_S\dashv\Disc_S\dashv\mathbb S_S
\]
are accessible.  Hence
\[
    \mathcal S_S=\Disc_S\mathbb S_S
\]
is accessible.  The cycle functor is the functorial cofiber
\[
    \mathcal Z_S
    =
    \cofib\bigl(\mathcal S_S\longrightarrow\id\bigr)
\]
in the functor category, so it is accessible.  By
Proposition~\ref{prop:infinitesimal-bnv-pure-cycles},
\(\mathcal Z_S\dashv j_S\) and \(j_S\) is fully faithful.  Thus
\(\mathcal R_S^{\mathrm{BNV}}\) is the essential image of an accessible
reflective localization of the presentable category \(\mathcal R_S^0\), and
is therefore presentable.  Its stability and exactness of the reflector were
proved in Proposition~\ref{prop:infinitesimal-bnv-pure-cycles}.

Finally, a left adjoint computes colimits in its target.  Applying the
reflector \(\mathcal Z_S\) to the ambient colimit of the diagram obtained through
\(j_S\) gives the displayed formula.
\end{proof}

\begin{theorem}[Upper BNV triangle]
\label{thm:infinitesimal-bnv-upper-triangle}
There is a canonical fibre and cofiber sequence
\[
    \boxed{
    \Disc_S\Omega\mathsf Z_S(E)
    \longrightarrow
    \mathcal A_S(E)
    \longrightarrow
    \mathcal Z_S(E).
    }
\]
\end{theorem}

\begin{proof}
Applying the exact functor $\mathbb S_S$ to the localization fibre sequence
\[
    \mathcal A_S(E)\longrightarrow E
    \xrightarrow{\eta_E}\mathcal H_SE
\]
gives a fibre sequence
\[
    \mathbb S_S\mathcal A_S(E)
    \longrightarrow
    \mathbb S_SE
    \xrightarrow{\mathbb S_S(\eta_E)}
    \mathbb H_SE.
\]
Applying the exact functor $\mathbb H_S$ to the cycle cofiber sequence gives
\[
    \mathbb S_SE
    \xrightarrow{\mathbb H_S(\epsilon_E)}
    \mathbb H_SE
    \longrightarrow
    \mathsf Z_S(E).
\]
By Lemma~\ref{lem:infinitesimal-bnv-adjacent-mates}, the two middle arrows
are the same canonical morphism $\chi_E$.  Consequently
\[
    \mathbb S_S\mathcal A_S(E)
    \simeq
    \fib(\chi_E)
    \simeq
    \Omega\cofib(\chi_E)
    \simeq
    \Omega\mathsf Z_S(E).
\]

The endofunctor $\mathcal Z_S$ is exact, being the cofiber of the natural
transformation between the exact functors $\mathcal S_S$ and the identity.
Apply it to the localization fibre sequence.  Since $\mathcal H_SE$ is local,
\[
    \mathcal Z_S(\mathcal H_SE)\simeq0,
\]
and therefore
\[
    \mathcal Z_S(\mathcal A_S(E))
    \xrightarrow{\simeq}
    \mathcal Z_S(E).
\]
The defining cycle sequence for $\mathcal A_S(E)$ is
\[
    \Disc_S\mathbb S_S\mathcal A_S(E)
    \longrightarrow
    \mathcal A_S(E)
    \longrightarrow
    \mathcal Z_S(\mathcal A_S(E)).
\]
Substituting the two canonical equivalences just obtained gives
\[
    \Disc_S\Omega\mathsf Z_S(E)
    \longrightarrow
    \mathcal A_S(E)
    \longrightarrow
    \mathcal Z_S(E),
\]
which is both a fibre and a cofiber sequence because $\mathcal R_S^0$ is
stable.
\end{proof}

\begin{theorem}[Brave-new BNV reconstruction square]
\label{thm:infinitesimal-bnv-reconstruction}
For every $E\in\mathcal R_S^0$, the square
\[
\begin{tikzcd}[column sep=large,row sep=large]
    E \ar[r,"R_E"] \ar[d,"\eta_E"']
        & \mathcal Z_S(E) \ar[d,"\eta_{\mathcal Z_S(E)}"]\\
    \Disc_S\mathsf U_S(E)
        \ar[r,"\Disc_S\phi_E"']
        & \Disc_S\mathsf Z_S(E)
\end{tikzcd}
\]
is a pullback square.  Equivalently,
\[
    \boxed{
    E
    \simeq
    \mathcal Z_S(E)
    \times_{\Disc_S\mathsf Z_S(E)}
    \Disc_S\mathsf U_S(E).
    }
\]
\end{theorem}

\begin{proof}
The square commutes by naturality of the localization unit.  The fibre of the
upper horizontal morphism is $\mathcal S_SE$.  By the exact triangle obtained
by applying $\mathbb H_S$ to the cycle sequence, and by
Lemma~\ref{lem:infinitesimal-bnv-adjacent-mates}, the fibre of the
characteristic morphism is canonically
\[
    \fib(\phi_E)\simeq\mathbb S_SE.
\]
Since $\Disc_S$ is exact, the fibre of the lower horizontal morphism is
therefore
\[
    \Disc_S\fib(\phi_E)
    \simeq
    \Disc_S\mathbb S_SE
    =
    \mathcal S_SE.
\]
Under these identifications, the induced comparison of horizontal fibres is
the localization unit of the already local object $\mathcal S_SE$, hence an
equivalence.  In a stable $\infty$-category a square is a pullback precisely
when the induced map of horizontal fibres is an equivalence.
\end{proof}

\begin{definition}[Characteristic pullback category]
\label{def:infinitesimal-bnv-characteristic-data}
Let $\mathcal T_S$ be the pullback $\infty$-category
\[
\begin{tikzcd}[column sep=large,row sep=large]
    \mathcal T_S \ar[r] \ar[d]
        & \Fun(\Delta^1,\SH(S))
          \ar[d,"{(\operatorname{ev}_0,\operatorname{ev}_1)}"]\\
    \SH(S)\times\mathcal R_S^{\mathrm{BNV}}
        \ar[r,"{(\id,\,\mathbb H_S)}"']
        & \SH(S)\times\SH(S).
\end{tikzcd}
\]
Thus an object is a triple $(U,P,\phi)$ with
\[
    U\in\SH(S),
    \qquad
    P\in\mathcal R_S^{\mathrm{BNV}},
    \qquad
    \phi:U\longrightarrow\mathbb H_SP.
\]
\end{definition}

\begin{theorem}[BNV lax-pullback classification]
\label{thm:infinitesimal-bnv-classification}
There is a canonical equivalence of $\infty$-categories
\[
    \boxed{
    \mathcal R_S^0\simeq\mathcal T_S.
    }
\]
The inverse construction is
\[
    E(U,P,\phi)
    :=
    P
    \times_{\Disc_S\mathbb H_SP}
    \Disc_SU.
\]
For this object,
\[
    \mathsf U_S(E(U,P,\phi))\simeq U,
    \qquad
    \mathcal Z_S(E(U,P,\phi))\simeq P,
\]
and under these identifications its characteristic morphism is canonically
\[
    \phi_{E(U,P,\phi)}\simeq\phi.
\]
\end{theorem}

\begin{proof}
The reconstruction theorem sends $E$ functorially to
\[
    \left(
      \mathsf U_S(E),
      \mathcal Z_S(E),
      \phi_E
    \right)
\]
and reconstructs $E$ from that triple.

Conversely, put
\[
    E:=E(U,P,\phi)
    =P\times_{\Disc_S\mathbb H_SP}\Disc_SU.
\]
Exactness of $\mathbb H_S$ gives
\[
\begin{aligned}
    \mathbb H_SE
    &\simeq
    \mathbb H_SP\times_{\mathbb H_SP}U\\
    &\simeq U.
\end{aligned}
\]
Since $\mathbb S_S$ preserves limits, $\mathbb S_S\Disc_S\simeq\id$, and
$P$ is BNV-pure,
\[
    \mathbb S_SE
    \simeq
    0\times_{\mathbb H_SP}U
    \simeq
    \fib(\phi).
\]
Let
\[
    \pi_P:E\longrightarrow P
\]
be the first projection.  Its fibre is canonically
\[
    \fib(\pi_P)
    \simeq
    \Disc_S\fib(\phi)
    \simeq
    \Disc_S\mathbb S_SE.
\]
The composite
\[
    \Disc_S\mathbb S_SE
    \xrightarrow{\epsilon_E}
    E
    \xrightarrow{\pi_P}
    P
\]
corresponds under $\Disc_S\dashv\mathbb S_S$ to a morphism
$\mathbb S_SE\to\mathbb S_SP\simeq0$, and is therefore null.  Hence the
coreflection counit factors through $\fib(\pi_P)$.  Under the preceding
identification, applying $\mathbb S_S$ to this factorization gives the
identity of $\mathbb S_SE$; consequently the factorization is an equivalence
because both source and target lie in the essential image of $\Disc_S$.
Thus
\[
    \Disc_S\mathbb S_SE
    \longrightarrow
    E
    \xrightarrow{\pi_P}
    P
\]
is the cycle cofiber sequence, and therefore
\[
    \mathcal Z_S(E)\simeq P.
\]

It remains to recover the characteristic morphism.  Under the canonical
identification
\[
    \mathbb H_SE
    \simeq
    \mathbb H_SP\times_{\mathbb H_SP}U
    \simeq U,
\]
the morphism $\mathbb H_S(\pi_P)$ is the first projection
\[
    \mathbb H_SP\times_{\mathbb H_SP}U
    \longrightarrow
    \mathbb H_SP.
\]
The canonical equivalence from $U$ to this pullback is induced by the pair
$(\phi,\id_U)$.  Hence the first projection is exactly $\phi$.  Since
$\pi_P$ identifies with the cycle morphism $R_E$ under
$\mathcal Z_S(E)\simeq P$, one obtains
\[
    \phi_E=\mathbb H_S(R_E)\simeq\phi.
\]
The constructions are functorial in morphisms of triples.  Together with the
reconstruction square, these identifications show that the two displayed
constructions are mutually inverse equivalences of $\infty$-categories.
\end{proof}

\begin{proposition}[Colimits of characteristic triples]
\label{prop:infinitesimal-bnv-characteristic-colimits}
Let \(K\) be a small \(\infty\)-category and let
\[
    k\longmapsto(U_k,P_k,\phi_k)
\]
be a \(K\)-diagram in \(\mathcal T_S\).  Put
\[
    U:=\colim_k^{\SH(S)}U_k,
    \qquad
    P:=\colim_k^{\mathcal R_S^{\mathrm{BNV}}}P_k.
\]
The cocone maps \(P_k\to P\) induce a canonical comparison
\[
    \lambda_P:
    \colim_k\mathbb H_SP_k
    \longrightarrow
    \mathbb H_SP.
\]
Let
\[
    \overline\phi:
    U
    \simeq
    \colim_kU_k
    \longrightarrow
    \colim_k\mathbb H_SP_k
\]
be induced by the \(\phi_k\).  Then the colimit of the diagram in
\(\mathcal T_S\) is canonically
\[
    \boxed{
    \colim_k^{\mathcal T_S}(U_k,P_k,\phi_k)
    \simeq
    \bigl(
      U,\,
      P,\,
      \lambda_P\circ\overline\phi
    \bigr).
    }
\]
\end{proposition}

\begin{proof}
For each \(k\), let
\[
    E_k\in\mathcal R_S^0
\]
be the object corresponding to \((U_k,P_k,\phi_k)\) under
Theorem~\ref{thm:infinitesimal-bnv-classification}, and form
\[
    E:=\colim_k^{\mathcal R_S^0}E_k.
\]
Since \(\mathbb H_S\) is a left adjoint,
\[
    \mathbb H_SE
    \simeq
    \colim_k\mathbb H_SE_k
    \simeq
    \colim_kU_k
    =
    U.
\]
Since
\[
    \mathcal Z_S:
    \mathcal R_S^0
    \longrightarrow
    \mathcal R_S^{\mathrm{BNV}}
\]
is also a left adjoint,
\[
    \mathcal Z_SE
    \simeq
    \colim_k^{\mathcal R_S^{\mathrm{BNV}}}\mathcal Z_SE_k
    \simeq
    \colim_k^{\mathcal R_S^{\mathrm{BNV}}}P_k
    =
    P.
\]

For each \(k\), naturality of the cycle map gives a commutative square
\[
\begin{tikzcd}[column sep=large,row sep=large]
    E_k \ar[r,"R_{E_k}"] \ar[d]
        & P_k \ar[d]\\
    E \ar[r,"R_E"']
        & P.
\end{tikzcd}
\]
Applying \(\mathbb H_S\), and using
\(\phi_{E_k}=\mathbb H_S(R_{E_k})\), shows that the restriction of
\(\phi_E=\mathbb H_S(R_E)\) along \(U_k\to U\) is the composite
\[
    U_k
    \xrightarrow{\phi_k}
    \mathbb H_SP_k
    \longrightarrow
    \mathbb H_SP.
\]
By the universal property of \(U=\colim_kU_k\), this identifies
\[
    \phi_E
    \simeq
    \lambda_P\circ\overline\phi.
\]
Hence, under the equivalence of
Theorem~\ref{thm:infinitesimal-bnv-classification}, the object \(E\)
corresponds to
\[
    \bigl(
      U,\,
      P,\,
      \lambda_P\circ\overline\phi
    \bigr).
\]
Therefore the displayed triple is the required colimit in
\(\mathcal T_S\).
\end{proof}

\section{The affine line and universal BNV integration}
\label{sec:infinitesimal-bnv-integration}

The interval object in the Thom-stable differential category is the spectral
affine line \(\mathbb A^1_S\).  The Thom coordinate has already been stabilized
by \(L_{\mathrm{Th},S}\); the interval and stabilization coordinate therefore
play distinct roles.

Let
\[
  q_S:\mathbb A^1_S\longrightarrow S
\]
be affine-line projection.  Composition and base change define
\[
  c_{q_S}:\operatorname{Sm}^{\fp}_{/\mathbb A^1_S}
  \longrightarrow\operatorname{Sm}^{\fp}_{/S},
  \qquad
  b_{q_S}:\operatorname{Sm}^{\fp}_{/S}
  \longrightarrow\operatorname{Sm}^{\fp}_{/\mathbb A^1_S},
\]
with \(c_{q_S}\dashv b_{q_S}\).

\begin{proposition}[Affine-line adjunction on Thom-stable coefficients]
\label{prop:infinitesimal-raw-affine-line-adjunction}
There is an exact adjunction
\[
  q_S^{*,0}:
  \mathcal R_S^0
  \rightleftarrows
  \mathcal R_{\mathbb A^1_S}^0:
  q_{S,*}^{0}.
\]
The left adjoint is strong symmetric monoidal and satisfies
\[
  q_S^{*,0}\mathbb T_S^0
  \simeq
  \mathbb T_{\mathbb A^1_S}^0.
\]
Under the fully faithful inclusions of the Thom-stable categories into their
raw prespectrum presentation categories, the adjunction is induced levelwise by
\[
  (q_S^*F)(Y)=F(c_{q_S}Y),
  \qquad
  (q_{S,*}G)(X)=G(X\times_S\mathbb A^1_S).
\]
\end{proposition}

\begin{proof}
On stable spectral Nisnevich sheaves, restriction along
\(c_{q_S}^{\op}\) is left Kan extension along \(b_{q_S}^{\op}\), because
\(b_{q_S}^{\op}\dashv c_{q_S}^{\op}\).  Its right adjoint is restriction along
\(b_{q_S}^{\op}\).  Base change and composition preserve Nisnevich
distinguished squares, so both functors restrict to sheaves.  The base-change
functor preserves products and terminal objects, hence its left Kan extension
is strong symmetric monoidal.  Exactness follows from stability.

Base change carries \(\mathbb A^1_S\) and \(\mathbb G_{m,S}\) to the
corresponding objects over \(\mathbb A^1_S\), so
\[
  q_S^*\mathbb T_S^{\mathrm{raw}}
  \simeq
  \mathbb T_{\mathbb A^1_S}^{\mathrm{raw}}.
\]
The adjunction therefore lifts levelwise to raw symmetric Thom prespectra, and
strong monoidality identifies
\[
  q_S^*(\lambda_{m,G})
  \simeq
  \lambda_{m,q_S^*G}.
\]
Thus the raw left adjoint preserves Thom-stabilization equivalences.

To see that it also preserves Thom-local objects, use smoothness of
$q_S$.  Composition with $q_S$ gives a stable Nisnevich smooth direct image
\[
  q_{S,\sharp,\Nis}:
  \mathcal N_{\mathbb A^1_S}^{\Sp}
  \rightleftarrows
  \mathcal N_S^{\Sp}:
  q_S^*.
\]
It sends a suspension representable over $\mathbb A^1_S$ to the same smooth
spectral scheme regarded over $S$.  The smooth projection formula gives
\[
  q_{S,\sharp,\Nis}
  (\mathbb T_{\mathbb A^1_S}^{\mathrm{raw}}\otimes A)
  \simeq
  \mathbb T_S^{\mathrm{raw}}\otimes q_{S,\sharp,\Nis}A,
\]
so this adjunction lifts to raw symmetric Thom prespectra.  If $E$ is
Thom-local over $S$ and $\lambda_{m,G}$ is a stabilization generator over
$\mathbb A^1_S$, then
\[
\begin{aligned}
  \map(\lambda_{m,G},q_S^*E)
  &\simeq
  \map(q_{S,\sharp}\lambda_{m,G},E).
\end{aligned}
\]
The raw smooth direct image identifies
\[
  q_{S,\sharp}\lambda_{m,G}
  \simeq
  \lambda_{m,q_{S,\sharp,\Nis}G}.
\]
For $G\in\mathcal G_{\mathbb A^1_S}$ the object
$q_{S,\sharp,\Nis}G$ belongs to $\mathcal G_S$, because the composite of a
smooth finite-presentation morphism with $q_S$ is again smooth of finite
presentation.  The displayed mapping map is therefore an equivalence by
Thom-locality of $E$.  Hence the raw $q_S^*$ carries Thom-local objects to
Thom-local objects.

The raw right adjoint also preserves Thom-local objects.  If $E$ is Thom-local
over $\mathbb A^1_S$, adjunction identifies mapping a stabilization generator
$\lambda_{m,G}$ over $S$ into $q_{S,*}E$ with mapping
\[
  q_S^*(\lambda_{m,G})
  \simeq
  \lambda_{m,q_S^*G}
\]
into $E$, which is an equivalence.  Thus both sides of the raw adjunction
restrict to the Thom-local subcategories.  This gives the displayed levelwise
adjunction on $\mathcal R^0$, rather than merely an adjunction obtained after
applying the reflector.  The Thom-coordinate formula follows from the raw
coordinate equivalence and strong monoidality.
\end{proof}

\begin{definition}[Affine-line cylinder]
\label{def:infinitesimal-affine-line-cylinder}
Define the exact affine-line cylinder
\[
  \boxed{
  \mathbb I_S
  :=
  q_{S,*}^0q_S^{*,0}:
  \mathcal R_S^0\longrightarrow\mathcal R_S^0.
  }
\]
The adjunction unit gives
\[
  \operatorname{pr}^*:\id\longrightarrow\mathbb I_S.
\]
The zero and unit sections give endpoint transformations
\[
  0^*,1^*:\mathbb I_S\longrightarrow\id
\]
with specified homotopies
\[
  0^*\operatorname{pr}^*\simeq\id,
  \qquad
  1^*\operatorname{pr}^*\simeq\id.
\]
On a Thom-stable prespectrum represented in the raw presentation category,
\[
  (\mathbb I_SE)_m(X)
  \simeq
  E_m(X\times_S\mathbb A^1_S).
\]
\end{definition}

\begin{lemma}[Cylinder criterion for motivic locality]
\label{lem:infinitesimal-cylinder-locality}
For $E\in\mathcal R_S^0$, the following conditions are equivalent:
\begin{enumerate}
\item $E$ is motivic local;
\item the affine-line projection unit is an equivalence
\[
    \boxed{
    \operatorname{pr}^*:E\xrightarrow{\simeq}\mathbb I_SE.
    }
\]
\end{enumerate}
\end{lemma}

\begin{proof}
Let
\[
  K_{m,X,a}
  :=
  F^0_{m,S}\bigl(\Sigma^a\Sigma^\infty_+h_X\bigr).
\]
For a Thom-local object $E$, the localization adjunction and the free-level
adjunction give
\[
\begin{aligned}
  \map_{\mathcal R_S^0}(K_{m,X,a},E)
  &\simeq
  \map_{\mathcal N_S^{\Sp}}
  \bigl(\Sigma^a\Sigma^\infty_+h_X,E_m\bigr).
\end{aligned}
\]
Therefore locality of $E$ with respect to
$\alpha^0_{m,X,a}$ is exactly the condition that
\[
  E_m(X)
  \longrightarrow
  E_m(X\times_S\mathbb A^1_S)
\]
be an equivalence for every $m,X,a$.  By
Proposition~\ref{prop:infinitesimal-raw-affine-line-adjunction}, these are
precisely the levelwise components of
\[
  \operatorname{pr}^*:E\longrightarrow\mathbb I_SE.
\]
The fully faithful inclusion of $\mathcal R_S^0$ into the raw prespectrum
category detects equivalences levelwise.  Hence $E$ is local for all the
$\alpha^0_{m,X,a}$ if and only if $\operatorname{pr}^*$ is an equivalence.
\end{proof}

\begin{lemma}[Normalized path between specified inverses]
\label{lem:infinitesimal-normalized-inverse-path}
Let $u:A\xrightarrow{\simeq}B$ be an equivalence in an $\infty$-category and
let $r_0,r_1:B\to A$ be morphisms equipped with specified homotopies
\[
    \rho_0:r_0u\simeq\id_A,
    \qquad
    \rho_1:r_1u\simeq\id_A.
\]
Then the space of paths $h:r_0\simeq r_1$ equipped with a coherent homotopy,
relative to $\rho_0$ and $\rho_1$, from $h\circ u$ to the constant path at
$\id_A$ is contractible.
\end{lemma}

\begin{proof}
Precomposition with $u$ gives an equivalence
\[
    u^*:\Map(B,A)\xrightarrow{\simeq}\Map(A,A).
\]
Let
\[
    K:=\fib_{\id_A}(u^*).
\]
Since $u^*$ is an equivalence, $K$ is contractible.  The specified pairs
$(r_0,\rho_0)$ and $(r_1,\rho_1)$ are points of $K$, so the path space between
them in $K$ is contractible.  Unwinding the definition of a path in this
homotopy fibre gives precisely a path $h:r_0\simeq r_1$ together with the
stated coherent normalization relative to the specified endpoint
homotopies.
\end{proof}

\begin{proposition}[Canonical normalized endpoint path on motivic-local objects]
\label{prop:infinitesimal-local-endpoint-path}
For every $M\in\SH(S)$, the projection map
\[
    \operatorname{pr}^*:
    \Disc_SM
    \xrightarrow{\simeq}
    \mathbb I_S\Disc_SM
\]
is an equivalence.  There is a contractibly unique natural path
\[
    h_M^{\mathrm{loc}}:0^*\simeq1^*
\]
between the endpoint maps
\[
    0^*,1^*:
    \mathbb I_S\Disc_SM
    \longrightarrow
    \Disc_SM
\]
whose precomposition with $\operatorname{pr}^*$ is coherently the constant
path at the identity relative to the two specified retraction homotopies.
\end{proposition}

\begin{proof}
The object $\Disc_SM$ is motivic local, so
Lemma~\ref{lem:infinitesimal-cylinder-locality} gives a natural equivalence
\[
    \operatorname{pr}^*:
    \Disc_SM
    \xrightarrow{\simeq}
    \mathbb I_S\Disc_SM.
\]
Apply
Lemma~\ref{lem:infinitesimal-normalized-inverse-path} in the functor category
\[
    \Fun(\SH(S),\mathcal R_S^0)
\]
to the natural equivalence $\Disc_S\to\mathbb I_S\Disc_S$ and its two
specified retractions $0^*$ and $1^*$.  Working in the functor category
supplies naturality and all higher coherences.
\end{proof}

\begin{lemma}[Mapping out of cofibers and into fibers]
\label{lem:infinitesimal-bnv-mapping-fiber-cofiber}
Let $\mathcal D$ be a stable $\infty$-category.
\begin{enumerate}
\item For every cofiber sequence
\[
    A\xrightarrow{u}B\xrightarrow{q}C
\]
and every $Z\in\mathcal D$, there is a canonical equivalence
\[
    \Map_{\mathcal D}(C,Z)
    \simeq
    \fib_{0}
    \left(
      \Map_{\mathcal D}(B,Z)
      \xrightarrow{u^*}
      \Map_{\mathcal D}(A,Z)
    \right).
\]
Thus a morphism $C\to Z$ is equivalently a morphism $b:B\to Z$ together
with a specified nullhomotopy of $bu$.
\item For every fiber sequence
\[
    F\xrightarrow{a}Y\xrightarrow{v}Z
\]
and every $X\in\mathcal D$, there is a canonical equivalence
\[
    \Map_{\mathcal D}(X,F)
    \simeq
    \fib_{0}
    \left(
      \Map_{\mathcal D}(X,Y)
      \xrightarrow{v_*}
      \Map_{\mathcal D}(X,Z)
    \right).
\]
Thus a morphism $X\to F$ is equivalently a morphism $x:X\to Y$ together
with a specified nullhomotopy of $vx$.
\end{enumerate}
Both equivalences are natural in all variables.
\end{lemma}

\begin{proof}
Applying $\Map_{\mathcal D}(-,Z)$ to the cofiber square defining $C$ turns
that colimit square into a limit square of spaces and gives the first
equivalence.  Applying $\Map_{\mathcal D}(X,-)$ to the fiber square defining
$F$ preserves that limit square and gives the second.  Naturality is inherited
from functoriality of mapping spaces.
\end{proof}

\begin{lemma}[Normalized BNV factorization data]
\label{lem:infinitesimal-bnv-normalized-factorization}
Put
\[
    \delta:=1^*-0^*:\mathbb I_S\longrightarrow\id.
\]
The normalized endpoint path of
Proposition~\ref{prop:infinitesimal-local-endpoint-path} canonically
determines, through contractible spaces of compatible choices, a natural
transformation
\[
    \alpha:
    \mathbb I_S\mathcal Z_S
    \longrightarrow
    \id
\]
and a specified nullhomotopy
\[
    \overline\mu:
    \eta\alpha\simeq0
\]
such that
\[
    \alpha\circ\mathbb I_SR
    \simeq
    \delta.
\]
Moreover the construction is normalized in the sense that
\[
    \alpha\circ\operatorname{pr}^*_{\mathcal Z_S}
    \simeq0,
\]
with its canonical basepoint nullhomotopy, and
\[
    \overline\mu\circ\operatorname{pr}^*_{\mathcal Z_S}
\]
is the constant nullhomotopy of that zero transformation.
\end{lemma}

\begin{proof}
Work in the stable functor $\infty$-category
\[
    \mathcal F_S
    :=
    \Fun(\mathcal R_S^0,\mathcal R_S^0).
\]
The cycle and deformation constructions give a cofiber sequence and a fibre
sequence
\[
    \mathcal S_S
    \xrightarrow{\epsilon}
    \id
    \xrightarrow{R}
    \mathcal Z_S
\]
and
\[
    \mathcal A_S
    \xrightarrow{a}
    \id
    \xrightarrow{\eta}
    \mathcal H_S.
\]
Since $\mathbb I_S$ is exact, the first becomes a cofiber sequence
\[
    \mathbb I_S\mathcal S_S
    \xrightarrow{\mathbb I_S\epsilon}
    \mathbb I_S
    \xrightarrow{\mathbb I_SR}
    \mathbb I_S\mathcal Z_S.
\]

The normalized endpoint path on the essential image of $\Disc_S$ gives a
specified natural nullhomotopy
\[
    \kappa^{\mathrm{loc}}:
    \delta\Disc_S\simeq0
\]
in $\Fun(\SH(S),\mathcal R_S^0)$.  Whiskering with $\mathbb S_S$ and
$\mathbb H_S$ gives specified nullhomotopies
\[
    \kappa_{\mathcal S}:
    \delta\mathcal S_S\simeq0,
    \qquad
    \kappa_{\mathcal H}:
    \delta\mathcal H_S\simeq0.
\]
Their normalization says that after precomposition with the corresponding
projection transformation $\operatorname{pr}^*$ they are the constant
nullhomotopies relative to the two specified endpoint retraction homotopies.

Naturality of $\delta$ with respect to
$\epsilon:\mathcal S_S\to\id$ gives
\[
    \delta\circ\mathbb I_S\epsilon
    =
    \epsilon\circ\delta\mathcal S_S.
\]
Postcomposing $\kappa_{\mathcal S}$ with $\epsilon$ therefore gives a
specified nullhomotopy
\[
    \nu:
    \delta\circ\mathbb I_S\epsilon
    \simeq0.
\]
By the cofiber part of
Lemma~\ref{lem:infinitesimal-bnv-mapping-fiber-cofiber}, the pair
$(\delta,\nu)$ determines, through a contractible space of compatible
choices, a natural transformation
\[
    \alpha:
    \mathbb I_S\mathcal Z_S
    \longrightarrow
    \id
\]
with a specified homotopy
\[
    \alpha\circ\mathbb I_SR
    \simeq
    \delta.
\]

We now construct the nullhomotopy of $\eta\alpha$.  Naturality of $\delta$
with respect to $\eta:\id\to\mathcal H_S$ gives
\[
    \eta\circ\delta
    =
    \delta\mathcal H_S\circ\mathbb I_S\eta.
\]
Precomposing $\kappa_{\mathcal H}$ with $\mathbb I_S\eta$ gives
\[
    \mu_0:
    \eta\delta\simeq0.
\]

The composite
\[
    \eta\epsilon:
    \mathcal S_S\longrightarrow\mathcal H_S
\]
takes values in the essential image of $\Disc_S$.  By
Lemma~\ref{lem:infinitesimal-bnv-adjacent-mates}, it is the image under
$\Disc_S$ of the natural characteristic transformation
$\mathbb S_S\to\mathbb H_S$.  Naturality of the normalized endpoint path on
the essential image of $\Disc_S$ therefore supplies the coherent
identification
\[
    (\eta\epsilon)\circ\kappa_{\mathcal S}
    \simeq
    \kappa_{\mathcal H}\circ
    \mathbb I_S(\eta\epsilon).
\]
The left-hand side is $\eta\circ\nu$, while the right-hand side is
$\mu_0\circ\mathbb I_S\epsilon$.  Hence
\[
    \mu_0\circ\mathbb I_S\epsilon
    \simeq
    \eta\circ\nu.
\]
Under the cofiber mapping equivalence, this is precisely the compatibility
required for the nullhomotopy $\mu_0$ of $\eta\delta$ to descend to a
specified nullhomotopy
\[
    \overline\mu:
    \eta\alpha\simeq0.
\]
The space of such compatible descended nullhomotopies is contractible because
the cofiber mapping equivalence is an equivalence of spaces.

It remains to record the normalization.  Naturality of
$\operatorname{pr}^*:\id\to\mathbb I_S$ with respect to $R$ gives
\[
    \mathbb I_SR\circ\operatorname{pr}^*
    =
    \operatorname{pr}^*_{\mathcal Z_S}\circ R.
\]
Under the cofiber mapping equivalence defining $\alpha$, the composite
\[
    \alpha\circ\operatorname{pr}^*_{\mathcal Z_S}
\]
is represented by
\[
    \delta\circ\operatorname{pr}^*:\id\longrightarrow\id
\]
together with the nullhomotopy induced by $\nu$.  The two specified endpoint
retraction homotopies identify
$\delta\circ\operatorname{pr}^*$ with zero, and the normalization of
$\kappa_{\mathcal S}$ identifies the induced nullhomotopy on
$\mathcal S_S$ with the constant one.  Thus the corresponding point of the
homotopy fibre is its basepoint, and
\[
    \alpha\circ\operatorname{pr}^*_{\mathcal Z_S}
    \simeq0
\]
with its canonical basepoint nullhomotopy.

Finally, the normalization of $\kappa_{\mathcal H}$, together with the
coherent naturality just used for $\eta\epsilon$, identifies
\[
    \overline\mu\circ\operatorname{pr}^*_{\mathcal Z_S}
\]
with the constant nullhomotopy of this zero transformation.  This proves all
assertions.
\end{proof}

\begin{theorem}[Universal brave-new BNV integration]
\label{thm:infinitesimal-universal-integration}
For every $E\in\mathcal R_S^0$ there is a canonical natural transformation
\[
    \boxed{
    \int_{\mathbb A^1}^{\mathrm{BNV}}:
    \mathbb I_S\mathcal Z_S(E)
    \longrightarrow
    \mathcal A_S(E).
    }
\]
It satisfies the universal homotopy formula
\[
    \boxed{
    1^*-0^*
    \simeq
    a_E
    \circ
    \int_{\mathbb A^1}^{\mathrm{BNV}}
    \circ
    \mathbb I_SR_E,
    }
\]
as maps $\mathbb I_SE\to E$, where
$a_E:\mathcal A_S(E)\to E$ is the deformation inclusion.  It is normalized
by
\[
    \int_{\mathbb A^1}^{\mathrm{BNV}}
    \circ
    \operatorname{pr}^*_{\mathcal Z_S(E)}
    \simeq0.
\]
More precisely, after fixing the normalized endpoint contractions of
Proposition~\ref{prop:infinitesimal-local-endpoint-path}, the cofiber
factorization and the subsequent lift through the deformation fibre used to
construct $\int_{\mathbb A^1}^{\mathrm{BNV}}$ each lie in a contractible space of compatible
choices.
\end{theorem}

\begin{proof}
Work in the stable functor $\infty$-category
\[
    \mathcal F_S
    :=
    \Fun(\mathcal R_S^0,\mathcal R_S^0).
\]
By Lemma~\ref{lem:infinitesimal-bnv-normalized-factorization}, the normalized
endpoint path determines a natural transformation
\[
    \alpha:
    \mathbb I_S\mathcal Z_S
    \longrightarrow
    \id
\]
and a specified nullhomotopy
\[
    \overline\mu:
    \eta\alpha\simeq0
\]
through contractible spaces of compatible choices, with
\[
    \alpha\circ\mathbb I_SR
    \simeq
    1^*-0^*.
\]

Apply the fibre part of
Lemma~\ref{lem:infinitesimal-bnv-mapping-fiber-cofiber} in $\mathcal F_S$ to
the fibre sequence
\[
    \mathcal A_S
    \xrightarrow{a}
    \id
    \xrightarrow{\eta}
    \mathcal H_S.
\]
The pair $(\alpha,\overline\mu)$ determines, through a contractible space of
compatible choices, a natural transformation
\[
    \int_{\mathbb A^1}^{\mathrm{BNV}}:
    \mathbb I_S\mathcal Z_S
    \longrightarrow
    \mathcal A_S
\]
whose composite with $a$ is $\alpha$.  Therefore
\[
    a\circ
    \int_{\mathbb A^1}^{\mathrm{BNV}}
    \circ
    \mathbb I_SR
    \simeq
    \alpha\circ\mathbb I_SR
    \simeq
    1^*-0^*.
\]
Evaluating at $E$ gives the universal homotopy formula.

For the normalization, Lemma~\ref{lem:infinitesimal-bnv-normalized-factorization}
gives
\[
    \alpha\circ\operatorname{pr}^*_{\mathcal Z_S}
    \simeq0
\]
with its canonical basepoint nullhomotopy, and identifies
\[
    \overline\mu\circ\operatorname{pr}^*_{\mathcal Z_S}
\]
with the constant nullhomotopy of that zero transformation.  Under the fibre
mapping equivalence used above, this pair is the basepoint.  Hence
\[
    \int_{\mathbb A^1}^{\mathrm{BNV}}
    \circ
    \operatorname{pr}^*_{\mathcal Z_S}
    \simeq0.
\]
Evaluating at $E$ yields the stated normalization.  Since the construction is
performed in the functor $\infty$-category and both mapping-space
factorizations are equivalences, naturality, higher coherence, and the
asserted contractible uniqueness are part of the construction.
\end{proof}

\section*{Part V: Thom-stable functoriality and the smooth Umkehr lift}
\addcontentsline{toc}{section}{Part V: Thom-stable functoriality and the smooth Umkehr lift}

\section{Base functoriality of Thom-stable differential coefficients}
\label{sec:bn-diff-thom-stable-base-functoriality}

Let \(f:T\to S\) be a morphism of qcqs spectral schemes in the standing
motivic scope.  Base change defines
\[
  u_f:\operatorname{Sm}^{\fp}_{/S}
  \longrightarrow
  \operatorname{Sm}^{\fp}_{/T},
  \qquad
  X\longmapsto X_T.
\]

\begin{theorem}[Inverse and direct image before affine-line localization]
\label{thm:bn-diff-thom-stable-base-adjunction}
There is a canonical exact adjunction
\[
  \boxed{
  f^{*,0}:
  \mathcal R_S^0
  \rightleftarrows
  \mathcal R_T^0:
  f_*^0.
  }
\]
The functor \(f^{*,0}\) preserves all small colimits, is strong symmetric
monoidal, and satisfies
\[
  f^{*,0}\mathbb T_S^0\simeq\mathbb T_T^0.
\]
For composable morphisms \(U\xrightarrow{g}T\xrightarrow{f}S\) there are
canonical coherent equivalences
\[
  (fg)^{*,0}\simeq g^{*,0}f^{*,0},
  \qquad
  (fg)_*^0\simeq f_*^0g_*^0.
\]
After affine-line homotopification these are the base functors of Chapter~6:
\[
  \boxed{
  \mathbb H_Tf^{*,0}
  \simeq
  f^*\mathbb H_S,
  \qquad
  f^{*,0}\Disc_S
  \simeq
  \Disc_Tf^*.
  }
\]
The direct image preserves motivic-local objects and therefore
\[
  f_*^0\Disc_T\simeq\Disc_Sf_*.
\]
\end{theorem}

\begin{proof}
On stable presheaves, left Kan extension along \(u_f^{\op}\) is left adjoint to
restriction along \(u_f^{\op}\).  It carries suspension representables by
\[
  \Sigma^\infty_+h_X
  \longmapsto
  \Sigma^\infty_+h_{X_T}.
\]
Base change preserves spectral Nisnevich distinguished squares and the empty
object, so the left adjoint preserves stable Nisnevich equivalences.  The right
adjoint preserves Nisnevich-local objects because its value on a distinguished
square over \(S\) is the value of the original sheaf on its base change over
\(T\).  Thus the adjunction descends to
\[
  f^*_{\Nis}:\mathcal N_S^{\Sp}
  \rightleftarrows
  \mathcal N_T^{\Sp}:f_{*,\Nis}.
\]
Since base change preserves fibre products and terminal objects,
\(f^*_{\Nis}\) is strong symmetric monoidal.  Exactness follows from stability.
It carries the defining cofiber of \(\mathbb T_S^{\mathrm{raw}}\) to that of
\(\mathbb T_T^{\mathrm{raw}}\), hence
\[
  f^*_{\Nis}\mathbb T_S^{\mathrm{raw}}
  \simeq
  \mathbb T_T^{\mathrm{raw}}.
\]
Consequently the adjunction lifts to the raw symmetric Thom prespectrum
categories.

Strong monoidality identifies the pullback of every stabilization morphism with
the corresponding stabilization morphism over \(T\):
\[
  f^*(\lambda_{m,G})
  \simeq
  \lambda_{m,f^*G}.
\]
Therefore the left adjoint descends through \(L_{\mathrm{Th}}\).  If \(E\) is
Thom-stable over \(T\), adjunction identifies mapping a stabilization generator
into \(f_*E\) with mapping its pullback into \(E\), so the right adjoint also
preserves Thom-stable objects.  This yields the displayed adjunction on
\(\mathcal R^0\).

The construction on smooth-site functors is strictly compatible with
composition up to the canonical Kan-extension equivalences; their coherences
survive stable Nisnevich sheafification, passage to symmetric prespectra and
Thom stabilization.  This gives the composition statements.

The same representable calculation carries every affine-line generator over
\(S\) to the corresponding affine-line generator over \(T\).  Thus
\(f^{*,0}\) descends through the remaining \(\mathbb A^1\)-localization.  On
suspension representables the descended functor is exactly the stable motivic
inverse image of Chapter~6, and compact generation identifies the two functors.
This gives
\(\mathbb H_Tf^{*,0}\simeq f^*\mathbb H_S\).  Since
\(f^{*,0}\) carries local objects to local objects, the two local objects
\(f^{*,0}\Disc_SM\) and \(\Disc_Tf^*M\) have the same homotopification; full
faithfulness of \(\Disc_T\) gives their canonical equivalence.

Finally the pointwise stable direct image preserves affine-line invariance: for
an affine-line-local prespectrum \(E\) over \(T\), evaluating \(f_*^0E\) on
\(X\times_S\mathbb A^1_S\to X\) is evaluating \(E\) on its base change
\(X_T\times_T\mathbb A^1_T\to X_T\), which is an equivalence.  Thus
\(f_*^0\) preserves local objects.  Taking the restricted right adjoint gives
\(f_*^0\Disc_T\simeq\Disc_Sf_*\).
\end{proof}

\section{Smooth direct image before affine-line localization}
\label{sec:bn-diff-thom-stable-smooth-direct-image}

Let \(p:X\to Y\) be smooth and of finite presentation.  Composition with
\(p\) defines
\[
  c_p:\operatorname{Sm}^{\fp}_{/X}
  \longrightarrow
  \operatorname{Sm}^{\fp}_{/Y}.
\]

\begin{proposition}[Stable Nisnevich smooth direct image]
\label{prop:bn-diff-stable-nis-smooth-direct-image}
There is a colimit-preserving exact functor
\[
  p_{\sharp,\Nis}:\mathcal N_X^{\Sp}\longrightarrow\mathcal N_Y^{\Sp}
\]
left adjoint to \(p^*_{\Nis}\).  On suspension representables,
\[
  p_{\sharp,\Nis}\Sigma^\infty_+h_V
  \simeq
  \Sigma^\infty_+h_V,
\]
where \(V\) on the right is regarded as smooth over \(Y\) by composition.
For \(E\in\mathcal N_X^{\Sp}\) and \(F\in\mathcal N_Y^{\Sp}\) there is a canonical
projection equivalence
\[
  \boxed{
  p_{\sharp,\Nis}(E\otimes p^*_{\Nis}F)
  \simeq
  p_{\sharp,\Nis}(E)\otimes F.
  }
\]
In particular,
\[
  \boxed{
  p_{\sharp,\Nis}
  (\mathbb T_X^{\mathrm{raw}}\otimes E)
  \simeq
  \mathbb T_Y^{\mathrm{raw}}\otimes p_{\sharp,\Nis}E.
  }
\]
All equivalences are coherent under composition of smooth morphisms.
\end{proposition}

\begin{proof}
On stable presheaves, left Kan extension along \(c_p^{\op}\) sends a
representable over \(X\) to the same spectral scheme regarded over \(Y\).
A spectral Nisnevich distinguished square over \(X\), after composition with
\(p\), is the same distinguished square over \(Y\); its open immersion,
\'etale morphism, reduced closed complement and complementary equivalence are
unchanged.  The empty object is unchanged as well.  Hence the left Kan
extension preserves the stable Nisnevich generators and descends to
\(p_{\sharp,\Nis}\).  The smooth-site adjunction with base change descends at
the same time, giving \(p_{\sharp,\Nis}\dashv p^*_{\Nis}\).

For suspension representables \(h_V\) over \(X\) and \(h_W\) over \(Y\), both
sides of the projection morphism identify with
\[
  \Sigma^\infty_+h_{V\times_YW}.
\]
The two-variable functors involved preserve colimits separately and suspension
representables generate, so the projection morphism is an equivalence for all
\(E,F\).  Taking \(F=\mathbb T_Y^{\mathrm{raw}}\) and using
\(p^*\mathbb T_Y^{\mathrm{raw}}\simeq\mathbb T_X^{\mathrm{raw}}\) gives raw
Thom-linearity.  The coherence is inherited from Kan extension, the projection
formula and composition on smooth sites.
\end{proof}

\begin{theorem}[Thom-stable smooth direct image]
\label{thm:bn-diff-thom-stable-smooth-direct-image}
The preceding adjunction descends canonically to
\[
  \boxed{
  p_\sharp^0:
  \mathcal R_X^0
  \rightleftarrows
  \mathcal R_Y^0:
  p^{*,0}.
  }
\]
The functor \(p_\sharp^0\) preserves all small colimits and is exact.  For
smooth morphisms
\[
  X\xrightarrow{p}Y\xrightarrow{q}Z
\]
there are coherent equivalences
\[
  (qp)_\sharp^0\simeq q_\sharp^0p_\sharp^0.
\]
For every Cartesian square
\[
\begin{tikzcd}[column sep=large,row sep=large]
X' \ar[r,"g'"] \ar[d,"p'"'] & X \ar[d,"p"]\\
Y' \ar[r,"g"'] & Y
\end{tikzcd}
\]
with \(p\) smooth of finite presentation, the canonical smooth
Beck--Chevalley mate
\[
  \boxed{
  \operatorname{Ex}_{\sharp,p,g}^{0}:
  p_\sharp^{\prime,0}(g')^{*,0}
  \xrightarrow{\simeq}
  g^{*,0}p_\sharp^0
  }
\]
is an equivalence.  We write
\[
  \boxed{
  \beta_{\sharp,p,g}^{0}
  :=
  \left(\operatorname{Ex}_{\sharp,p,g}^{0}\right)^{-1}:
  g^{*,0}p_\sharp^0
  \xrightarrow{\simeq}
  p_\sharp^{\prime,0}(g')^{*,0}
  }
\]
for its inverse, which is the orientation used below when constructing the
exceptional-style smooth lift.
There is also a coherent projection formula
\[
  \boxed{
  p_\sharp^0(E\otimes p^{*,0}F)
  \xrightarrow{\simeq}
  p_\sharp^0E\otimes F.
  }
\]
After affine-line homotopification,
\[
  \boxed{
  \mathbb H_Yp_\sharp^0
  \simeq
  p_\sharp\mathbb H_X,
  }
\]
where the right-hand functor is the stable smooth direct image of Chapter~6.
\end{theorem}

\begin{proof}
Raw Thom-linearity in
Proposition~\ref{prop:bn-diff-stable-nis-smooth-direct-image} supplies a functor
on symmetric raw Thom prespectra by applying \(p_{\sharp,\Nis}\) levelwise and
transporting the module structure through
\[
  p_{\sharp,\Nis}
  ((\mathbb T_X^{\mathrm{raw}})^{\otimes r}\otimes E_m)
  \simeq
  (\mathbb T_Y^{\mathrm{raw}})^{\otimes r}
  \otimes p_{\sharp,\Nis}E_m.
\]
It carries \(\lambda_{m,G}\) to the corresponding stabilization morphism over
\(Y\), so it descends through \(L_{\mathrm{Th}}\).  The inverse image is already
constructed in Theorem~\ref{thm:bn-diff-thom-stable-base-adjunction}; the
linearity equivalences for the two sides are adjoint mates by the projection
formula, so the adjunction descends with its unit and counit.

Before Thom stabilization, smooth composition and the projection formula are
obtained on suspension representables from the corresponding geometric
identities.  For the Cartesian square above, the canonical smooth
Beck--Chevalley mate has the conventional direction
\[
    p'_{\sharp,\Nis}(g')^*_{\Nis}
    \longrightarrow
    g^*_{\Nis}p_{\sharp,\Nis}.
\]
On a suspension representable attached to $V\to X$, both sides identify
canonically with the suspension representable attached to
$V\times_YY'\to Y'$, and under these identifications the mate is the identity.
Both functors preserve colimits, so the mate is an equivalence on all stable
Nisnevich sheaves.  Its inverse is therefore canonical as well.

The smooth composition, Beck--Chevalley, inverse Beck--Chevalley, and
projection transformations are compatible with raw Thom-linearity and hence
with the stabilization generators.  Passing through \(L_{\mathrm{Th}}\)
gives the displayed equivalences.  The canonical mate satisfies horizontal
and vertical Beck--Chevalley pasting; inverting those coherent diagrams gives
the corresponding pasting laws for $\beta_{\sharp,p,g}^{0}$.

Finally, composition with a smooth morphism carries an affine-line projection
over \(X\) to the same affine-line projection regarded over \(Y\).  Thus
\(p_\sharp^0\) preserves affine-line equivalences and descends through
\(\mathbb H\).  On suspension motives the descended functor agrees with the
stable smooth direct image of Chapter~6; compact generation and preservation of
colimits identify the two functors.
\end{proof}

\section{Nisnevich descent for the Thom-stable coefficient system}
\label{sec:bn-diff-thom-stable-descent}

\begin{lemma}[Symmetric prespectra over coefficient pullbacks]
\label{lem:bn-diff-prespectra-pullback}
Let
\[
    \mathcal C
    \simeq
    \mathcal A\times_{\mathcal D}\mathcal B
\]
be a pullback in
$\operatorname{CAlg}(\operatorname{Pr}^{L}_{\mathrm{st}})$ whose tensor
products and colimits are computed componentwise.  Let
\[
    T=(T_{\mathcal A},T_{\mathcal B},\tau)\in\mathcal C
\]
be an object with its specified compatibility over $\mathcal D$.  Then there
are canonical symmetric monoidal equivalences
\[
    \boxed{
    \SymSeq(\mathcal C)
    \simeq
    \SymSeq(\mathcal A)
    \times_{\SymSeq(\mathcal D)}
    \SymSeq(\mathcal B)
    }
\]
and
\[
\boxed{
\begin{aligned}
    \Mod_{\mathbb S_T^{\mathrm{sym}}}
      (\SymSeq(\mathcal C))
    \simeq{}&
    \Mod_{\mathbb S_{T_{\mathcal A}}^{\mathrm{sym}}}
      (\SymSeq(\mathcal A))
\\
    &\times_{
      \Mod_{\mathbb S_{T_{\mathcal D}}^{\mathrm{sym}}}
      (\SymSeq(\mathcal D))
    }
    \Mod_{\mathbb S_{T_{\mathcal B}}^{\mathrm{sym}}}
      (\SymSeq(\mathcal B)),
\end{aligned}}
\]
where $T_{\mathcal D}$ denotes either compatible image of
$T_{\mathcal A}$ or $T_{\mathcal B}$ in $\mathcal D$.
\end{lemma}

\begin{proof}
For every $m\geq0$, the functor category
\[
    \Fun(B\Sigma_m,-)
\]
preserves limits of $\infty$-categories.  Taking the product over all levels
therefore identifies the underlying symmetric-sequence category of
$\mathcal C$ with the pullback of the corresponding symmetric-sequence
categories.  Additive Day convolution is built from the tensor products,
small colimits, and finite symmetric-group induction.  All of these
constructions are componentwise in the given pullback, so this equivalence is
symmetric monoidal.

Under it, the commutative algebra
\[
    \mathbb S_T^{\mathrm{sym}}(m)=T^{\otimes m}
\]
is exactly the compatible triple of the corresponding sphere commutative
algebras.  A module over this algebra consists of a symmetric sequence
together with the action morphisms
\[
    T^{\otimes r}\otimes E_m\longrightarrow E_{m+r}
\]
and their associative, symmetric, and unital coherence data.  Mapping spaces,
the action morphisms, and every coherence diagram are computed componentwise
in a pullback category.  Thus compatible module structures are exactly
pullbacks of the module structures over $\mathcal A$ and $\mathcal B$ with
their specified agreement over $\mathcal D$.  Relative tensor products are
computed by geometric realizations of the componentwise two-sided bar
constructions, so the resulting equivalence is symmetric monoidal.
\end{proof}

\begin{theorem}[Nisnevich descent before affine-line localization]
\label{thm:bn-diff-thom-stable-descent}
Let \(\{U_i\to S\}_{i\in I}\) be a finite spectral Nisnevich cover and let
\(U_\bullet\to S\) be its \v{C}ech nerve.  Restriction induces symmetric
monoidal equivalences
\[
  \boxed{
  \mathcal N_S^{\Sp}
  \xrightarrow{\simeq}
  \lim_{[n]\in\Delta}\mathcal N_{U_n}^{\Sp},
  }
\]
and
\[
  \boxed{
  \mathcal R_S^0
  \xrightarrow{\simeq}
  \lim_{[n]\in\Delta}\mathcal R_{U_n}^0.
  }
\]
Consequently restriction to a Nisnevich cover is conservative on
\(\mathcal R_S^0\), and the presheaf
\[
  S\longmapsto\Pic(\mathcal R_S^0)
\]
satisfies finite Nisnevich descent.
\end{theorem}

\begin{proof}
We first prove excision for a spectral Nisnevich distinguished square of bases
\[
\begin{tikzcd}[column sep=large,row sep=large]
W \ar[r,"k"] \ar[d,"q"'] & V \ar[d,"p"]\\
U \ar[r,"j"'] & X.
\end{tikzcd}
\]
Restriction gives a strong symmetric monoidal functor
\[
    R_{\Nis}:
    \mathcal N_X^{\Sp}
    \longrightarrow
    \mathcal N_U^{\Sp}
    \times_{\mathcal N_W^{\Sp}}
    \mathcal N_V^{\Sp}.
\]
An object of the pullback is a triple $(A,B,\alpha)$ with
\[
    \alpha:q^*A\xrightarrow{\simeq}k^*B.
\]
For $T\to X$ smooth and of finite presentation put
\[
    T_U:=T\times_XU,
    \qquad
    T_V:=T\times_XV,
    \qquad
    T_W:=T\times_XW
\]
and define
\[
    G(A,B,\alpha)(T)
    :=
    A(T_U)
    \times_{A(T_W)}
    B(T_V),
\]
where the map from $B(T_V)$ to $A(T_W)$ is restriction to $T_W$ followed by
$\alpha$.  Base change carries every Nisnevich distinguished square over $X$
to distinguished squares over $U$, $V$, and $W$.  Hence each of the three
spectrum-valued presheaves occurring in this finite pullback is a Nisnevich
sheaf, and so is $G(A,B,\alpha)$.

If $T\to U$, then $T_U\simeq T$ and $T_V\simeq T_W$, so the gluing
equivalence gives
\[
    G(A,B,\alpha)(T)\simeq A(T).
\]
Likewise restriction to $V$ recovers $B$.  Thus
\[
    R_{\Nis}G\simeq\id.
\]
Conversely, for $E\in\mathcal N_X^{\Sp}$,
\[
    GR_{\Nis}(E)(T)
    \simeq
    E(T_U)\times_{E(T_W)}E(T_V).
\]
The square $T_W,T_V,T_U,T$ is the base change of the original distinguished
square, so Nisnevich excision gives
\[
    E(T)\xrightarrow{\simeq}
    E(T_U)\times_{E(T_W)}E(T_V).
\]
Therefore
\[
    \boxed{
    \mathcal N_X^{\Sp}
    \xrightarrow{\simeq}
    \mathcal N_U^{\Sp}
    \times_{\mathcal N_W^{\Sp}}
    \mathcal N_V^{\Sp}.
    }
\]
The equivalence is symmetric monoidal because every restriction functor is
strong symmetric monoidal and the tensor product in the pullback is computed
componentwise.

The raw Thom coordinates restrict compatibly:
\[
    \mathbb T_X^{\mathrm{raw}}|_U\simeq\mathbb T_U^{\mathrm{raw}},
    \qquad
    \mathbb T_X^{\mathrm{raw}}|_V\simeq\mathbb T_V^{\mathrm{raw}},
    \qquad
    \mathbb T_X^{\mathrm{raw}}|_W\simeq\mathbb T_W^{\mathrm{raw}}.
\]
Apply Lemma~\ref{lem:bn-diff-prespectra-pullback} to the preceding symmetric
monoidal pullback of stable Nisnevich coefficient categories and this
compatible raw Thom coordinate.  It gives a canonical symmetric monoidal
equivalence
\[
    \boxed{
    \mathcal P_X^{\mathrm{pre}}
    \xrightarrow{\simeq}
    \mathcal P_U^{\mathrm{pre}}
    \times_{\mathcal P_W^{\mathrm{pre}}}
    \mathcal P_V^{\mathrm{pre}}.
    }
\]

We next show that this equivalence restricts to the Thom-stable local
subcategories.  Let $r:T\to S$ be any smooth morphism of finite presentation.
By Proposition~\ref{prop:bn-diff-stable-nis-smooth-direct-image}, the stable
Nisnevich adjunction
\[
    r_{\sharp,\Nis}\dashv r^*_{\Nis}
\]
is raw Thom-linear.  It therefore lifts to an adjunction on raw symmetric
Thom prespectra.  Let $E$ be Thom-local over $S$ and let
$\lambda_{m,G}$ be a stabilization generator over $T$.  The lifted adjunction
gives
\[
    \map(\lambda_{m,G},r^*E)
    \simeq
    \map(r_\sharp\lambda_{m,G},E).
\]
Raw Thom-linearity identifies
\[
    r_\sharp\lambda_{m,G}
    \simeq
    \lambda_{m,r_{\sharp,\Nis}G}.
\]
For $G\in\mathcal G_T$, the object $r_{\sharp,\Nis}G$ belongs to
$\mathcal G_S$: on a suspension representable it is the same smooth spectral
scheme regarded over $S$ by composition, and integer suspensions are
preserved.  Hence the last mapping map is an equivalence by Thom-locality of
$E$.  Thus raw smooth pullback preserves Thom-local objects.  This applies to
$j,p,k,q$, since open immersions and spectral \'etale morphisms, and their base
changes, are smooth of finite presentation.

Conversely, let $E\in\mathcal P_X^{\mathrm{pre}}$ and suppose that its
restrictions to $U$ and $V$ are Thom-local.  For every stabilization generator
$\lambda_{m,G}$ over $X$, the map of mapping spectra which tests locality is,
under the pullback description of $\mathcal P_X^{\mathrm{pre}}$, the pullback
of the corresponding testing maps over $U$ and $V$, with their common
restriction over $W$.  Those two maps are equivalences by hypothesis, hence
the global testing map is an equivalence.  Thus $E$ is Thom-local.  We obtain
\[
    \boxed{
    \mathcal R_X^0
    \xrightarrow{\simeq}
    \mathcal R_U^0
    \times_{\mathcal R_W^0}
    \mathcal R_V^0.
    }
\]

For the empty spectral scheme the smooth site contains only the empty object,
and the empty-cover condition forces a stable Nisnevich sheaf to take it to
zero.  Hence $\mathcal N_\varnothing^{\Sp}$ is the terminal stable presentable
category, and the same is true after formation of raw Thom prespectra and Thom
stabilization.

The Nisnevich distinguished squares of spectral bases form the complete and
regular cd-structure of Chapter~6.  To pass from the distinguished-square
calculation to full descent, let $\mathcal D$ be any test object of
$\operatorname{CAlg}(\operatorname{Pr}^{L}_{\mathrm{st}})$ and apply the
representable functor
\[
    \Map
    \left(
      \mathcal D,-
    \right)
\]
to either coefficient system $S\mapsto\mathcal N_S^{\Sp}$ or
$S\mapsto\mathcal R_S^0$.  The preceding calculations show that the resulting
space-valued presheaf takes the empty base to a point and every Nisnevich
distinguished square to a Cartesian square.  The cd-structure criterion of
Chapter~6 therefore makes it a Nisnevich sheaf.  Since this holds for every
$\mathcal D$ and the Yoneda embedding preserves limits, both coefficient
systems satisfy base-Nisnevich descent.  In particular, for the finite cover
of the theorem,
\[
    \mathcal N_S^{\Sp}
    \xrightarrow{\simeq}
    \lim_{[n]\in\Delta}\mathcal N_{U_n}^{\Sp},
    \qquad
    \mathcal R_S^0
    \xrightarrow{\simeq}
    \lim_{[n]\in\Delta}\mathcal R_{U_n}^0.
\]

Conservativity follows immediately from the second limit description.  Finally,
formation of the Picard $\infty$-groupoid preserves limits of symmetric
monoidal $\infty$-categories: invertibility is detected componentwise and the
space of inverse data for an invertible object is contractible.  Applying this
to the displayed descent equivalence gives finite Nisnevich descent for
$S\mapsto\Pic(\mathcal R_S^0)$.
\end{proof}

\section{Differential Thom objects, exact additivity, and virtual twists}
\label{sec:bn-diff-differential-thom}

\begin{definition}[Raw Thom object of a vector bundle]
\label{def:bn-diff-raw-vector-thom}
Let \(E\) be a finite locally free \(\mathcal O_X\)-module.  Write
\(V_X(E)\to X\) for the corresponding spectral vector bundle and identify the
zero section with \(X\subset V_X(E)\).  Define
\[
  \operatorname{th}^{\mathrm{raw}}_X(E)
  :=
  \cofib\left(
    \Sigma^\infty_+h_{V_X(E)\setminus X}
    \longrightarrow
    \Sigma^\infty_+h_{V_X(E)}
  \right)
  \in\mathcal N_X^{\Sp}
\]
and its Thom-stable realization by
\[
  \boxed{
  \operatorname{th}^0_X(E)
  :=
  Q_X\operatorname{th}^{\mathrm{raw}}_X(E)
  \in\mathcal R_X^0.
  }
\]
\end{definition}

\begin{proposition}[Whitney sum and pullback]
\label{prop:bn-diff-differential-thom-additivity}
For finite locally free modules \(E,F\) on \(X\) there is a canonical coherent
equivalence
\[
  \boxed{
  \operatorname{th}^0_X(E\oplus F)
  \simeq
  \operatorname{th}^0_X(E)\otimes
  \operatorname{th}^0_X(F).
  }
\]
For every morphism \(g:Y\to X\),
\[
  \boxed{
  g^{*,0}\operatorname{th}^0_X(E)
  \simeq
  \operatorname{th}^0_Y(g^*E).
  }
\]
For the trivial rank-one bundle,
\[
  \operatorname{th}^0_X(\mathcal O_X)
  \simeq
  \mathbb T_X^0.
\]
\end{proposition}

\begin{proof}
The vector bundle of a direct sum is
\(V(E\oplus F)=V(E)\times_XV(F)\).  The complement of its zero section is the
union
\[
  (V(E)\setminus X)\times_XV(F)
  \;\cup\;
  V(E)\times_X(V(F)\setminus X).
\]
The two opens and their intersection form a Nisnevich, in fact Zariski,
distinguished gluing diagram.  The pushout-product formula for cofibres in the
stable symmetric monoidal sheaf category therefore identifies
\[
  \operatorname{th}^{\mathrm{raw}}_X(E\oplus F)
  \simeq
  \operatorname{th}^{\mathrm{raw}}_X(E)
  \otimes
  \operatorname{th}^{\mathrm{raw}}_X(F).
\]
Applying the strong symmetric monoidal functor \(Q_X\) gives the first formula.
The pullback formula follows because vector bundles, zero sections and their
open complements commute with base change and because
\(g^{*,0}Q_X\simeq Q_Yg^*\).  For \(E=\mathcal O_X\), the defining quotient is
exactly the raw Thom coordinate.
\end{proof}

\begin{theorem}[Invertibility of differential Thom objects]
\label{thm:bn-diff-differential-thom-invertibility}
For every finite locally free module \(E\) on a qcqs spectral scheme \(X\),
\[
  \boxed{
  \operatorname{th}^0_X(E)
  \in
  \Pic(\mathcal R_X^0).
  }
\]
The resulting symmetric monoidal functor
\[
  \operatorname{th}^0_X:
  (\operatorname{Vect}(X)^\simeq,\oplus)
  \longrightarrow
  \Pic(\mathcal R_X^0)
\]
is natural under pullback.
\end{theorem}

\begin{proof}
Choose a finite Zariski cover \(\{U_i\to X\}\) trivializing \(E\).  On
\(U_i\), if \(E\) has rank \(r_i\), repeated Whitney sum gives
\[
  \operatorname{th}^0_{U_i}(E|_{U_i})
  \simeq
  (\mathbb T_{U_i}^0)^{\otimes r_i},
\]
which is invertible by
Theorem~\ref{thm:infinitesimal-raw-thom-stabilization}.  The cover is
Nisnevich, and Theorem~\ref{thm:bn-diff-thom-stable-descent} identifies
\(\Pic(\mathcal R_X^0)\) with the limit of the Picard groupoids on its
\v{C}ech nerve.  The canonical pullback compatibilities of
Proposition~\ref{prop:bn-diff-differential-thom-additivity} therefore descend
the local inverse objects and their evaluation and coevaluation maps to a
global inverse of \(\operatorname{th}^0_X(E)\).

Whitney sum and its associativity, symmetry and unit coherences give the stated
symmetric monoidal functor on the maximal subgroupoid of vector bundles.
Pullback naturality is
Proposition~\ref{prop:bn-diff-differential-thom-additivity}.
\end{proof}

\begin{lemma}[Affine split exactness for vector bundles]
\label{lem:bn-diff-affine-vect-split}
Let \(U=\Spec A\) be affine.  Every exact sequence
\[
  P'\longrightarrow P\longrightarrow P''
\]
in the exact \(\infty\)-category \(\operatorname{Vect}(U)\) splits.  Hence the canonical split-exact comparison
\[
  \boxed{
  \gamma_U:
  \bigl(\operatorname{Vect}(U)^\simeq,\oplus\bigr)^{\mathrm{gp}}
  \xrightarrow{\simeq}
  K(U)
  }
\]
is an equivalence.  It is natural under pullback between affine spectral
schemes.
\end{lemma}

\begin{proof}
Identify finite locally free quasi-coherent modules on \(U\) with finite
locally free \(A\)-modules.  By Chapter~1 these are finite projective, hence
retracts of finite free connective \(A\)-modules; their duals are again finite
projective and connective.  In particular \(P''\) is dualizable, and there is a
canonical equivalence of mapping spectra
\[
  \map_A(P'',P')
  \simeq
  (P'')^\vee\otimes_AP'.
\]
Both factors on the right are connective, so
\[
  \pi_{-1}\map_A(P'',P')=0.
\]
Applying \(\map_A(P'',-)\) to the cofiber sequence gives a fibre sequence
\[
  \map_A(P'',P')
  \longrightarrow
  \map_A(P'',P)
  \longrightarrow
  \map_A(P'',P'').
\]
The obstruction to lifting \(\id_{P''}\) across the last map lies in
\(\pi_{-1}\map_A(P'',P')\), and therefore vanishes.  Choose the resulting
section \(s:P''\to P\).  The induced map
\[
  P'\oplus P''\longrightarrow P
\]
is a map of cofiber sequences which is the identity on the first and third
terms, hence is an equivalence.  Thus the exact structure is split.

Let
\[
    \gamma_U:
    \bigl(\operatorname{Vect}(U)^\simeq,\oplus\bigr)^{\mathrm{gp}}
    \longrightarrow
    K(U)
\]
be the canonical split-exact comparison.  Before group completion it is
induced simplicially by the cumulative-direct-sum construction
\[
    (P_1,\ldots,P_n)
    \longmapsto
    \left(
      0\to P_1\to
      P_1\oplus P_2\to\cdots\to
      P_1\oplus\cdots\oplus P_n
    \right)
\]
in the Waldhausen \(S_\bullet\)-construction.  Since every admissible exact
sequence in \(\operatorname{Vect}(U)\) is split by the first part of the proof,
the split-exact group-completion theorem identifies $\gamma_U$ as an
equivalence.

No choice of splitting enters the definition of $\gamma_U$.  The
cumulative-direct-sum construction is natural under exact additive functors,
and pullback of finite locally free modules is exact and preserves direct
sums.  Hence the equivalences $\gamma_U$ are natural under pullback between
affine spectral schemes.
\end{proof}

\begin{proposition}[Affine differential \(J\)-homomorphism]
\label{prop:bn-diff-affine-j}
For every affine spectral scheme \(U\) there is a canonical natural map of
grouplike \(E_\infty\)-spaces
\[
  \boxed{
  J_{U,\mathrm{aff}}^0:
  K(U)
  \longrightarrow
  \Pic(\mathcal R_U^0)
  }
\]
characterized by
\[
  J_{U,\mathrm{aff}}^0([E])
  \simeq
  \operatorname{th}^0_U(E)
\]
for every finite locally free \(E\) on \(U\).
\end{proposition}

\begin{proof}
By Theorem~\ref{thm:bn-diff-differential-thom-invertibility}, Whitney sum gives
an \(E_\infty\)-map
\[
  \operatorname{Vect}(U)^\simeq
  \longrightarrow
  \Pic(\mathcal R_U^0).
\]
The target is grouplike.  Therefore the universal property of group completion
gives a canonical \(E_\infty\)-map
\[
    \overline{\operatorname{th}}_U^0:
    \bigl(\operatorname{Vect}(U)^\simeq,\oplus\bigr)^{\mathrm{gp}}
    \longrightarrow
    \Pic(\mathcal R_U^0).
\]
Define
\[
    J_{U,\mathrm{aff}}^0
    :=
    \overline{\operatorname{th}}_U^0\circ\gamma_U^{-1},
\]
where $\gamma_U$ is the canonical natural equivalence of
Lemma~\ref{lem:bn-diff-affine-vect-split}.  The inverse of a natural
equivalence in the functor $\infty$-category is itself natural through a
contractible space of choices.  The normalization on vector bundles is built
into the group-completion construction, and pullback naturality follows from
the naturality of $\gamma_U$ and of the Thom functor.
\end{proof}

\begin{lemma}[Affine-basis Picard descent]
\label{lem:bn-diff-affine-picard-descent}
Let \(X\) be qcqs and let \(\operatorname{AffOp}(X)\) be a small affine
Zariski basis of \(X\).  Restriction induces a canonical equivalence
\[
  \boxed{
  \Pic(\mathcal R_X^0)
  \xrightarrow{\simeq}
  \lim_{U\in\operatorname{AffOp}(X)}
  \Pic(\mathcal R_U^0),
  }
\]
where the limit is taken over the restriction diagram.  The same statement
holds with \(\mathcal R^0\) replaced by \(\SH\).
\end{lemma}

\begin{proof}
Choose a finite affine Zariski cover of \(X\).  Because \(X\) is
quasi-separated, every finite intersection of members of the cover is
quasi-compact and admits a finite affine Zariski cover.  Iterating these finite
affine refinements in the \v{C}ech nerve gives a simplicial affine basis
refinement.

Theorem~\ref{thm:bn-diff-thom-stable-descent} says that
\(S\mapsto\Pic(\mathcal R_S^0)\) satisfies descent for every finite
Nisnevich cover, hence for every finite affine Zariski cover.  Applying descent
successively to the chosen cover and to the affine refinements of its finite
intersections identifies \(\Pic(\mathcal R_X^0)\) with the limit of its values
on the affine Zariski basis.  Independence of the chosen finite cover follows
by passing to common finite affine refinements.  This gives the displayed
canonical equivalence.

For \(\SH\), use the Nisnevich descent and Picard descent of Chapter~6 and the
same basis-refinement argument.
\end{proof}

\begin{theorem}[Differential motivic \(J\)-homomorphism]
\label{thm:bn-diff-differential-j}
For every qcqs spectral scheme \(X\) there is a canonical natural map of
grouplike \(E_\infty\)-spaces
\[
  \boxed{
  J_X^0:
  K(X)
  \longrightarrow
  \Pic(\mathcal R_X^0).
  }
\]
It is characterized by
\[
  \boxed{
  J_X^0([E])
  \simeq
  \operatorname{th}^0_X(E)
  }
\]
for finite locally free \(E\), and for every morphism \(g:Y\to X\) between
qcqs spectral schemes in the standing motivic scope there is a canonical
coherent equivalence
\[
  \boxed{
  g^{*,0}J_X^0(\xi)
  \simeq
  J_Y^0(g^*\xi).
  }
\]
After affine-line homotopification, \(J_X^0\) is the motivic
\(J\)-homomorphism of Chapter~6:
\[
  \boxed{
  \mathbb H_XJ_X^0(\xi)
  \simeq
  J_X(\xi).
  }
\]
For \(\xi\in K(X)\), define the virtual differential Thom twist
\[
  \Sigma^{\xi,0}(M)
  :=
  M\otimes J_X^0(\xi).
\]
These are exact autoequivalences with coherent equivalences
\[
  \Sigma^{\xi+\eta,0}
  \simeq
  \Sigma^{\xi,0}\Sigma^{\eta,0},
  \qquad
  \Sigma^{0,0}\simeq\id,
  \qquad
  \Sigma^{-\xi,0}\simeq(\Sigma^{\xi,0})^{-1},
\]
and
\[
  \boxed{
  g^{*,0}\Sigma^{\xi,0}
  \simeq
  \Sigma^{g^*\xi,0}g^{*,0}.
  }
\]
Their motivic shadows are the virtual Thom twists of Chapter~6:
\[
  \boxed{
  \mathbb H_X\Sigma^{\xi,0}
  \simeq
  \Sigma^\xi\mathbb H_X.
  }
\]
\end{theorem}

\begin{proof}
For every affine open \(U\subseteq X\), restriction in vector-bundle
\(K\)-theory gives
\[
  K(X)\longrightarrow K(U).
\]
The affine maps of Proposition~\ref{prop:bn-diff-affine-j} are compatible under
further affine restriction.  Hence there is a canonical composite
\[
\begin{aligned}
  K(X)
  &\longrightarrow
  \lim_{U\in\operatorname{AffOp}(X)}K(U)\\
  &\xrightarrow{\ \lim J_{U,\mathrm{aff}}^0\ }
  \lim_{U\in\operatorname{AffOp}(X)}\Pic(\mathcal R_U^0)\\
  &\xrightarrow{\simeq}
  \Pic(\mathcal R_X^0),
\end{aligned}
\]
where the final equivalence is
Lemma~\ref{lem:bn-diff-affine-picard-descent}.  Define this composite to be
\(J_X^0\).  Notice that no descent equivalence for \(K(X)\) is used or
asserted; only the canonical restriction map into the affine limit is needed.

If \(E\) is finite locally free, then for every affine open \(U\subseteq X\),
\[
  J_{U,\mathrm{aff}}^0([E|_U])
  \simeq
  \operatorname{th}^0_U(E|_U)
  \simeq
  \operatorname{th}^0_X(E)|_U.
\]
Affine-basis Picard descent therefore identifies
\(J_X^0([E])\) with \(\operatorname{th}^0_X(E)\).

Let \(g:Y\to X\) be such a morphism.  To compare \(g^{*,0}J_X^0(\xi)\) with
\(J_Y^0(g^*\xi)\), restrict to an affine open \(V\subseteq Y\).  Choose a
finite affine open cover of \(X\); the inverse images cover \(V\), and since
\(V\) is quasi-compact they admit a finite affine refinement
\(\{V_a\to V\}\) such that each \(V_a\) maps into an affine open
\(U_a\subseteq X\).  On every \(V_a\), the two objects are identified by the
affine pullback naturality of Proposition~\ref{prop:bn-diff-affine-j}.  Picard
descent on \(V\), followed by affine-basis descent on \(Y\), gives the stated
canonical equivalence.  Identity, composition, and higher pullback coherences
are checked after the same affine refinements and are unique by descent.

Apply the strong symmetric monoidal functor \(\mathbb H_X\).  On an affine
open \(U\), both
\[
  \mathbb H_UJ_{U,\mathrm{aff}}^0,
  \qquad
  J_U:
  K(U)\longrightarrow\Pic(\SH(U))
\]
are maps of grouplike \(E_\infty\)-spaces extending the same symmetric
monoidal Thom functor
\[
  E\longmapsto\operatorname{th}_U(E).
\]
By Lemma~\ref{lem:bn-diff-affine-vect-split} and the universal property of
group completion, they are canonically equivalent.  The affine-basis Picard
descent statement for \(\SH\) now identifies
\(\mathbb H_XJ_X^0\) with the motivic \(J_X\) globally.

The formulas for virtual twists are the functorial consequences of the map of
grouplike \(E_\infty\)-spaces \(J_X^0\).  Tensoring with an invertible object
is an exact autoequivalence.  Pullback naturality gives the pullback formula,
and the motivic comparison of \(J\)-objects gives the final displayed
homotopification equivalence.
\end{proof}

\begin{proposition}[Exact differential Thom additivity and Waldhausen coherence]
\label{prop:bn-diff-exact-thom-additivity}
For every exact sequence
\[
  0\longrightarrow E'
  \longrightarrow E
  \longrightarrow E''
  \longrightarrow0
\]
of finite locally free modules on a qcqs spectral scheme \(X\), there is a
canonical equivalence
\[
  \boxed{
  \operatorname{th}^0_X(E)
  \xrightarrow{\simeq}
  \operatorname{th}^0_X(E')\otimes
  \operatorname{th}^0_X(E'').
  }
\]
It is natural under morphisms of exact sequences and compatible with direct
sums.  More generally, for a filtration
\[
  0=E_0\longrightarrow E_1\longrightarrow\cdots\longrightarrow E_n
\]
and \(Q_{a,b}:=E_b/E_a\), every subdivision
\(a=a_0<a_1<\cdots<a_r=b\) carries a canonical equivalence
\[
  \boxed{
  \operatorname{th}^0_X(Q_{a,b})
  \xrightarrow{\simeq}
  \bigotimes_{\nu=1}^{r}
  \operatorname{th}^0_X(Q_{a_{\nu-1},a_\nu}),
  }
\]
and these equivalences satisfy the full associativity, refinement,
degeneracy, direct-sum, and higher coherence supplied by the Waldhausen
\(S_\bullet\)-construction.  For the standard split exact sequence
\[
    0\longrightarrow E'
    \longrightarrow E'\oplus E''
    \longrightarrow E''
    \longrightarrow0
\]
the displayed additivity equivalence is the Whitney-sum equivalence of
Proposition~\ref{prop:bn-diff-differential-thom-additivity}.  More generally,
if an exact sequence is equipped with a specified splitting
\[
    \sigma:E'\oplus E''\xrightarrow{\simeq}E
\]
compatible with the inclusion of $E'$ and the quotient to $E''$, its
additivity equivalence is the transport of the Whitney-sum equivalence
along $\operatorname{th}^0_X(\sigma)$, naturally in morphisms of split exact
sequences.
\end{proposition}

\begin{proof}
The exact sequence defines its canonical Waldhausen additivity path
\[
  [E]\simeq[E']+[E'']
\]
in \(K(X)\).  Apply the map of grouplike \(E_\infty\)-spaces
\(J_X^0\) of Theorem~\ref{thm:bn-diff-differential-j}.  Using
\(J_X^0([F])\simeq\operatorname{th}^0_X(F)\) gives the displayed Thom
additivity equivalence.

Naturality under morphisms of exact sequences and compatibility with direct
sums are inherited from the functoriality and \(E_\infty\)-structure of the
Waldhausen \(S_\bullet\)-construction.  For a filtration, the simplices of
\(S_\bullet\operatorname{Vect}(X)\) supply the canonical higher additivity
simplices relating the classes of the successive quotients under arbitrary
subdivision.  Applying the single map \(J_X^0\) carries all of these simplices
to the Picard \(\infty\)-groupoid, producing the stated equivalences together
with their higher coherence.

For the standard split exact sequence, restrict to an affine Zariski basis.
There the equivalence $\gamma_U$ of
Lemma~\ref{lem:bn-diff-affine-vect-split} is represented by the
cumulative-direct-sum simplices.  Consequently the Waldhausen additivity path
of the standard split sequence is sent by
$J_{U,\mathrm{aff}}^0$ to the Whitney-sum equivalence.  Affine-basis Picard
descent identifies the resulting global equivalence with the Whitney-sum
equivalence of
Proposition~\ref{prop:bn-diff-differential-thom-additivity}.

If the exact sequence is equipped with a specified splitting
$\sigma:E'\oplus E''\xrightarrow{\simeq}E$ compatible with the inclusion and
quotient, the standard split sequence maps to the given exact sequence by a
morphism of exact sequences.  Naturality of
the Waldhausen additivity path and of $J_X^0$ therefore transports the
Whitney-sum equivalence along $\operatorname{th}^0_X(\sigma)$.  The same
naturality gives compatibility with morphisms of split exact sequences.
\end{proof}

\begin{remark}[The level-zero realization is not asserted fully faithful]
\label{rem:bn-diff-Q-not-fully-faithful}
The functor
\[
  Q_S=L_{\mathrm{Th},S}F_{0,S}:\mathcal N_S^{\Sp}\to\mathcal R_S^0
\]
is the realization functor used by later coefficient constructions.  Unlike
the free level-zero functor into the raw prespectrum presentation category, it
is not asserted to be fully faithful: Thom stabilization can change mapping
objects.  This distinction is essential when intrinsic de Rham coefficients
are realized in the stable differential category.
\end{remark}

\section{The smooth Thom-stable exceptional lift}
\label{sec:bn-diff-smooth-thom-stable-umkehr}

Let
\[
  p:X\longrightarrow Y
\]
be smooth, separated, and of finite presentation.  Write
\[
  T_p:=[\Omega_{X/Y}]\in K(X)
\]
for the tangent class used in Chapter~8.

\begin{definition}[Smooth Thom-stable Umkehr lift]
\label{def:bn-diff-smooth-thom-stable-umkehr}
Define
\[
  \boxed{
  \widetilde p_!
  :=
  p_\sharp^0\Sigma^{-T_p,0}:
  \mathcal R_X^0\longrightarrow\mathcal R_Y^0
  }
\]
and
\[
  \boxed{
  \widetilde p^{\,!}
  :=
  \Sigma^{T_p,0}p^{*,0}:
  \mathcal R_Y^0\longrightarrow\mathcal R_X^0.
  }
\]
\end{definition}

\begin{theorem}[Differential smooth adjunction and motivic shadow]
\label{thm:bn-diff-smooth-thom-stable-umkehr}
There is a canonical adjunction
\[
  \boxed{
  \widetilde p_!
  \dashv
  \widetilde p^{\,!}.
  }
\]
The left adjoint is exact and preserves all small colimits.  There are canonical
coherent equivalences
\[
  \boxed{
  \mathbb H_Y\widetilde p_!
  \simeq
  p_!\mathbb H_X,
  \qquad
  \mathbb H_X\widetilde p^{\,!}
  \simeq
  p^!\mathbb H_Y.
  }
\]
Thus \(\widetilde p_!\) is a canonically constructed Thom-stable lift of the
smooth exceptional direct image of Chapter~8.
\end{theorem}

\begin{proof}
For \(E\in\mathcal R_X^0\) and \(F\in\mathcal R_Y^0\), invertibility of the
Thom twist and \(p_\sharp^0\dashv p^{*,0}\) give
\[
\begin{aligned}
\Map_Y(\widetilde p_!E,F)
&\simeq
\Map_Y(p_\sharp^0\Sigma^{-T_p,0}E,F)\\
&\simeq
\Map_X(\Sigma^{-T_p,0}E,p^{*,0}F)\\
&\simeq
\Map_X(E,\Sigma^{T_p,0}p^{*,0}F),
\end{aligned}
\]
which proves the adjunction.

Using Theorem~\ref{thm:bn-diff-thom-stable-smooth-direct-image} and
Theorem~\ref{thm:bn-diff-differential-j},
\[
\begin{aligned}
\mathbb H_Y\widetilde p_!
&\simeq
\mathbb H_Yp_\sharp^0\Sigma^{-T_p,0}\\
&\simeq
p_\sharp\Sigma^{-T_p}\mathbb H_X\\
&\simeq
p_!\mathbb H_X,
\end{aligned}
\]
where the last equivalence is Corollary~8.10.1.  For the right adjoint, use the
direct calculation
\[
\begin{aligned}
\mathbb H_X\widetilde p^{\,!}
&\simeq
\mathbb H_X\Sigma^{T_p,0}p^{*,0}\\
&\simeq
\Sigma^{T_p}p^*\mathbb H_Y\\
&\simeq
p^!\mathbb H_Y,
\end{aligned}
\]
where the last equivalence is smooth purity from Theorem~8.10.1.  Taking the
right-adjoint mate of the already constructed left comparison instead gives
the complementary local-object equivalence
\[
  \widetilde p^{\,!}\Disc_Y
  \simeq
  \Disc_Xp^!.
\]
\end{proof}

\begin{theorem}[Composition, base change, and projection for the smooth lift]
\label{thm:bn-diff-smooth-thom-stable-coherence}
For smooth separated finite-presentation morphisms
\[
  X\xrightarrow{p}Y\xrightarrow{q}Z
\]
there is a canonical coherent equivalence
\[
  \boxed{
  \widetilde q_!\widetilde p_!
  \xrightarrow{\simeq}
  \widetilde{(qp)}_!.
  }
\]
For every Cartesian square
\[
\begin{tikzcd}[column sep=large,row sep=large]
X' \ar[r,"g'"] \ar[d,"p'"'] & X \ar[d,"p"]\\
Y' \ar[r,"g"'] & Y
\end{tikzcd}
\]
with \(p\) smooth, separated and of finite presentation, there is a canonical
Beck--Chevalley equivalence
\[
  \boxed{
  g^{*,0}\widetilde p_!
  \xrightarrow{\simeq}
  \widetilde p'_!(g')^{*,0}.
  }
\]
For \(L\in\Pic(\mathcal R_Y^0)\),
\[
  \widetilde p_!(E\otimes p^{*,0}L)
  \simeq
  \widetilde p_!E\otimes L.
\]
All transformations satisfy the identity, composition, horizontal and vertical
pasting coherences inherited from smooth base change, cotangent transitivity and
Thom additivity.
\end{theorem}

\begin{proof}
Cotangent transitivity for smooth morphisms gives the canonical path in
\(K(X)\)
\[
  T_{qp}\simeq T_p+p^*T_q.
\]
Applying the natural map \(J_X^0\) of
Theorem~\ref{thm:bn-diff-differential-j}, together with pullback naturality,
gives a canonical coherent equivalence of differential Thom twists
\[
  \Sigma^{-T_{qp},0}
  \simeq
  \Sigma^{-T_p,0}\Sigma^{-p^*T_q,0}.
\]
The higher associativity of this equivalence is the image under \(J^0\) of the
higher cotangent-transitivity coherence, equivalently of the Waldhausen
coherence recorded in
Proposition~\ref{prop:bn-diff-exact-thom-additivity}.

Using the differential projection formula and composition of smooth direct
image,
\[
\begin{aligned}
\widetilde q_!\widetilde p_!E
&=
q_\sharp^0\Sigma^{-T_q,0}
  p_\sharp^0\Sigma^{-T_p,0}E\\
&\simeq
q_\sharp^0p_\sharp^0
  \Sigma^{-p^*T_q,0}\Sigma^{-T_p,0}E\\
&\simeq
(qp)_\sharp^0\Sigma^{-T_{qp},0}E
=
\widetilde{(qp)}_!E.
\end{aligned}
\]
The coherence is inherited from coherent smooth composition and projection,
the coherently associative cotangent-transitivity paths, and the natural
\(E_\infty\)-map \(J^0\).

For base change, the inverse smooth Beck--Chevalley equivalence of
Theorem~\ref{thm:bn-diff-thom-stable-smooth-direct-image} gives
\[
  \beta_{\sharp,p,g}^{0}:
  g^{*,0}p_\sharp^0
  \xrightarrow{\simeq}
  p_\sharp^{\prime,0}(g')^{*,0}.
\]
Cotangent base change gives \(T_{p'}\simeq(g')^*T_p\), and pullback
compatibility of the differential \(J\)-homomorphism gives
\[
  (g')^{*,0}\Sigma^{-T_p,0}
  \simeq
  \Sigma^{-T_{p'},0}(g')^{*,0}.
\]
Pasting these equivalences proves Beck--Chevalley.  The projection formula is
the smooth projection formula combined with additivity of Thom twists.  All
pasting statements follow from the coherences already constructed for those
components.
\end{proof}

\begin{corollary}[Affine-line Beck--Chevalley for the smooth lift]
\label{cor:bn-diff-smooth-affine-line-beck-chevalley}
Let \(q_Y:\mathbb A^1_Y\to Y\) and \(q_X:\mathbb A^1_X\to X\) be the affine-line
projections and let
\[
  p_{\mathbb A^1}:\mathbb A^1_X
  =X\times_Y\mathbb A^1_Y
  \longrightarrow
  \mathbb A^1_Y
\]
be the base change of \(p\).  Then there is a canonical equivalence
\[
  \boxed{
  \beta_p:
  q_Y^{*,0}\widetilde p_!
  \xrightarrow{\simeq}
  \widetilde{p_{\mathbb A^1}}{}_!\,q_X^{*,0}.
  }
\]
It is coherent with the zero and unit section squares, with identity morphisms,
and with composition of smooth separated morphisms.
\end{corollary}

\begin{proof}
Apply Theorem~\ref{thm:bn-diff-smooth-thom-stable-coherence} to the Cartesian
affine-line square.  The section compatibilities and composition coherences are
the corresponding pasting compatibilities of smooth base change and the
pullback naturality of the tangent Thom twist.
\end{proof}

\section*{Part VI: BNV completion, obstruction, and verified Umkehr comparison}
\addcontentsline{toc}{section}{Part VI: BNV completion, obstruction, and verified Umkehr comparison}

\section{The canonical BNV correction of a constructed stable lift}
\label{sec:bn-diff-canonical-bnv-correction}

The purpose of this section is to replace any putative BNV admissibility
condition by a canonical completion construction and then to measure whether
the original geometric lift already agrees with that completion.  We start
with an actual exact functor which has already been constructed and whose
motivic shadow has already been identified.  No class of ``admissible''
morphisms or extra cycle datum is introduced.

Let
\[
  T:\mathcal R_X^0\longrightarrow\mathcal R_Y^0
\]
be an exact functor, let
\[
  F:\SH(X)\longrightarrow\SH(Y)
\]
be exact, and suppose a natural equivalence
\[
  \vartheta_T:
  \mathbb H_YT
  \xrightarrow{\simeq}
  F\mathbb H_X
\]
has been constructed.  For the applications below,
\(T=\widetilde p_!\), \(F=p_!\), and
\(\vartheta_T\) is Theorem~\ref{thm:bn-diff-smooth-thom-stable-umkehr}.

\begin{construction}[Canonical cycle realization attached to a lift]
\label{con:bn-diff-canonical-cycle-realization}
Define
\[
  P_T:=\mathcal Z_YT:
  \mathcal R_X^0\longrightarrow\mathcal R_Y^{\mathrm{BNV}}.
\]
There is a canonical natural transformation
\[
  \boxed{
  \psi_T:
  F\mathbb H_X
  \longrightarrow
  \mathbb H_YP_T
  }
\]
given by
\[
  F\mathbb H_X
  \xrightarrow{\vartheta_T^{-1}}
  \mathbb H_YT
  \xrightarrow{\mathbb H_YR^Y_T}
  \mathbb H_Y\mathcal Z_YT.
\]
For \(E\in\mathcal R_X^0\), define the \emph{canonical BNV correction} of
\(T\) by
\[
\boxed{
\begin{aligned}
  \widehat T(E)
  &:={}
  P_T(\mathcal Z_XE)
  \\[-1mm]
  &\quad\times_{
    \Disc_Y\mathbb H_YP_T(\mathcal Z_XE)
  }
  \Disc_YF\mathbb H_XE,
\end{aligned}}
\]
where the map from the second factor is induced by
\[
\begin{aligned}
  F\mathbb H_XE
  &\xrightarrow{F(\phi_E)}
  F\mathbb H_X\mathcal Z_XE\\
  &\xrightarrow{\psi_T(\mathcal Z_XE)}
  \mathbb H_YP_T(\mathcal Z_XE).
\end{aligned}
\]
\end{construction}

\begin{proposition}[Sectors of the canonical correction]
\label{prop:bn-diff-canonical-correction-sectors}
The functor \(\widehat T\) is exact.  There are canonical equivalences
\[
  \boxed{
  \mathbb H_Y\widehat T
  \simeq
  F\mathbb H_X,
  \qquad
  \mathcal Z_Y\widehat T
  \simeq
  \mathcal Z_YT\mathcal Z_X.
  }
\]
Under these identifications, its characteristic morphism is
\[
  F\mathbb H_XE
  \xrightarrow{F(\phi_E)}
  F\mathbb H_X\mathcal Z_XE
  \xrightarrow{\psi_T(\mathcal Z_XE)}
  \mathbb H_Y\mathcal Z_YT\mathcal Z_XE.
\]
\end{proposition}

\begin{proof}
The object \(P_T(\mathcal Z_XE)=\mathcal Z_YT\mathcal Z_XE\) is BNV-pure by
construction.  The formula for \(\widehat T(E)\) is exactly the inverse
construction in Theorem~\ref{thm:infinitesimal-bnv-classification} applied to
the displayed characteristic triple.  That theorem gives the two sector
identifications and the characteristic formula.  Exactness follows because
\(T\), \(F\), \(\mathbb H\), \(\Disc\), and \(\mathcal Z\) are exact and finite
pullbacks are finite colimits in a stable category.
\end{proof}

\begin{proposition}[The canonical correction preserves the motivic-local sector]
\label{prop:bn-diff-canonical-correction-local}
For every constructed exact lift \(T\) with motivic shadow
\(\vartheta_T:\mathbb H_YT\simeq F\mathbb H_X\), the canonical correction
\[
  \widehat T:\mathcal R_X^0\longrightarrow\mathcal R_Y^0
\]
carries motivic-local objects to motivic-local objects.
\end{proposition}

\begin{proof}
Let \(L\in\mathcal R_X^0\) be motivic local.  Since \(L\) belongs to the
essential image of \(\Disc_X\), the coreflection counit
\[
  \mathcal S_XL\longrightarrow L
\]
is an equivalence.  Its cofiber is therefore zero:
\[
  \mathcal Z_XL\simeq0.
\]
Proposition~\ref{prop:bn-diff-canonical-correction-sectors} gives
\[
\begin{aligned}
  \mathcal Z_Y\widehat T(L)
  &\simeq
  \mathcal Z_YT\mathcal Z_XL\\
  &\simeq
  \mathcal Z_YT(0)
  \simeq0,
\end{aligned}
\]
where the last equivalence uses exactness of \(T\).

For any \(M\in\mathcal R_Y^0\), the cycle cofiber sequence
\[
  \mathcal S_YM\longrightarrow M\longrightarrow\mathcal Z_YM
\]
shows that \(\mathcal Z_YM\simeq0\) precisely when the coreflection counit is
an equivalence, equivalently when \(M\) lies in the essential image of
\(\Disc_Y\).  Hence \(\widehat T(L)\) is motivic local.
\end{proof}

\begin{definition}[BNV cycle defect]
\label{def:bn-diff-cycle-defect}
Define the natural transformation
\[
  \boxed{
  \omega_T:
  \mathcal Z_YT
  \longrightarrow
  \mathcal Z_YT\mathcal Z_X
  }
\]
by
\[
  \omega_T(E):=\mathcal Z_YT(R_E^X).
\]
We call \(\omega_T\) the \emph{BNV cycle defect} of the constructed lift
\(T\).
\end{definition}


\begin{definition}[BNV obstruction object]
\label{def:bn-diff-bnv-obstruction-object}
Define the exact functor-valued obstruction object of $T$ by
\[
    \boxed{
    \operatorname{Obs}_{\mathrm{BNV}}(T)
    :=
    \cofib(\omega_T).
    }
\]
Thus $\operatorname{Obs}_{\mathrm{BNV}}(T)$ is an exact functor
$\mathcal R_X^0\to\mathcal R_Y^{\mathrm{BNV}}$.
\end{definition}

\begin{proposition}[Exact form of the BNV obstruction]
\label{prop:bn-diff-bnv-obstruction-object}
There are canonical exact triangles of functors
\[
    \mathcal Z_YT\mathcal S_X
    \longrightarrow
    \mathcal Z_YT
    \xrightarrow{\omega_T}
    \mathcal Z_YT\mathcal Z_X
\]
and
\[
    \mathcal Z_YT
    \xrightarrow{\omega_T}
    \mathcal Z_YT\mathcal Z_X
    \longrightarrow
    \operatorname{Obs}_{\mathrm{BNV}}(T).
\]
Consequently
\[
    \boxed{
    \operatorname{Obs}_{\mathrm{BNV}}(T)
    \simeq
    \Sigma\,\mathcal Z_YT\mathcal S_X.
    }
\]
In particular,
\[
    \operatorname{Obs}_{\mathrm{BNV}}(T)\simeq0
    \quad\Longleftrightarrow\quad
    \omega_T\text{ is an equivalence}.
\]
\end{proposition}

\begin{proof}
Apply the exact functor $\mathcal Z_YT$ to the cycle cofiber sequence over $X$,
\[
    \mathcal S_X
    \longrightarrow
    \id
    \xrightarrow{R^X}
    \mathcal Z_X.
\]
The resulting second morphism is exactly
$\mathcal Z_YT(R^X)=\omega_T$, giving the first exact triangle.  The second is
the defining cofiber triangle of $\operatorname{Obs}_{\mathrm{BNV}}(T)$.
In a stable category the cofiber of the second morphism in an exact triangle is
the suspension of the first term, which gives the displayed equivalence and
the vanishing criterion.
\end{proof}

\begin{theorem}[Canonical comparison with the BNV correction]
\label{thm:bn-diff-canonical-correction-comparison}
There is a canonical natural transformation
\[
  \boxed{
  \kappa_T:T\longrightarrow\widehat T
  }
\]
whose underlying motivic component is the identity under
\(\vartheta_T\) and whose BNV-cycle component is \(\omega_T\).  Consequently
\[
  \boxed{
  \kappa_T\text{ is an equivalence}
  \quad\Longleftrightarrow\quad
  \omega_T\text{ is an equivalence}.
  }
\]
\end{theorem}

\begin{proof}
Under the BNV classification, \(T(E)\) has characteristic triple
\[
  \bigl(
    F\mathbb H_XE,
    \mathcal Z_YT(E),
    \mathbb H_Y(R^Y_{T(E)})\circ\vartheta_{T,E}^{-1}
  \bigr).
\]
The correction \(\widehat T(E)\) has the triple of
Proposition~\ref{prop:bn-diff-canonical-correction-sectors}.  Take the identity
on the underlying component and
\[
  \omega_T(E)=\mathcal Z_YT(R_E^X)
\]
on the cycle component.  Naturality of the target cycle map gives
\[
  \mathcal Z_YT(R_E^X)\circ R^Y_{T(E)}
  \simeq
  R^Y_{T(\mathcal Z_XE)}\circ T(R_E^X).
\]
Apply \(\mathbb H_Y\) and use naturality of \(\vartheta_T\).  The resulting
homotopy is exactly the compatibility square between the two characteristic
maps.  Hence the pair determines a morphism of characteristic triples, and
Theorem~\ref{thm:infinitesimal-bnv-classification} gives \(\kappa_T\).

A morphism in the characteristic pullback category is an equivalence exactly
when its underlying and cycle components are equivalences.  The first component
is the identity, so \(\kappa_T\) is an equivalence exactly when \(\omega_T\) is.
\end{proof}

\section{Recognition of the BNV obstruction}
\label{sec:bn-diff-bnv-obstruction-recognition}

\begin{theorem}[Equivalent forms of the BNV cycle obstruction]
\label{thm:bn-diff-bnv-obstruction}
In the situation above, the following conditions are equivalent.
\begin{enumerate}
\item The cycle defect
\[
  \omega_T:
  \mathcal Z_YT\longrightarrow\mathcal Z_YT\mathcal Z_X
\]
is an equivalence.
\item The BNV obstruction object vanishes:
\[
  \operatorname{Obs}_{\mathrm{BNV}}(T)\simeq0.
\]
\item The functor \(T\) carries every motivic-local object of
\(\mathcal R_X^0\) to a motivic-local object of \(\mathcal R_Y^0\).
\item The canonical localization morphism
\[
  \boxed{
  \ell_T:
  T\Disc_X
  \longrightarrow
  \Disc_YF
  }
\]
obtained from the localization unit and \(\vartheta_T\) is an equivalence.
\item The canonical comparison
\[
  \kappa_T:T\longrightarrow\widehat T
\]
is an equivalence.
\end{enumerate}
If these conditions hold and \(T\) preserves all small colimits, then
\(\widehat T\simeq T\) preserves all small colimits and admits a right adjoint.
\end{theorem}

\begin{proof}
The equivalence of (1) and (2), together with
\[
  \operatorname{Obs}_{\mathrm{BNV}}(T)
  \simeq
  \Sigma\,\mathcal Z_YT\mathcal S_X,
\]
is Proposition~\ref{prop:bn-diff-bnv-obstruction-object}.  Hence (1) holds
exactly when
\[
  \mathcal Z_YT\mathcal S_XE\simeq0
\]
for every \(E\).

An object \(M\in\mathcal R_Y^0\) is motivic local if and only if
\(\mathcal Z_YM\simeq0\).  Indeed, the defining cofiber sequence
\[
  \mathcal S_YM\longrightarrow M\longrightarrow\mathcal Z_YM
\]
shows that \(\mathcal Z_YM\simeq0\) exactly when the coreflection counit is an
equivalence, which is exactly membership in the essential image of
\(\Disc_Y\).  The objects \(\mathcal S_XE\) range through all motivic-local
objects: for local \(L\), the coreflection counit gives
\(\mathcal S_XL\simeq L\).  Thus (1) and (3) are equivalent.

For \(M\in\SH(X)\), define \(\ell_T(M)\) as the composite localization unit
\[
  T\Disc_XM
  \longrightarrow
  \Disc_Y\mathbb H_YT\Disc_XM
  \xrightarrow{\Disc_Y\vartheta_T}
  \Disc_YFM,
\]
using \(\mathbb H_X\Disc_X\simeq\id\).  The target is local and the displayed
map is precisely the localization unit of \(T\Disc_XM\) after the canonical
identification of its localization with \(\Disc_YFM\).  Therefore
\(\ell_T(M)\) is an equivalence exactly when \(T\Disc_XM\) is local.  This is
condition (3), so (3) and (4) are equivalent.

The equivalence of (1) and (5) is
Theorem~\ref{thm:bn-diff-canonical-correction-comparison}.  The final statement
follows because a colimit-preserving functor between presentable categories
admits a right adjoint.
\end{proof}

\begin{corollary}[BNV idempotence of the canonical correction]
\label{cor:bn-diff-canonical-correction-idempotent}
For every constructed exact lift \(T\) with motivic shadow
\(\mathbb H_YT\simeq F\mathbb H_X\), the canonical correction has vanishing
BNV obstruction:
\[
  \boxed{
  \operatorname{Obs}_{\mathrm{BNV}}(\widehat T)\simeq0.
  }
\]
Equivalently, there are canonical equivalences
\[
  \omega_{\widehat T}:
  \mathcal Z_Y\widehat T
  \xrightarrow{\simeq}
  \mathcal Z_Y\widehat T\mathcal Z_X,
\]
\[
  \widehat T\Disc_X
  \xrightarrow{\simeq}
  \Disc_YF,
\]
and
\[
  \boxed{
  \kappa_{\widehat T}:
  \widehat T
  \xrightarrow{\simeq}
  \widehat{\widehat T}.
  }
\]
Thus the canonical BNV correction is idempotent up to its canonical
equivalence.
\end{corollary}

\begin{proof}
By Proposition~\ref{prop:bn-diff-canonical-correction-local},
\(\widehat T\) preserves motivic-local objects.  Its motivic shadow is the
canonical equivalence
\[
  \mathbb H_Y\widehat T\simeq F\mathbb H_X
\]
of Proposition~\ref{prop:bn-diff-canonical-correction-sectors}.  Apply
Theorem~\ref{thm:bn-diff-bnv-obstruction} to \(\widehat T\).  Conditions
(1)--(5) of that theorem all hold, giving the displayed equivalences.
\end{proof}

\begin{remark}[Scope of the completion]
The idempotence statement is an exact BNV statement.  No preservation of
arbitrary colimits by \(\widehat T\), and hence no right adjoint for
\(\widehat T\), is asserted without further input.  When
\(\kappa_T:T\to\widehat T\) is an equivalence and \(T\) preserves all small
colimits, these properties follow from
Theorem~\ref{thm:bn-diff-bnv-obstruction}.
\end{remark}

\begin{remark}[What is and is not proved for a general smooth morphism]
\label{rem:bn-diff-general-smooth-obstruction}
For the constructed smooth lift
\[
  T=\widetilde p_!,
  \qquad
  F=p_!,
\]
all data entering Theorem~\ref{thm:bn-diff-bnv-obstruction} have been
constructed.  In particular, the canonical BNV completion
\[
  \widehat{\widetilde p_!}:
  \mathcal R_X^0\longrightarrow\mathcal R_Y^0
\]
exists, is exact, preserves motivic-local objects, and has motivic shadow
\[
  \mathbb H_Y\widehat{\widetilde p_!}
  \simeq
  p_!\mathbb H_X.
\]
The chapter does not assert that the canonical comparison
\[
  \kappa_p:
  \widetilde p_!
  \longrightarrow
  \widehat{\widetilde p_!}
\]
is an equivalence for every smooth separated finite-presentation morphism.
Equivalently, it does not assert that \(\omega_p\) is always invertible.  The
canonically attached functor
\[
  \operatorname{Obs}_{\mathrm{BNV}}(p)
  :=
  \cofib\left(
    \mathcal Z_Y\widetilde p_!
    \xrightarrow{\omega_p}
    \mathcal Z_Y\widetilde p_!\mathcal Z_X
  \right)
  \simeq
  \Sigma\,\mathcal Z_Y\widetilde p_!\mathcal S_X
\]
therefore measures the failure of the original geometric lift to agree with
its canonical BNV completion.  Its vanishing is equivalent to
\(\kappa_p\) being an equivalence.  Thus the obstruction is not an obstruction
to the existence of a BNV completion and no unspecified BNV-cycle pushforward
or ``admissibility'' condition is introduced.
\end{remark}

\begin{theorem}[Mapping space out of the canonical BNV correction]
\label{thm:bn-diff-correction-mapping-space}
Let \(E\in\mathcal R_X^0\) and \(G\in\mathcal R_Y^0\).  Put
\[
  U_G:=\mathbb H_YG,
  \qquad
  P_G:=\mathcal Z_YG,
  \qquad
  \phi_G:U_G\to\mathbb H_YP_G.
\]
Then there is a canonical homotopy pullback
\[
\boxed{
\begin{aligned}
\Map_{\mathcal R_Y^0}(\widehat T(E),G)
\simeq{}&
\Map_{\SH(Y)}(F\mathbb H_XE,U_G)
\\
&\times_{
\Map_{\SH(Y)}(F\mathbb H_XE,\mathbb H_YP_G)
}
\Map_{\mathcal R_Y^0}
(\mathcal Z_YT\mathcal Z_XE,P_G),
\end{aligned}}
\]
where the first map to the middle term is postcomposition with \(\phi_G\), and
the second is induced by the characteristic morphism of
\(\widehat T(E)\) followed by \(\mathbb H_Y\) of the cycle map.
\end{theorem}

\begin{proof}
Apply the mapping-space formula in the pullback category
\(\mathcal T_Y\) of Theorem~\ref{thm:infinitesimal-bnv-classification} to the
characteristic triple of
Proposition~\ref{prop:bn-diff-canonical-correction-sectors} and the triple
\((U_G,P_G,\phi_G)\).
\end{proof}

\section{Cycle, deformation, and interval exchange when the obstruction vanishes}
\label{sec:bn-diff-bnv-exchange-from-defect}

This section concerns the original constructed lift \(T\).  The equivalent
conditions of Theorem~\ref{thm:bn-diff-bnv-obstruction} are properties of its
canonical comparison \(T\to\widehat T\), not additional structure.  When they
hold, the original lift already agrees with its BNV completion.  The canonical
correction \(\widehat T\) itself satisfies the same local-preservation
criterion unconditionally by
Proposition~\ref{prop:bn-diff-canonical-correction-local}.

\begin{proposition}[Canonical BNV exchanges from local preservation]
\label{prop:bn-diff-bnv-exchanges-from-local-preservation}
Suppose that \(T\) preserves motivic-local objects, equivalently that the
conditions of Theorem~\ref{thm:bn-diff-bnv-obstruction} hold.  Then there are
canonical natural transformations
\[
  \theta_H:
  T\mathcal H_X
  \xrightarrow{\simeq}
  \mathcal H_YT,
\]
\[
  \theta_Z:
  T\mathcal Z_X
  \longrightarrow
  \mathcal Z_YT,
\]
and a canonical equivalence
\[
  \boxed{
  \theta_A:
  T\mathcal A_X
  \xrightarrow{\simeq}
  \mathcal A_YT.
  }
\]
They satisfy
\[
  \theta_H\circ T\eta_X\simeq\eta_YT,
  \qquad
  \theta_Z\circ TR_X\simeq R_YT,
  \qquad
  a_YT\circ\theta_A\simeq Ta_X.
\]
No additional cycle transport is chosen.
\end{proposition}

\begin{proof}
Since \(T\) preserves local objects,
Theorem~\ref{thm:bn-diff-bnv-obstruction} gives the equivalence
\[
  T\Disc_X\simeq\Disc_YF.
\]
Whiskering with \(\mathbb H_X\) and using \(\vartheta_T\) gives
\[
  T\mathcal H_X
  =T\Disc_X\mathbb H_X
  \simeq
  \Disc_YF\mathbb H_X
  \simeq
  \Disc_Y\mathbb H_YT
  =\mathcal H_YT,
\]
compatibly with localization units.  Taking fibres of the two compatible maps
to the local sector gives \(\theta_A\).

For the cycle map, apply \(T\) to
\[
  \mathcal S_XE\longrightarrow E\longrightarrow\mathcal Z_XE.
\]
The object \(T\mathcal S_XE\) is local.  The target
\(\mathcal Z_YT(E)\) is BNV-pure.  Hence
\[
\begin{aligned}
\map_{\mathcal R_Y^0}
(T\mathcal S_XE,\mathcal Z_YT(E))
&\simeq
\map_{\SH(Y)}
(\mathbb H_YT\mathcal S_XE,
 \mathbb S_Y\mathcal Z_YT(E))\\
&\simeq0.
\end{aligned}
\]
Thus the composite
\[
  T\mathcal S_XE\to T(E)\xrightarrow{R^Y_{T(E)}}\mathcal Z_YT(E)
\]
has a contractibly unique nullhomotopy and therefore factors, through a
contractible space of choices, across the cofiber \(T\mathcal Z_XE\).  This is
\(\theta_Z\), and its compatibility with \(R\) is built into the
factorization.  The remaining compatibilities follow from the fibre and
cofiber constructions.
\end{proof}

\begin{corollary}[Unconditional BNV exchanges for the canonical correction]
\label{cor:bn-diff-canonical-correction-exchanges}
For every constructed exact lift \(T\) with motivic shadow
\(\mathbb H_YT\simeq F\mathbb H_X\), the canonical correction \(\widehat T\)
carries canonical transformations
\[
  \widehat T\mathcal H_X
  \xrightarrow{\simeq}
  \mathcal H_Y\widehat T,
\]
\[
  \widehat T\mathcal Z_X
  \longrightarrow
  \mathcal Z_Y\widehat T,
\]
and
\[
  \widehat T\mathcal A_X
  \xrightarrow{\simeq}
  \mathcal A_Y\widehat T,
\]
with the unit, cycle, and deformation compatibilities of
Proposition~\ref{prop:bn-diff-bnv-exchanges-from-local-preservation}.
\end{corollary}

\begin{proof}
Apply Proposition~\ref{prop:bn-diff-bnv-exchanges-from-local-preservation} to
\(\widehat T\), using
Proposition~\ref{prop:bn-diff-canonical-correction-local}.
\end{proof}

For a general exact functor \(T\), no cylinder exchange is introduced as
additional structure.  In the geometric application below, the affine-line
Beck--Chevalley equivalence of the already constructed smooth lift canonically
produces the required cylinder exchange by mate formation, including its
projection, endpoint, and normalization coherences.

\section{The smooth lift has a constructed normalized cylinder exchange}
\label{sec:bn-diff-smooth-cylinder-exchange}

\begin{theorem}[Normalized cylinder exchange for the smooth Thom-stable lift]
\label{thm:bn-diff-smooth-cylinder-exchange}
% Retain the former definition label so references from later draft chapters
% continue to resolve after replacing the extra datum by this construction.
\label{def:bn-diff-normalized-cylinder-exchange}
For every smooth separated finite-presentation
\(p:X\to Y\), the affine-line Beck--Chevalley equivalence
\(\beta_p\) of
Corollary~\ref{cor:bn-diff-smooth-affine-line-beck-chevalley} canonically
constructs a normalized cylinder exchange
\[
  \boxed{
  \theta_I^p:
  \widetilde p_!\mathbb I_X
  \longrightarrow
  \mathbb I_Y\widetilde p_!.
  }
\]
These exchanges are coherent under composition of smooth separated
finite-presentation morphisms.
\end{theorem}

\begin{proof}
Let
\(p_{\mathbb A^1}:\mathbb A^1_X\to\mathbb A^1_Y\) be the affine-line base
change.  Under
\(q_Y^{*,0}\dashv q_{Y,*}^0\), define \(\theta_I^p\) to be the right mate of
\[
\begin{aligned}
&q_Y^{*,0}\widetilde p_!
 q_{X,*}^0q_X^{*,0}
\\
&\xrightarrow{\ \beta_p\ }
\widetilde{p_{\mathbb A^1}}{}_!
 q_X^{*,0}q_{X,*}^0q_X^{*,0}
\\
&\xrightarrow{\ \widetilde{p_{\mathbb A^1}}{}_!(\epsilon_X)\ }
\widetilde{p_{\mathbb A^1}}{}_!q_X^{*,0}
\\
&\xrightarrow{\ \beta_p^{-1}\ }
q_Y^{*,0}\widetilde p_!,
\end{aligned}
\]
where \(\epsilon_X:q_X^{*,0}q_{X,*}^0\to\id\) is the counit.  The triangle
identities show that precomposition with the cylinder unit is carried to the
cylinder unit over \(Y\).  Coherence of \(\beta_p\) with the zero and unit
section squares gives the two endpoint identities.  These statements respect
the specified retraction homotopies because they are obtained by pasting the
same adjunction units and counits.  Mate formation preserves horizontal and
vertical pasting, so composition coherence follows from that of \(\beta_p\).
\end{proof}

\begin{theorem}[BNV integration naturality for the constructed smooth cylinder exchange]
\label{thm:bn-diff-bnv-umkehr-compat}
Let \(p:X\to Y\) be smooth, separated, and of finite presentation, and let
\[
  \theta_I^p:
  \widetilde p_!\mathbb I_X
  \longrightarrow
  \mathbb I_Y\widetilde p_!
\]
be the cylinder exchange constructed in
Theorem~\ref{thm:bn-diff-smooth-cylinder-exchange}.  If the canonical
comparison with the BNV completion
\[
  \kappa_p:
  \widetilde p_!
  \longrightarrow
  \widehat{\widetilde p_!}
\]
is an equivalence, equivalently if the canonical cycle defect
\[
  \omega_p:
  \mathcal Z_Y\widetilde p_!
  \longrightarrow
  \mathcal Z_Y\widetilde p_!\mathcal Z_X
\]
is an equivalence, define
\[
  \theta_{IZ}^p
  :=
  \mathbb I_Y(\theta_Z)\circ\theta_I^p\mathcal Z_X:
  \widetilde p_!\mathbb I_X\mathcal Z_X
  \longrightarrow
  \mathbb I_Y\mathcal Z_Y\widetilde p_!.
\]
Then the universal BNV integrals satisfy the canonical commutative square
\[
\begin{tikzcd}[column sep=huge,row sep=large]
\widetilde p_!\mathbb I_X\mathcal Z_X
  \ar[r,"\widetilde p_!\int_{\mathbb A^1,X}^{\mathrm{BNV}}"]
  \ar[d,"\theta_{IZ}^p"']
&
\widetilde p_!\mathcal A_X
  \ar[d,"\theta_A","\simeq"']
\\
\mathbb I_Y\mathcal Z_Y\widetilde p_!
  \ar[r,"\int_{\mathbb A^1,Y}^{\mathrm{BNV}}"']
&
\mathcal A_Y\widetilde p_!.
\end{tikzcd}
\]
Thus neither a normalized cylinder exchange nor an independent
``BNV-compatibility path'' is additional input for the smooth Umkehr
construction: the former is the mate of the constructed affine-line
Beck--Chevalley equivalence, and the latter follows from the universal BNV
factorization.
\end{theorem}

\begin{proof}
Put \(\delta_X=1_X^*-0_X^*\) and \(\delta_Y=1_Y^*-0_Y^*\).  The endpoint
compatibilities constructed in
Theorem~\ref{thm:bn-diff-smooth-cylinder-exchange} give
\[
  \delta_Y\widetilde p_!\circ\theta_I^p
  \simeq
  \widetilde p_!\delta_X.
\]
Since \(\omega_p\) is an equivalence,
Proposition~\ref{prop:bn-diff-bnv-exchanges-from-local-preservation}
constructs the cycle and deformation exchanges
\[
  \theta_Z:
  \widetilde p_!\mathcal Z_X
  \longrightarrow
  \mathcal Z_Y\widetilde p_!,
  \qquad
  \theta_A:
  \widetilde p_!\mathcal A_X
  \xrightarrow{\simeq}
  \mathcal A_Y\widetilde p_!.
\]
The normalized BNV factorization of
Lemma~\ref{lem:infinitesimal-bnv-normalized-factorization} is characterized by
factoring \(\delta\) through \(\mathbb I\mathcal Z\), together with the
specified normalized nullhomotopy after localization.  The exchange maps just
constructed identify the localization fibre sequences, while
\(\theta_I^p\), being the mate of the affine-line Beck--Chevalley equivalence,
preserves projection, endpoints, and their specified retractions.
Consequently applying \(\widetilde p_!\) to the normalized factorization over
\(X\) and transporting through \(\theta_{IZ}^p\) gives the same endpoint
difference and the same normalized local nullhomotopy as the factorization
over \(Y\).

The cofiber mapping equivalence of
Lemma~\ref{lem:infinitesimal-bnv-mapping-fiber-cofiber} shows that the two
factorizations agree through a contractible space of choices.  The BNV
integral is the unique compatible lift of this factorization through
\[
  \mathcal A\longrightarrow\id\longrightarrow\mathcal H.
\]
Using the fibre part of the same mapping lemma and the equivalence
\(\theta_A\) proves the displayed square.
\end{proof}

\begin{corollary}[Smooth raw BNV--Umkehr comparison with the canonical completion]
\label{cor:bn-diff-smooth-bnv-umkehr-obstruction}
For a smooth separated finite-presentation morphism \(p:X\to Y\), let
\[
  \kappa_p:
  \widetilde p_!
  \longrightarrow
  \widehat{\widetilde p_!}
\]
be the canonical comparison.  If \(\kappa_p\) is an equivalence, equivalently
if the canonical cycle defect
\[
  \omega_p:
  \mathcal Z_Y\widetilde p_!
  \longrightarrow
  \mathcal Z_Y\widetilde p_!\mathcal Z_X
\]
is an equivalence, then the original geometric lift \(\widetilde p_!\) itself
is a BNV differential refinement of motivic \(p_!\):
\[
  \mathbb H_Y\widetilde p_!\simeq p_!\mathbb H_X,
\]
it admits the canonical cycle and deformation exchanges, and its constructed
cylinder exchange satisfies the BNV integration square of
Theorem~\ref{thm:bn-diff-bnv-umkehr-compat}.
\end{corollary}

\begin{proof}
The motivic comparison is
Theorem~\ref{thm:bn-diff-smooth-thom-stable-umkehr}.  The fact that the defect is an equivalence
is equivalent to preservation of local objects by
Theorem~\ref{thm:bn-diff-bnv-obstruction}.  Apply
Proposition~\ref{prop:bn-diff-bnv-exchanges-from-local-preservation},
Theorem~\ref{thm:bn-diff-smooth-cylinder-exchange}, and
Theorem~\ref{thm:bn-diff-bnv-umkehr-compat}.
\end{proof}

\section{Split finite \texorpdfstring{\'etale}{etale} transfer}
\label{sec:bn-diff-split-finite-etale}

\begin{lemma}[Disjoint-union decomposition of Thom-stable coefficients]
\label{lem:bn-diff-disjoint-union-decomposition}
Let
\[
  X=\coprod_{i=1}^{d}Y_i
\]
be a finite disjoint union of qcqs spectral schemes.  Restriction along the
open-and-closed component inclusions induces canonical symmetric monoidal
equivalences
\[
  \boxed{
  \mathcal N_X^{\Sp}
  \xrightarrow{\simeq}
  \prod_{i=1}^{d}\mathcal N_{Y_i}^{\Sp},
  }
\]
and
\[
  \boxed{
  \mathcal R_X^0
  \xrightarrow{\simeq}
  \prod_{i=1}^{d}\mathcal R_{Y_i}^0.
  }
\]
Under the second equivalence, motivic-local objects and the BNV functors are
componentwise.
\end{lemma}

\begin{proof}
Let \(V\to X\) be smooth and of finite presentation and put
\[
  V_i:=V\times_XY_i.
\]
Then
\[
  V\simeq\coprod_iV_i,
\]
and the component inclusions \(V_i\hookrightarrow V\) form a finite Zariski
cover with pairwise empty intersections.  If \(E\in\mathcal N_X^{\Sp}\), stable
Nisnevich descent and \(E(\varnothing)\simeq0\) therefore give
\[
  E(V)\simeq\prod_iE(V_i).
\]
Consequently restriction
\[
  \mathcal N_X^{\Sp}\longrightarrow\prod_i\mathcal N_{Y_i}^{\Sp}
\]
is fully faithful.  Conversely, for a tuple \((E_i)_i\), define
\[
  E(V):=\prod_iE_i(V_i).
\]
Base change of a Nisnevich distinguished square along \(Y_i\to X\) is a
Nisnevich distinguished square over \(Y_i\); finite products of spectra
preserve the corresponding limit diagrams.  Thus \(E\) is a stable
Nisnevich sheaf, and this construction is inverse to restriction.  Since each
component restriction is strong symmetric monoidal and the family is jointly
conservative, the equivalence is symmetric monoidal.

The raw Thom coordinate restricts to the tuple of raw Thom coordinates on the
components.  Formation of symmetric sequences commutes with finite products,
the sphere commutative algebra is componentwise, and modules over a finite
product of such algebras form the product of the corresponding module
categories.  Hence the same argument gives a symmetric monoidal equivalence of
raw symmetric Thom prespectrum categories.  The adjoint Thom-stability
conditions
\[
  E_m\xrightarrow{\simeq}
  \operatorname{Hom}(\mathbb T^{\mathrm{raw}},E_{m+1})
\]
are componentwise, so the equivalence restricts to the displayed equivalence
on \(\mathcal R^0\).

The affine-line generators are likewise componentwise.  Hence motivic
localness is componentwise.  The functors \(\mathbb H\), \(\Disc\), and
\(\mathbb S\), and therefore \(\mathcal H\), \(\mathcal S\),
\(\mathcal A\), and \(\mathcal Z\), identify with the corresponding
componentwise functors by uniqueness of the adjacent adjoints and by the fibre
and cofiber definitions.
\end{proof}

\begin{theorem}[Split finite \'etale differential transfer]
\label{thm:bn-diff-split-finite-etale}
Let
\[
  p:
  \coprod_{i=1}^{d}Y
  \longrightarrow
  Y
\]
be the fold map, regarded as a finite \'etale spectral morphism.  Then
\[
  T_p=0
\]
and, under the canonical symmetric monoidal equivalence
\[
  \mathcal R_{\coprod_iY}^0
  \simeq
  \prod_{i=1}^{d}\mathcal R_Y^0,
\]
one has
\[
  p^{*,0}(E)=(E,\ldots,E),
  \qquad
  \boxed{
  \widetilde p_!(E_1,\ldots,E_d)
  \simeq
  \bigoplus_{i=1}^{d}E_i.
  }
\]
The functor \(\widetilde p_!\) preserves motivic-local objects, and therefore
its BNV cycle defect is an equivalence
\[
  \boxed{
  \omega_p:
  \mathcal Z_Y\widetilde p_!
  \xrightarrow{\simeq}
  \mathcal Z_Y\widetilde p_!\mathcal Z_X.
  }
\]
Equivalently, the canonical comparison with the BNV completion is an
equivalence:
\[
  \boxed{
  \kappa_p:
  \widetilde p_!
  \xrightarrow{\simeq}
  \widehat{\widetilde p_!}.
  }
\]
Moreover the BNV reconstruction sectors and the affine-line cylinder commute with this
transfer through canonical equivalences:
\[
  \widetilde p_!\mathcal A_X
  \simeq
  \mathcal A_Y\widetilde p_!,
  \qquad
  \widetilde p_!\mathcal Z_X
  \simeq
  \mathcal Z_Y\widetilde p_!,
\]
\[
  \widetilde p_!\mathbb I_X
  \simeq
  \mathbb I_Y\widetilde p_!.
\]
Consequently
\[
\begin{tikzcd}[column sep=huge,row sep=large]
\widetilde p_!\mathbb I_X\mathcal Z_X
  \ar[r,"\widetilde p_!\int_X^{\mathrm{BNV}}"]
  \ar[d,"\simeq"']
&
\widetilde p_!\mathcal A_X
  \ar[d,"\simeq"]
\\
\mathbb I_Y\mathcal Z_Y\widetilde p_!
  \ar[r,"\int_Y^{\mathrm{BNV}}"']
&
\mathcal A_Y\widetilde p_!.
\end{tikzcd}
\]
After homotopification this transfer is exactly the finite \'etale motivic
Umkehr:
\[
  \boxed{
  \mathbb H_Y\widetilde p_!
  \simeq
  p_!\mathbb H_X
  \simeq
  p_\sharp\mathbb H_X.
  }
\]
Since \(p\) is proper, the last functor is also canonically equivalent to
\(p_*\mathbb H_X\).
\end{theorem}

\begin{proof}
The cotangent complex of an \'etale morphism vanishes, so \(T_p=0\) and
\(\widetilde p_!=p_\sharp^0\).  By
Lemma~\ref{lem:bn-diff-disjoint-union-decomposition}, pullback is the diagonal
functor
\[
  \mathcal R_Y^0
  \longrightarrow
  \prod_i\mathcal R_Y^0.
\]
Its left adjoint is finite coproduct, which in the stable category
\(\mathcal R_Y^0\) is finite direct sum.  This proves the formula for
\(\widetilde p_!\).

A motivic-local object over \(X\) is componentwise local by the same lemma.
The local subcategory of \(\mathcal R_Y^0\) is closed under all small colimits
by Lemma~\ref{lem:infinitesimal-compact-stable-bireflection}; in particular a
finite direct sum of local objects is local.  Hence \(\widetilde p_!\)
preserves motivic-local objects.  Theorem~\ref{thm:bn-diff-bnv-obstruction}
therefore identifies its cycle defect with the displayed equivalence, and
Theorem~\ref{thm:bn-diff-canonical-correction-comparison} identifies
\(\kappa_p\) with the displayed equivalence to the canonical BNV completion.

Under the product decomposition all BNV functors are componentwise.  The
transfer is finite direct sum.  Exact functors between stable categories
preserve finite direct sums, while the right adjoints occurring in the BNV
triple preserve finite products, which are the same finite direct sums.
Therefore the deformation and cycle exchanges are equivalences.  The
relative affine-line projection over a finite disjoint union is again the
finite disjoint union of the component projections, so the affine-line
adjunction and \(\mathbb I\) are componentwise as well.  This gives the
cylinder equivalence.  The displayed integration square is then
Theorem~\ref{thm:bn-diff-bnv-umkehr-compat}.

Finally Theorem~\ref{thm:bn-diff-smooth-thom-stable-umkehr} identifies the
homotopification with \(p_!\).  \'Etale purity gives
\(p_!\simeq p_\sharp\), and proper normalization for finite morphisms gives
\(p_!\simeq p_*\) in \(\SH(Y)\).
\end{proof}

\section{Nisnevich-locally split finite \texorpdfstring{\'etale}{etale} transfer}
\label{sec:bn-diff-locally-split-finite-etale}

\begin{lemma}[Finite refinement of base Nisnevich covering sieves]
\label{lem:bn-diff-base-nisnevich-finite-refinement}
Every Nisnevich covering sieve of a qcqs spectral scheme contains a finite
Nisnevich covering family obtained from the two-member distinguished covers
of Chapter~6 by pullback and finite composition.  Consequently, if a finite
\'etale morphism
\[
    p:X\longrightarrow Y
\]
is Nisnevich-locally split, then there is a finite Nisnevich cover
\[
    \{Y_\alpha\longrightarrow Y\}_{\alpha\in A}
\]
such that every base change
\[
    X\times_YY_\alpha\longrightarrow Y_\alpha
\]
is a finite disjoint union of copies of $Y_\alpha$.
\end{lemma}

\begin{proof}
Let $\mathcal J_{\mathrm{fin}}$ be the class of sieves on qcqs spectral
schemes which contain a finite family obtained from identity covers and the
two-member spectral Nisnevich distinguished covers of Chapter~6 by pullback
and finite composition.  Identity families belong to
$\mathcal J_{\mathrm{fin}}$.  Pullback stability of distinguished squares
shows that $\mathcal J_{\mathrm{fin}}$ is stable under pullback.  To verify the
sieve transitivity axiom, let $R$ be a sieve in
$\mathcal J_{\mathrm{fin}}$ on $X$ and let $Q$ be a sieve on $X$ such that
for every morphism $U\to X$ belonging to $R$, the pullback sieve $Q|_U$ lies
in $\mathcal J_{\mathrm{fin}}$.  Choose a finite distinguished-composite
covering family
\[
    \{U_i\to X\}_{i\in I}
\]
contained in $R$.  For each of the finitely many $i$, choose a finite
distinguished-composite covering family
\[
    \{V_{ij}\to U_i\}_{j\in J_i}
\]
contained in $Q|_{U_i}$.  The finite family of composites
\[
    \{V_{ij}\to U_i\to X\}_{i,j}
\]
is contained in $Q$ and is again obtained from distinguished covers by
pullback and finite composition.  Hence $Q$ belongs to
$\mathcal J_{\mathrm{fin}}$.  Thus $\mathcal J_{\mathrm{fin}}$ is a
Grothendieck topology.

By construction it contains the distinguished families which generate the
base Nisnevich topology of Chapter~6.  Hence every base Nisnevich covering
sieve belongs to $\mathcal J_{\mathrm{fin}}$ and therefore contains such a
finite covering family.  Conversely, every family used to define
$\mathcal J_{\mathrm{fin}}$ is a Nisnevich covering family, so this finitary
topology is exactly the base Nisnevich topology.

Now suppose $p$ is Nisnevich-locally split.  The sieve of morphisms
$Y'\to Y$ for which $X\times_YY'\to Y'$ is a finite disjoint union of copies
of $Y'$ is covering and is stable under further pullback.  Choose a finite
Nisnevich covering family contained in this sieve by the first part.  It has
the asserted splitting property.
\end{proof}

\begin{theorem}[Descent of the finite \'etale differential refinement]
\label{thm:bn-diff-locally-split-finite-etale}
Let \(p:X\to Y\) be finite \'etale and Nisnevich-locally split, meaning that
the sieve of morphisms $Y'\to Y$ over which
\[
    X\times_YY'\longrightarrow Y'
\]
is a finite disjoint union of copies of $Y'$ is covering.  Then the canonical
smooth Thom-stable lift
\[
  \widetilde p_!=p_\sharp^0
\]
preserves motivic-local objects.  Equivalently,
\[
  \boxed{
  \omega_p:
  \mathcal Z_Y\widetilde p_!
  \xrightarrow{\simeq}
  \mathcal Z_Y\widetilde p_!\mathcal Z_X.
  }
\]
Equivalently again, the canonical comparison
\[
  \boxed{
  \kappa_p:
  \widetilde p_!
  \xrightarrow{\simeq}
  \widehat{\widetilde p_!}
  }
\]
is an equivalence.  Hence the original geometric lift already agrees with its
canonical BNV completion, is a genuine BNV differential refinement of the
brave-new motivic finite \'etale Umkehr, and the universal BNV interval
integral is natural under \(\widetilde p_!\).
\end{theorem}

\begin{proof}
By Lemma~\ref{lem:bn-diff-base-nisnevich-finite-refinement}, choose a finite
Nisnevich cover
\[
    \{Y_\alpha\to Y\}_{\alpha\in A}
\]
such that every base change
\[
    p_\alpha:
    X\times_YY_\alpha\longrightarrow Y_\alpha
\]
is a finite disjoint union of copies of \(Y_\alpha\).

Let \(L\in\mathcal R_X^0\) be motivic local.  Put
\(g_\alpha:Y_\alpha\to Y\) and let
\(g'_\alpha:X\times_YY_\alpha\to X\) be the base change.  Pullback preserves
motivic-local objects by
Theorem~\ref{thm:bn-diff-thom-stable-base-adjunction}, so
\((g'_\alpha)^{*,0}L\) is local.  Smooth base change for the Thom-stable lift
gives
\[
  g_\alpha^{*,0}\widetilde p_!L
  \simeq
  \widetilde p_{\alpha,!}(g'_\alpha)^{*,0}L.
\]
The right-hand side is local by
Theorem~\ref{thm:bn-diff-split-finite-etale}, since \(p_\alpha\) is split.

Consider the localization unit
\[
  \eta_L:
  \widetilde p_!L
  \longrightarrow
  \Disc_Y\mathbb H_Y\widetilde p_!L.
\]
Using the canonical compatibilities
\[
  g_\alpha^{*,0}\Disc_Y
  \simeq
  \Disc_{Y_\alpha}g_\alpha^*,
  \qquad
  \mathbb H_{Y_\alpha}g_\alpha^{*,0}
  \simeq
  g_\alpha^*\mathbb H_Y,
\]
naturality identifies \(g_\alpha^{*,0}\eta_L\) with the localization unit of
the already local object
\(g_\alpha^{*,0}\widetilde p_!L\).  Hence every
\(g_\alpha^{*,0}\eta_L\) is an equivalence.

By Theorem~\ref{thm:bn-diff-thom-stable-descent}, restriction to the finite
Nisnevich cover \(\{Y_\alpha\to Y\}\) is jointly conservative on
\(\mathcal R_Y^0\).  Therefore \(\eta_L\) itself is an equivalence.  Thus
\(\widetilde p_!L\) is motivic local for every motivic-local \(L\).

Theorem~\ref{thm:bn-diff-bnv-obstruction} now gives the displayed equivalence
for \(\omega_p\), and
Theorem~\ref{thm:bn-diff-canonical-correction-comparison} gives the displayed
equivalence for \(\kappa_p\).  The motivic comparison, constructed normalized
cylinder exchange and BNV integration naturality follow from
Theorem~\ref{thm:bn-diff-smooth-thom-stable-umkehr},
Theorem~\ref{thm:bn-diff-smooth-cylinder-exchange}, and
Theorem~\ref{thm:bn-diff-bnv-umkehr-compat}.
\end{proof}

\section*{Part VII: Separation, realization interface, and final package}
\addcontentsline{toc}{section}{Part VII: Separation, realization interface, and final package}

\section{Separation of the infinitesimal and stable layers}
\label{sec:infinitesimal-layer-separation}

The chapter has produced two distinct adjoint patterns:
\[
  i_!\dashv i^*\dashv i_*\dashv i^!
\]
on the space-valued finite-infinitesimal layer and
\[
  \mathbb H_S\dashv\Disc_S\dashv\mathbb S_S
\]
on the Thom-stable differential and motivic coefficient layer.  Their roles
remain different.

\begin{proposition}[Non-identification of the coreflections and localizations]
\label{prop:infinitesimal-nonidentification}
There is no canonical natural equivalence between
\[
  \CrysCoref=i_*i^!
\]
and
\[
  \mathcal S_S=\Disc_S\mathbb S_S,
\]
because they are endofunctors of different categories.  Likewise the native
de Rham reflector \(L_{\dR}\) and BNV homotopification \(\mathbb H_S\) are
distinct localizations.  More precisely, \(L_{\dR}\) is generated by spectral
square-zero immersions in the native space-valued topos, whereas \(\mathbb H_S\)
is generated by stable affine-line projections in the already Thom-stable
category \(\mathcal R_S^0\).
\end{proposition}

\begin{proof}
The crystalline coreflection acts on
\(\mathcal X_{\Sp}^{\mathrm{inf}}\), while \(\mathcal S_S\) acts on
\(\mathcal R_S^0\).  No equivalence between these source categories is part of
the construction.  The two localization statements follow directly from
Theorem~\ref{thm:infinitesimal-dr-localization} and
Definition~\ref{def:infinitesimal-motivic-local-raw}.  Thom stabilization is
not a generator of \(\mathbb H_S\): it has already been imposed by
\(L_{\mathrm{Th},S}\).
\end{proof}

\section{The stable realization interface for intrinsic de Rham coefficients}
\label{sec:bn-diff-realization-interface}

\begin{proposition}[Canonical realization into the differential stable layer]
\label{prop:bn-diff-stable-realization-interface}
For every base \(S\) in the standing motivic scope, the constructed functor
\[
  \boxed{
  Q_S=L_{\mathrm{Th},S}F_{0,S}:
  \mathcal N_S^{\Sp}
  \longrightarrow
  \mathcal R_S^0
  }
\]
is exact, colimit-preserving and strong symmetric monoidal.  It is natural
under pullback:
\[
  f^{*,0}Q_S\simeq Q_Tf^*_{\Nis}.
\]
For a smooth finite-presentation morphism it is compatible with smooth direct
image in the sense that the canonical Thom-linear comparison is inherited from
Proposition~\ref{prop:bn-diff-stable-nis-smooth-direct-image}.
No full-faithfulness statement for \(Q_S\) is made.
\end{proposition}

\begin{proof}
Exactness, colimit preservation and strong symmetric monoidality are immediate
from the corresponding properties of \(F_{0,S}\) and
\(L_{\mathrm{Th},S}\).  Pullback compatibility follows because pullback on raw
prespectra carries stabilization generators to stabilization generators and
commutes with the level-zero free functor.  The smooth-direct-image statement
is the same raw Thom-linearity used to construct \(p_\sharp^0\).  The final
sentence records Remark~\ref{rem:bn-diff-Q-not-fully-faithful}.
\end{proof}

\begin{theorem}[Canonical intrinsic de Rham realization]
\label{thm:bn-diff-intrinsic-dr-realization}
Let \(\mathcal D_{\dR}^{\mathrm{st}}(S)\) and \(\Gamma_{\dR,S}\) be the
stable intrinsic de Rham coefficient category and coefficient-sheaf functor of
Section~\ref{sec:infinitesimal-stable-dr-interface}.  Define
\[
  \boxed{
  \operatorname{Real}_{\dR,S}^0
  :=
  Q_S\Gamma_{\dR,S}:
  \mathcal D_{\dR}^{\mathrm{st}}(S)
  \longrightarrow
  \mathcal R_S^0.
  }
\]
This functor is exact.  For every smooth finite-presentation morphism
\(f:T\to S\), there is a canonical coherent equivalence
\[
  \boxed{
  f^{*,0}\operatorname{Real}_{\dR,S}^0
  \xrightarrow{\simeq}
  \operatorname{Real}_{\dR,T}^0f_{\dR}^{*,\mathrm{st}}.
  }
\]
Thus the intrinsic finite-infinitesimal geometry has a constructed stable
realization into the Thom-stable BNV coefficient layer.  No full-faithfulness,
essential-surjectivity, blanket colimit-preservation, or identification of
\(L_{\dR}\) with \(\mathbb H_S\) is asserted.
\end{theorem}

\begin{proof}
The functor \(\Gamma_{\dR,S}\) is exact by
Proposition~\ref{prop:infinitesimal-stable-dr-coefficient-sheaf}, and \(Q_S\)
is exact by the preceding proposition, so their composite is exact.  For a
smooth finite-presentation morphism, paste the canonical equivalences
\[
  f^{*,0}Q_S\simeq Q_Tf^*_{\Nis}
\]
and
\[
  f^*_{\Nis}\Gamma_{\dR,S}
  \simeq
  \Gamma_{\dR,T}f_{\dR}^{*,\mathrm{st}}.
\]
Their identity, composition, and higher coherences are inherited from the
coherences of the two factors.  The final sentence records precisely the scope
of the constructed bridge.
\end{proof}

\begin{remark}[Relation with differential function spectra]
The BNV reconstruction square
\[
  E\simeq
  \mathcal Z_S(E)
  \times_{\Disc_S\mathsf Z_S(E)}
  \Disc_S\mathsf U_S(E)
\]
is the stable homotopy-pullback pattern underlying the Hopkins--Singer
differential-function construction and its sheaf-of-spectra formulation in
Bunke--Nikolaus--V\"olkl
\cite{HopkinsSinger2005}\cite[Proposition~3.5]{BunkeNikolausVoelkl}.  The
present construction differs in two geometric respects: the stable sheaf theory
is the spectral Nisnevich theory, and Thom periodicity is imposed before the
BNV homotopification.  The smooth lift
\(\widetilde p_!=p_\sharp^0\Sigma^{-T_p,0}\) is then constructed from the same
spectral geometry which produces motivic smooth purity, while its canonical
comparison
\[
  \widetilde p_!\longrightarrow\widehat{\widetilde p_!}
\]
and the equivalent defect \(\omega_p\) record whether this geometric lift
already agrees with its canonical BNV completion, rather than hiding the issue
inside an extra transfer datum.
\end{remark}

\section{The brave-new differential-cohomology and Umkehr package}
\label{sec:bn-diff-final}

\begin{theorem}[Brave-new differential cohomology and verified differential Umkehr]
\label{thm:bn-diff-final-package}
For every spectral base in the ranges stated in the individual constructions,
the preceding sections establish the following.

\begin{enumerate}
\item The big affine spectral Nisnevich \(\infty\)-topos and the
finite-infinitesimal arrow topos are constructed from the spectral Nisnevich
geometry of Chapters~1--3, and restriction along the identity-arrow embedding
sits in
\[
  i_!\dashv i^*\dashv i_*\dashv i^!.
\]

\item The square-zero-generated spectral de Rham reflector \(L_{\dR}\) is an
accessible left-exact localization, with relative de Rham objects, finite
crystals, formal lifting theory, smooth effectivity and the reduced-classical
comparison constructed above.

\item Stabilizing the relative de Rham-local slice gives
\[
  \mathcal D_{\dR}^{\mathrm{st}}(S)
  =
  \operatorname{Sp}((\mathcal X_{\Sp,/S}^{\dR})_*),
\]
and affine-smooth evaluation constructs an exact coefficient-sheaf functor
\[
  \Gamma_{\dR,S}:
  \mathcal D_{\dR}^{\mathrm{st}}(S)
  \longrightarrow
  \mathcal N_S^{\Sp}.
\]
Composing with \(Q_S\) gives the canonical stable realization
\[
  \operatorname{Real}_{\dR,S}^0:
  \mathcal D_{\dR}^{\mathrm{st}}(S)
  \longrightarrow
  \mathcal R_S^0,
\]
compatible with smooth pullback.  This bridge does not identify
\(L_{\dR}\) with BNV affine-line homotopification.

\item For qcqs \(S\) in the motivic range, stable spectral Nisnevich sheaves
form the stable presentably symmetric monoidal category
\[
  \mathcal N_S^{\Sp}
  =\operatorname{Shv}_{\Nis}
  (\operatorname{Sm}^{\fp}_{/S};\Sp).
\]
Its raw Thom object \(\mathbb T_S^{\mathrm{raw}}\) defines the symmetric Thom
prespectrum presentation category \(\mathcal P_S^{\mathrm{pre}}\).

\item Localization only at the raw Thom stabilization morphisms produces the
Thom-stable differential category
\[
  \boxed{
  L_{\mathrm{Th},S}:
  \mathcal P_S^{\mathrm{pre}}
  \longrightarrow
  \mathcal R_S^0.
  }
\]
This localization is exact and strong symmetric monoidal, and
\[
  \mathbb T_S^0
  =L_{\mathrm{Th},S}F_{0,S}(\mathbb T_S^{\mathrm{raw}})
\]
is tensor-invertible.  Formal raw Thom periodization is related to this
category by canonical colimit-preserving strong symmetric monoidal functors
\[
  \mathcal N_S^{\Sp}[(\mathbb T_S^{\mathrm{raw}})^{-1}]
  \xrightarrow{\ \Psi_S\ }
  \mathcal R_S^0
  \xrightarrow{\ \Phi_S\ }
  \mathcal N_S^{\Sp}[(\mathbb T_S^{\mathrm{raw}})^{-1}],
  \qquad
  \Phi_S\Psi_S\simeq\id,
\]
compatible with $Q_S$ and the formal periodization functor.  No equivalence
between these two categories is asserted before affine-line localization,
because symmetry of the raw coordinate has not been established at that
stage.  No affine-line localization is used to construct
$\mathcal R_S^0$.

\item Affine-line homotopification of the already Thom-stable category is
canonically the brave-new stable motivic category:
\[
  \boxed{
  \mathbb H_S:\mathcal R_S^0\longrightarrow\SH(S).
  }
\]
The inclusion of local objects occurs in the exact triple
\[
  \boxed{
  \mathbb H_S\dashv\Disc_S\dashv\mathbb S_S.
  }
\]
Thus the BNV deformation sector measures failure of affine-line invariance
inside a Thom-stable theory, rather than a mixture of affine-line and Thom
stabilization defects.

\item Every \(E\in\mathcal R_S^0\) has canonically constructed underlying,
secondary, deformation, BNV-cycle and localized-cycle objects
\[
  \mathsf U_S(E),\quad
  \mathsf S_S(E),\quad
  \mathcal A_S(E),\quad
  \mathcal Z_S(E),\quad
  \mathsf Z_S(E),
\]
and is reconstructed by
\[
  \boxed{
  E
  \simeq
  \mathcal Z_S(E)
  \times_{\Disc_S\mathsf Z_S(E)}
  \Disc_S\mathsf U_S(E).
  }
\]
Equivalently, \(\mathcal R_S^0\) is the characteristic pullback category of
triples
\[
  U\in\SH(S),
  \qquad
  P\in\mathcal R_S^{\mathrm{BNV}},
  \qquad
  \phi:U\to\mathbb H_SP.
\]

\item The Thom-stable affine-line adjunction constructs the exact cylinder
\[
  \mathbb I_S=q_{S,*}^0q_S^{*,0}
\]
and the universal BNV integration transformation
\[
  \boxed{
  \int_{\mathbb A^1}^{\mathrm{BNV}}:
  \mathbb I_S\mathcal Z_S
  \longrightarrow
  \mathcal A_S
  }
\]
with normalized homotopy formula
\[
  1^*-0^*
  \simeq
  a_S\circ
  \int_{\mathbb A^1}^{\mathrm{BNV}}
  \circ\mathbb I_SR_S.
\]

\item The assignments \(S\mapsto\mathcal R_S^0\) have arbitrary inverse and
direct image in the standing base range, smooth direct image
\[
  p_\sharp^0\dashv p^{*,0},
\]
smooth base change, the smooth projection formula, coherent composition and
finite Nisnevich descent.  Their affine-line homotopifications recover the
corresponding stable motivic operations of Chapter~6.

\item Every finite locally free module \(E\) determines an invertible
differential Thom object
\[
  \operatorname{th}_S^0(E)\in\Pic(\mathcal R_S^0),
\]
and vector-bundle algebraic \(K\)-theory acts through a natural differential
\(J\)-homomorphism
\[
  \boxed{
  J_S^0:K(S)\longrightarrow\Pic(\mathcal R_S^0).
  }
\]
It is constructed affinely from split-exact vector-bundle \(K\)-theory and
Whitney-sum Thom objects and then descended by affine Zariski-basis Picard
descent, deduced from finite Nisnevich descent.  Applying it to the Waldhausen additivity simplices gives exact
Thom additivity with its full filtration coherence.  Homotopification carries
these twists to the motivic Thom twists of Chapter~6.

\item For every smooth separated finite-presentation morphism
\(p:X\to Y\), there is a canonically constructed Thom-stable exceptional lift
\[
  \boxed{
  \widetilde p_!
  =p_\sharp^0\Sigma^{-T_p,0}
  \dashv
  \Sigma^{T_p,0}p^{*,0}
  =\widetilde p^{\,!}.
  }
\]
It satisfies coherent composition, base change and projection formulas, and
its motivic shadow is exactly the smooth exceptional adjunction of Chapter~8:
\[
  \boxed{
  \mathbb H_Y\widetilde p_!
  \simeq
  p_!\mathbb H_X,
  \qquad
  \mathbb H_X\widetilde p^{\,!}
  \simeq
  p^!\mathbb H_Y.
  }
\]

\item Every constructed exact lift
\(T:\mathcal R_X^0\to\mathcal R_Y^0\) with constructed motivic shadow
\(\mathbb H_YT\simeq F\mathbb H_X\) has a canonical BNV correction
\(\widehat T\), a canonical comparison
\[
  \kappa_T:T\to\widehat T,
\]
and a canonical cycle defect
\[
  \boxed{
  \omega_T:
  \mathcal Z_YT
  \longrightarrow
  \mathcal Z_YT\mathcal Z_X.
  }
\]
Its canonical obstruction object is
\[
  \boxed{
  \operatorname{Obs}_{\mathrm{BNV}}(T)
  :=\cofib(\omega_T)
  \simeq
  \Sigma\,\mathcal Z_YT\mathcal S_X.
  }
\]
The correction \(\widehat T\) always preserves motivic-local objects and is
canonically BNV-idempotent:
\[
  \operatorname{Obs}_{\mathrm{BNV}}(\widehat T)\simeq0,
  \qquad
  \widehat{\widehat T}\simeq\widehat T.
\]
For the original lift \(T\), the following are equivalent:
\(\omega_T\) is an equivalence;
\(\operatorname{Obs}_{\mathrm{BNV}}(T)\simeq0\); \(T\) preserves
motivic-local objects; the localization comparison
\(T\Disc_X\to\Disc_YF\) is an equivalence; and
\(\kappa_T\) is an equivalence.  Thus the obstruction measures precisely
whether the original lift already agrees with its canonical BNV completion;
no additional ``cycle data'' or admissibility axiom is introduced.

\item For the smooth lift \(\widetilde p_!\), the canonical BNV completion
\[
  \widehat{\widetilde p_!}
\]
exists unconditionally, is exact, preserves motivic-local objects, and has
motivic shadow \(p_!\mathbb H_X\).  Affine-line Beck--Chevalley and the
normalized cylinder exchange for the original geometric lift are constructed
unconditionally from smooth base change and the tangent Thom twist.  Whenever
\[
  \kappa_p:
  \widetilde p_!
  \xrightarrow{\simeq}
  \widehat{\widetilde p_!},
\]
equivalently whenever \(\omega_p\) is an equivalence, the original lift itself
carries the canonical BNV cycle and deformation exchanges and
\[
\begin{tikzcd}[column sep=huge,row sep=large]
\widetilde p_!\mathbb I_X\mathcal Z_X
  \ar[r,"\widetilde p_!\int_X^{\mathrm{BNV}}"]
  \ar[d]
&
\widetilde p_!\mathcal A_X
  \ar[d,"\simeq"]
\\
\mathbb I_Y\mathcal Z_Y\widetilde p_!
  \ar[r,"\int_Y^{\mathrm{BNV}}"']
&
\mathcal A_Y\widetilde p_!
\end{tikzcd}
\]
commutes canonically.

\item If \(p:\coprod_{i=1}^dY\to Y\) is split finite \'etale, then
\[
  \widetilde p_!(E_1,\ldots,E_d)
  \simeq
  \bigoplus_iE_i,
  \qquad
  \omega_p:
  \mathcal Z_Y\widetilde p_!
  \xrightarrow{\simeq}
  \mathcal Z_Y\widetilde p_!\mathcal Z_X,
\]
and therefore
\[
  \kappa_p:
  \widetilde p_!
  \xrightarrow{\simeq}
  \widehat{\widetilde p_!}.
\]
The BNV reconstruction sectors and affine-line cylinder commute with this transfer, so the
BNV interval integral is natural under it.  After homotopification it is
precisely the brave-new motivic finite \'etale Umkehr.

\item More generally, if a finite \'etale morphism is Nisnevich-locally split,
then Nisnevich descent and smooth base change imply that
\(\widetilde p_!\) preserves motivic-local objects.  Hence
the canonical cycle defect \(\omega_p\) and the canonical comparison
\[
  \kappa_p:
  \widetilde p_!
  \xrightarrow{\simeq}
  \widehat{\widetilde p_!}
\]
are equivalences, and the same differential-Umkehr and BNV-integration
compatibility holds globally.
\end{enumerate}

No differential-cohesive axiom, differential six-functor axiom, blanket
realization principle, admissibility condition, or unspecified cycle-transfer
datum is added.  For a general smooth separated finite-presentation morphism,
the canonical BNV completion \(\widehat{\widetilde p_!}\) exists and preserves
motivic-local objects.  The explicit obstruction
\(\operatorname{Obs}_{\mathrm{BNV}}(p)\), equivalently the failure of
\(\kappa_p:\widetilde p_!\to\widehat{\widetilde p_!}\) to be an equivalence,
measures whether the original geometric lift already equals that completion.
The chapter proves this comparison is an equivalence in the finite \'etale
classes above and makes no general vanishing claim beyond those constructions.
\end{theorem}
